% !TeX program = lualatex
% One-source manuscript. The original Nous PNG follows the end of the document.
% Run LuaLaTeX three times in a writable directory. The release uses TeX Live 2026.
% For stable timestamps, set SOURCE_DATE_EPOCH=1795582800 and FORCE_SOURCE_DATE=1.
\documentclass[12pt,letterpaper]{article}
\usepackage{iftex}\RequireLuaTeX
% Suppress the path-dependent PDF trailer ID for reproducible release builds.
\pdfvariable suppressoptionalinfo 512\relax
\usepackage[margin=1.02in,headheight=16pt,headsep=14pt,footskip=28pt]{geometry}
\usepackage{fontspec}\setmainfont{Latin Modern Roman}\setsansfont{Latin Modern Sans}
\usepackage{amsmath,amssymb,mathtools,booktabs,tabularx,array,longtable,graphicx,microtype,xcolor,fancyhdr,enumitem,hyperref,bookmark}
\definecolor{ink}{HTML}{183831}\definecolor{accent}{HTML}{095D58}\definecolor{mute}{HTML}{455E55}
\hypersetup{colorlinks=true,urlcolor=accent,linkcolor=accent,pdftitle={A regular chart at four},pdfauthor={Justin Candler and Uriel}}
\setlength{\parindent}{0pt}\setlength{\parskip}{5pt plus .5pt}\setlength{\emergencystretch}{3em}
\linespread{1.10}\raggedbottom\pagestyle{fancy}\fancyhf{}\fancyhead[L]{\small\sffamily\color{mute}Nous Independent Research}\fancyhead[R]{\small\sffamily\color{mute}A regular chart at four}\fancyfoot[C]{\small\sffamily\thepage}
\renewcommand{\headrulewidth}{.3pt}\setcounter{tocdepth}{1}\setlist{itemsep=2pt,topsep=4pt}
\newcolumntype{Y}{>{\raggedright\arraybackslash}X}

\newcommand{\srclink}[2]{\par\noindent{\small\sffamily\href{https://github.com/nousentllc/riemann-hypothesis/blob/3f39c057bb593d610473bd9d477cad518a6ba21c/research/uriel/experiments/uncapped-growth/#1}{#2}}\par}
\newcommand{\passmark}{\begingroup\color{accent}\rule{7pt}{7pt}\endgroup}
\newcommand{\failmark}{\begingroup\color{mute}\setlength{\fboxsep}{0pt}\fbox{\rule{0pt}{6.2pt}\rule{6.2pt}{0pt}}\endgroup}
\newcommand{\keyterm}[1]{\textbf{\color{ink}#1}}
\newcommand{\code}[1]{\texttt{\detokenize{#1}}}
\newcommand{\dossier}[2]{\clearpage\subsection*{#1}\phantomsection\label{#2}}
\newcommand{\reading}[2]{\par\smallskip\textsf{\color{accent}\bfseries #1}\enspace #2\par}
\newcommand{\ruleline}{\par\vspace{1pt}{\color{accent}\rule{\linewidth}{.35pt}}\par}
\newcommand{\yes}{\mathsf Y}\newcommand{\no}{\mathsf N}\newcommand{\unk}{\mathsf U}
\directlua{
 local source = assert(io.open(status.filename, "rb")); local rows={}; local take=false
 for line in source:lines() do
 if line:match("^BEGIN FIGURE nouslogo") then take=true
 elseif line:match("^END FIGURE") then take=false
 elseif take then table.insert(rows,line:match("^[0-9a-f]+")) end
 end
 source:close()
 function nousLogo()
 local hex=table.concat(rows); assert(string.len(hex)>100000)
 local bytes=hex:gsub("..",function(pair) return string.char(tonumber(pair,16)) end)
 local path=tex.jobname.."-embedded-nous.png"; local existing=io.open(path,"rb")
 if existing then local previous=existing:read("*a"); existing:close(); assert(previous==bytes)
 else local file=assert(io.open(path,"wb")); file:write(bytes); file:close()
 function cleanNousLogo() os.remove(path) end end
 node.write(img.node(img.scan{filename=path})) end
}


\newcommand{\id}[1]{{\small\nolinkurl{#1}}}
\newcommand{\sourceline}[2]{\par\begingroup\small\ttfamily\raggedright\hangindent=2.35em\hangafter=1\noindent\makebox[2.35em][l]{\scriptsize\sffamily\color{mute}#1}\strut#2\par\endgroup}
\newcommand{\sourceblock}{\par\begingroup\setlength{\parskip}{0pt}\linespread{0.96}\selectfont}
\newcommand{\sourceend}{\endgroup\par}
\hypersetup{pdftitle={A regular chart at four},pdfsubject={Readable Gaussian fiber charts and local Haar normalization}}

\begin{document}
\pagenumbering{roman}\thispagestyle{empty}
\begin{center}\vspace*{.46in}
{\sffamily\small\bfseries\color{accent}NOUS INDEPENDENT RESEARCH}\par\vspace{.51in}
{\fontsize{32}{38}\selectfont\bfseries\color{ink}A regular chart at four}\par\vspace{.24in}
\begin{minipage}{.91\linewidth}\centering\Large Skew Haar normalization on smaller Gaussian vector boxes over $\mathbb Q_2(\sqrt2)$\end{minipage}
\par\vspace{.35in}\resizebox{!}{1.12in}{\directlua{nousLogo()}}
\par\vspace{.45in}{\large Justin Candler with Uriel, AI research assistant}\par 25 November 2026
\end{center}\vfill
\begin{center}\small\sffamily Gaussian identity form of rank three at the ramified dyadic place\\
Original source 40e56bab \enspace / \enspace Readable source e063a923\end{center}
\clearpage
\section*{Abstract}
For the Gaussian identity form in three variables over $K=\mathbb Q_2(\sqrt2)$, we construct a regular-fiber chart at the skew target $4i$. On the vector box around $(2,0,0)$ of scalar width $8O$, the Gram image is exactly $4+32O$. Fixing the other five scalar coordinates reduces inversion to the surjective isometry $u\mapsto u+2u^2$ of $O$. A separate construction realizes the readable edition's hyperbolic chart in the original Gaussian identity space, using a square root of $-7$. The field trace-one element $1/2$ suffices for this construction, although the integral trace obstruction persists.

With the original source Haar fixed, the scalar and split-skew masses of the parent box are $2^{-26}$ and $2^{-27}$. Each nonempty child obtained by reduction modulo $16O$ carries $1/1024$ of the parent fiber mass. Exactly 1024 of 4096 children meet four. A translated pivot box misses four and meets 36. The readable parameter preimage of an active centered box of width $W$ has two Gaussian sides of width $W$ and one skew side of width $W/2$. Directly pushing its parameter cells through the frame reproduces the original vector children. An erroneous extra restriction of the skew side leaves only 256 distinct children.

We give separate Lean replays against the two pinned editions, including the readable definition. The measured object is a regular coordinate patch. The global readable-place disintegration at this ramified place has additional hypotheses.

\medskip\textbf{Human--AI collaboration.} Justin Candler posed the smaller-box and trace-obstruction questions, directed the research series and specified the evidence and publication requirements. Uriel/Codex derived the charts, authored the new Lean and numerical checks, ran the replays and prepared this manuscript. The upstream formalization is separately attributed. Independent human mathematical review remains open.
\clearpage
\begingroup\small\setlength{\parskip}{0pt}\tableofcontents\endgroup
\clearpage\pagenumbering{arabic}
\section{The local question}
At the dyadic place of $\mathbb Q(\sqrt2)$ we have a Gaussian Hermitian space with three coordinates. The vector $(2,0,0)$ has Hermitian norm four. In the readable edition the skew parameter for this norm is $4i$. The preceding paper measured the entire coordinate lattice in the fiber. Its two values, $257/128$ and $257/256$, leave open a geometric question: can we see regular vectors themselves vary with the target, and can we measure a proper piece of the fiber rather than only its total?

A whole-lattice computation cannot answer the latter question by division. A regular conditional measure can distribute its mass quite unevenly inside a lattice. Multiplying an already constructed measure by one half makes every set half as heavy by definition. A substantive test starts independently from the six-dimensional volume of a restricted set of vectors, follows its images under the actual Gram map and recovers the two normalizations from those observations. We do this for a compact open box and its children. It produces both positive pieces of the fiber and many genuine zeroes.

Let $K=\mathbb Q_2(\pi)$, $\pi^2=2$, and write $O$ for its ring of integers. Normalize the absolute value by $|2|=1/2$ and scalar Haar by $\lambda(O)=1$. The absolute value measures distances, whereas $\lambda(2O)=1/4$ measures volume. Confusing the two changes the answer at more than one point in the argument. Each scalar has two dyadic coefficients. The compact, open Gaussian order $B=O+Oi$ has volume one in the coordinates $(1,i)$. This paper always observes the same vectors in those coordinates.

Write a vector as six scalars $(x,w_1,\ldots,w_5)$. The Hermitian identity form becomes
\[
    Q(x,w)=x^2+\sum_{j=1}^5 w_j^2 .
\]
We study the compact open set $P=(2+8O)\times(8O)^5$. Its image turns out to be exactly $C=4+32O$. In the metric on $K$, $|4|=1/4$ and every perturbation in $32O$ has norm at most $1/32$. Thus $C$ is a neighborhood of a regular target and excludes zero. Its image under $s\mapsto is$ is a neighborhood in the one-dimensional skew line.

\section{Three assertions with different hypotheses}
The first assertion is analytic and local. On $P$, retaining all five transverse coordinates and replacing the first by the Gram value is a homeomorphism onto the product of $C$ with the transverse box, after the evident coordinate scalings. For each allowed target there is exactly one pivot above each transverse parameter. The inverse is constructed on a complete ball and is continuous in the target. No trace-one element of an order enters.

The second assertion concerns measures. The exact scalar-fiber mass of $P$ in this chart is $2^{-26}$. Its mass with the split skew Haar is $2^{-27}$. Requiring all six coordinates to match their center modulo sixteen gives $2^{-36}$ and $2^{-37}$. Requiring the first coordinate to be ten modulo sixteen gives zero at four. At target $36$, the same pivot class can carry positive mass. A rule which used only the valuation of the first coordinate could not distinguish these cases.

The third assertion concerns the readable proof. Its \id{flatChart} uses a hyperbolic plane and a residual Hermitian form, and its target is encoded with a skew unit. Its algebraic formula can be read with a field trace-one element $1/2$. The stronger \id{IntegralPlace} structure demands integral trace data. The two requirements concern different rings. We compare the defining expression, its unit-pivot hypothesis and a hyperbolic frame for the original identity form. Constructing an entire instance, with its additive character and Fourier estimates, would be a substantially stronger result. This distinction matters if a chart is to be used inside subsequent readable theorems.

\section{Why the first scalar can be solved}
Put $x=2+8u$ and $w_j=8z_j$. For $u,z_j\in O$ the expansion is
\[
 Q(2+8u,8z)=4+32\left(u+2u^2+2\sum_j z_j^2\right).
\]
The powers have independent origins. The cross term of the pivot square is $2\cdot2\cdot8u=32u$; all six quadratic perturbations have coefficient $64$. After division by $32$ the linear coefficient is one. The remaining coefficient is two, strictly smaller than one in absolute value. This provides a choice of coordinates in which inversion does not consume any further precision.

Let $h(u)=u+2u^2$. For $u,v\in O$ the difference factors as
\[
          h(u)-h(v)=(u-v)(1+2(u+v)).
\]
The parenthesized factor has norm one. Indeed $|2(u+v)|\le1/2$, so the two summands have unequal norms and the ultrametric inequality is an equality at the larger norm. It follows that $|h(u)-h(v)|=|u-v|$. This proves injectivity as well as preservation of each precision digit. It does not give surjectivity over an arbitrary incomplete domain.

For $b\in O$ consider the iteration $u\mapsto b-2u^2$. Its image stays in $O$. The same factorization without the leading one shows that its Lipschitz constant is at most $1/2$. The closed integer ball is complete. The contraction theorem supplies a fixed point, which solves $h(u)=b$. Thus $h$ is a surjective isometry and its inverse is continuous. The proof does not require a division by the derivative at each Newton step. Starting from zero already gives an explicit sequence of approximants whose error decays by at least one binary coefficient digit per iteration.

At target $4+32t$, fix $z$ and let $d(z)=2\sum z_j^2$. The solution is $u=h^{-1}(t-d(z))$. To recover the parameters from a vector in $P$, read $z_j=w_j/8$ and $t=(Q-4)/32$. Both divisions take place in the field on sets where the results are integral. This gives a joint homeomorphism, not just a fiberwise count. It also gives a simple surjectivity witness for every $t$: choose $z=0$ and solve the single pivot equation.

\section{Two independent regular coordinates}
The normalized chart acts by
\[
 (u,z)\longmapsto(h(u)+d(z),z).
\]
For fixed $z$ the first operation is a surjective isometry followed by translation. For fixed target the graph is a continuous injection of $O^5$ into the original vectors; the five transverse components divided by eight form its inverse on the image. These two descriptions permit different tests. A solver that produced the correct Gram value but duplicated a branch would pass the first equation and fail inverse uniqueness. A map with the correct pointwise inverse but an incorrect measure normalization would pass both equations and fail the volume tests.

There is also an exact local metric derivative before any reduction. For $x,y\in2+8O$, the sum $x+y$ belongs to $4+8O$ and has norm $1/4$. Thus
\[
          |Q(x,w)-Q(y,w)|=|x-y|/4 .
\]
The expression is independent of the five transverse coordinates, which cancel in the difference. If the pivot changes by $2^r O$, the target varies over $2^{r+2}O$. Every target residue of this finer width has one pivot residue above it, once the target and the transverse vector have been placed in the chart. The residue width in $x$ is four times the target width in metric units, hence sixteen times as large in scalar Haar volume.

The isometry argument gives more than a cardinality statement. If only finite permutations had been shown, existence of a compatible inverse through all levels would require a further argument. Here the field inverse directly identifies a nested sequence of target residues with a nested sequence of pivot residues. It also excludes extra solutions in the same vector box. Other boxes can contain other sheets of the same regular fiber, including the branch around $-2$. The chosen box deliberately isolates the branch through $+2$.

\section{The compact-open slab calculation}
Use $D_r(a)=a+2^rO$. Its Haar measure is $4^{-r}$ for every $a$. Let $S\subseteq(8O)^5$ be measurable and write $\kappa=\lambda^{\otimes5}$. Suppose first that the pivot must lie in $D_3(2)$. For a fixed target $s\in C$ and transverse point $w\in S$, the chart produces a root $y(w)$ in that pivot ball. If $r\ge3$, then
\[
 \{x\in D_3(2): Q(x,w)\in D_{r+2}(s)\}=D_r(y(w)).
\]
The forward inclusion follows from the metric derivative. The reverse inclusion also needs the whole smaller ball to lie in $D_3(2)$. This holds since two intersecting balls of nested radii in an ultrametric space are nested. The measure of the section is $4^{-r}$, independent of the unknown root. There is no need to compute its digits to integrate it.

Fubini now gives the six-dimensional slab volume $\kappa(S)4^{-r}$. Dividing by $\lambda(D_{r+2}(s))=4^{-(r+2)}$ gives the scalar density $16\kappa(S)$. The skew split measure, with Gaussian source held fixed, is $2j_*\lambda$ for $j(s)=is$. It gives the same target set twice the volume, hence the skew density is $8\kappa(S)$. Taking $S=(8O)^5$ yields $16\cdot 4^{-15}=2^{-26}$ and $8\cdot 4^{-15}=2^{-27}$. Every sufficiently small target slab gives these ratios exactly. These values are local contributions to the earlier whole-lattice numbers.

A direct consistency check avoids ratios. The volume of the parent vector box is $4^{-18}=2^{-36}$. The scalar image box $C$ has volume $4^{-5}=2^{-10}$. The skew image has volume $2^{-9}$. Multiplying $2^{-26}$ by $2^{-10}$ and $2^{-27}$ by $2^{-9}$ gives $2^{-36}$ in each convention. This also fixes the orientation of the Jacobian; a reciprocal guess would fail even on the full image box.

\section{What the resolving digit selects}
At the exact target four, $t=0$. Reducing $h(u)+d(z)=0$ modulo $2O$ gives $u=0$ in $O/2O$. Both quadratic contributions vanish. This fact is stronger than the requirement that $u$ have even constant coefficient. There are four elements of $O/2O$, because each of the two scalar coefficients has one binary digit. A pivot in the class of $\pi$ still fails this condition. Meanwhile all four choices in each transverse coordinate are allowed.

Passing from width eight to sixteen splits each scalar ball four ways. The vector box has $4^6=4096$ children; exactly $4^5=1024$ carry mass at four. Their pivot has the single allowed residue and their five transverse residues are arbitrary. All these children have equal mass since each has the same transverse Haar volume $4^{-20}$. The ratio between a surviving child's mass and the parent's is $4^{-5}=1/1024$. It is also the ratio at which the surviving children exhaust the parent. An equal allocation to all 4096 children would underestimate every surviving one by a factor of four and assign a positive mass to each empty one.

The criterion describes a change of target entirely within the same open chart. For $s=4+32t$ the surviving pivot digit is the class of $t$ modulo $2O$. As $t$ ranges over $0,1,\pi,1+\pi$ the active pivot cycles among four disjoint classes. In particular the box centered at ten, with transverse center zero and width sixteen, contains no vector at target four. At $s=36$, the normalized target is one and the same box becomes active. This switch is determined by the reduced polynomial. The center $10$ has square $100$. An exact witness above $36$ is $(-6,0,0)$, since $-6-10=-16$ and $(-6)^2=36$. Changing a transverse coordinate generally changes this root. Agreement modulo a selected precision is sufficient to choose a class, not to reconstruct the point.

Before the resolving digit the positive and empty children look identical to the target test. At target precision five all of $P$ maps to the single class of four modulo $32O$. In a fixed width-sixteen child the coarse scalar density is its volume $4^{-24}$ divided by $4^{-5}$, or $4^{-19}=2^{-38}$. At precision six the centered child's entire volume still lies in the class, giving scalar density $4^{-18}=2^{-36}$; the child centered at ten now contributes zero. The split densities are half of each of these. A positive finite-depth observation for a purportedly empty fiber can thus be genuine. It falsifies the inference from that observation to a pointwise mass.

\section{A hand computation at a translated transverse point}
Take the vector child with first coordinate $2\bmod16O$, first transverse coordinate $8(1+\pi)\bmod16O$, and remaining transverse coordinates $0\bmod16O$. The transverse vector at the displayed center has $z_1=1+\pi$. Its normalized offset is $2(1+\pi)^2=6+4\pi$. At target four we must solve $h(u)=-6-4\pi$. Modulo $2O$ the target is zero, so the branch belongs to the centered pivot child even though this transverse vector is a unit multiple of eight.

For a hand check to the next two levels, observe that $u$ is twice an integer of $O$, hence $2u^2\in8O$. Consequently
\[
       u\equiv -6-4\pi\equiv 2+4\pi\pmod{8O}.
\]
It follows that $x=2+8u\equiv18+32\pi\pmod{64O}$. Direct expansion gives $Q(18+32\pi,8+8\pi,0,\ldots,0)\equiv4\pmod{256O}$. One can perform this check without the inverse theorem: the first square is $2372+1152\pi$, the second $192+128\pi$, giving $2564+1280\pi=4+256(10+5\pi)$. The displayed vector need not be the exact root. Its congruence illustrates both the existence theorem and the two-digit shift between pivot and target precisions.

The target mass of this translated child is identical to the centered child's mass. Both have transverse volume $4^{-20}$, and in both cases the surviving pivot residue is the same. The roots nonetheless differ with the transverse vector. A uniform distribution on transverse coordinates therefore need not make the pivot coordinate uniform inside the local fiber. This is why reporting both the graph and the measure matters.
\section{The readable chart in its own variables}
The readable source starts with a commutative star ring $A$, a skew unit $\delta$, an element $\tau$ satisfying $\tau+\bar\tau=1$, and a Hermitian residual form $T$. On a hyperbolic plane plus $T$, write a vector as $(X,p,v)$, and use $b$ for the skew target. Then $s=\delta^{-1}b$ is the Hermitian target. For a unit pivot $p$, a free skew coordinate $a$, and a residual vector $v$, the readable formula is
\[
  X=\bar p^{-1}\tau\bigl(s-\bar v^{\mathsf T}Tv\bigr)+pa .
\]
The Gram value is $\bar Xp+\bar pX+\bar v^{\mathsf T}Tv$. Substitution leaves two copies of the Hermitian difference, with coefficients $\bar\tau$ and $\tau$. They add to one. The terms containing $a$ have opposite signs because $\bar a=-a$. This is the cancellation responsible for the chart equation. The source uses a ring inverse defined everywhere; the step that cancels $p$ requires the explicit unit hypothesis. Merely evaluating the formula at a nonunit would not establish the desired equality.

There are three distinct roles here for algebra, topology and integration. Forming the displayed point needs ring operations, a trace-one scalar and the unit inverse. Identifying its Gram value uses the Hermitian and skew symmetries. Pushing a source measure through this parameterization requires additional Haar and measurability data. In the readable source the class \id{Place} bundles further structure, including a character and a Fourier kernel bound. Its \id{IntegralPlace} extension adds integral hypotheses used elsewhere. The evaluation of this algebraic point does not use every field of those structures. We extract its formula into a separate definition with only the needed arguments and compare it by unfolding the pinned readable definition. This explains precisely which use of the source is possible even where no integral trace-one element exists.

In a single residual Gaussian coordinate let $p=p_0+ip_1$, $v=v_0+iv_1$, $a=i\alpha$ and $\tau=1/2$. For residual coefficient $+1$, put $D=p_0^2+p_1^2$ and $k=(s-v_0^2-v_1^2)/(2D)$. Then $X=(p_0k-p_1\alpha)+i(p_1k+p_0\alpha)$. For $D\ne0$ the inverse free parameter is
\[
        \alpha=(p_0\,\Im X-p_1\,\Re X)/D .
\]
A residual coefficient $-1$ changes the sign of the residual norm inside $k$. It leaves the denominator and recovery of $\alpha$ alone. We keep this sign explicit because the frame for the original identity form below has residual coefficient $-1$.

\section{Finding a hyperbolic frame in the original space}
A unit-pivot identity for an abstract split Hermitian form is not yet an identity for a particular Gaussian identity matrix. An explicit frame removes that ambiguity. The first ingredient is an element $c\in O$ with $c^2=-7$. The same subroutine used for the vector chart supplies one. Solve $h(u)=-1$ and put $c=1+4u$. Then $c^2=1+8h(u)=-7$. This proves existence in the local field without selecting a square root over the real numbers or adjoining a new scalar. The chosen root is characterized by $c\equiv1\pmod{4O}$. For this root, the normalized equation and uniqueness determine all its subsequent digits.

Write $b_0=1+i$ and $d_0=c+2i$. Their norms are $2$ and $-3$, respectively. Negative rational norm values cause no contradiction in a locally completed field. The norm here is the algebraic Gaussian norm $a\mapsto\bar a a$ with values in $K$. It is distinct from the nonnegative real absolute value used in the contraction argument. In particular $\bar b_0 b_0+\bar d_0d_0=-1$. Let $v_1=(1,0,0)$ and choose
\[
  e=(1,b_0,d_0),\qquad f=v_1-e/2,\qquad
  z=(0,\bar d_0,-\bar b_0).
\]
The Hermitian form has conjugate-linear first entry. It gives $h(e,e)=0$, $h(e,f)=1$, $h(f,f)=0$, $h(e,z)=h(f,z)=0$ and $h(z,z)=-1$. For instance, $h(e,z)=\bar b_0\bar d_0-\bar d_0\bar b_0=0$. The identity $h(f,f)=1-1/2-1/2+0=0$ uses $h(v_1,e)=1$. Arranging these six pairings as a matrix gives a hyperbolic plane and a residual line of norm $-1$.

The determinant of the matrix $M=[e\ f\ z]$ is $-1$. Add $e/2$ to its second column; the resulting second column is $v_1$. Expanding there gives $\bar b_0 b_0+\bar d_0 d_0=-1$. Thus $M$ is invertible over the field algebra and preserves the ambient volume. The latter statement follows either from the determinant modulus in an ambient place or by restriction of scalars: the $K$-determinant is the Gaussian norm of $-1$, equal to one. The second column contains $1/2$, so the matrix need not carry the full coordinate lattice to itself. Determinant and preservation of each coordinate box are separate questions. We compute all smaller box masses in the original coordinates.

At $(2,0,0)=e+2f$ the hyperbolic data are $(X,p,v)=(1,2,0)$. Taking $a=0$, $\tau=1/2$ and $s=4$ in the readable formula gives $X=(1/2)(1/2)4=1$. Thus the abstract hyperbolic center and the center of the contraction chart are the same vector. Nearby, the hyperbolic pivot is $h(e,V)$. For $V\in P$ its difference from two has real and imaginary coefficients in $8O$, since $e$ is integral. A Gaussian number with real part in $2+8O$ and imaginary part in $8O$ is a unit in the ambient algebra: its real part has norm $1/2$, its imaginary part norm at most $1/8$, and their squares cannot cancel. This argument uses invertibility in the field algebra. Requiring the pivot to be an integral unit would incorrectly exclude the center of the example.

\section{Why trace one still matters}
An element $a+ib$ of the coordinate order has trace $2a\in2O$. The element one does not lie in $2O$. In the larger cyclotomic order at this place the previous study found a trace image contained in $\pi O$, which likewise excludes one. The existence of a point near $4i$ cannot contradict either observation. Our trace-one scalar is $1/2$ in the ambient field. In the displayed hyperbolic chart it multiplies a difference of Hermitian target values. Neither this multiplication nor the equation that follows asks for an integral trace-one scalar.

The per-box calculation provides a second route around the same obstruction. It works in the original coordinates, where no trace splitting is necessary to write or invert $Q$. Once the vector source has been fixed, the earlier split calculation fixes the measure on the skew targets. Therefore a missing integral-place instance cannot make the per-box uncertainty arbitrary. It obstructs a particular upstream specialization. A compatible regular kernel on this local patch still has to obey the actual source volumes. Conversely, the existence of our local chart gives no additive character, no Plancherel bound and no estimate for an Euler product. Those are independent responsibilities of the original theorem interfaces.

The choice of a patch matters as much as the choice of ring. The hyperbolic formula covers the locus where its chosen pivot is a unit. The contraction chart additionally fixes $x\bmod 8O$ and the five transverse residues. These restrictions put us inside one sheet, make its inverse single-valued and give a useful compact test for integration. An argument on this patch cannot be extended to the whole level set by translating the vector box arbitrarily; translation changes a quadratic form. A global comparison would require more charts and agreement on overlaps or a separate uniqueness result applying to the global kernels.

\section{Why the surface measure has a factor of sixteen}
Write $\kappa=\lambda^{\otimes5}$ on physical transverse coordinates. The local scalar-fiber measure assigns a vector set its transverse preimage under the graph, with multiplier sixteen. Where the skew base is $2j_*\lambda$, the multiplier is eight. This sounds large until the two places where scale enters are kept distinct. The pivot-to-target map contracts metric distances by four. In the two-dimensional dyadic coefficient lattice that corresponds to a volume contraction by sixteen. Moving from scalar to split skew Haar then doubles the denominator separately. Multiplying the two metrics or calling the entire operation a coordinate relabeling loses one of these steps.

Here is a direct derivation in normalized coordinates. The pushforward of the probability Haar on $O^6$ under $(u,z)\mapsto(h(u)+d(z),z)$ is again the probability product. To see why, first observe that closed balls of a fixed coefficient precision have a common mass. Any surjective isometry sends each ball to a ball of that precision. The closed balls over all precisions form a basis, are closed under nonempty intersections and determine a finite Borel measure. The distribution of $h(u)$ is therefore the same as that of $u$. Translating by $d(z)$ also preserves the conditional measure for each fixed $z$. Applying Fubini yields the joint assertion even when the transverse test and the target test are dependent.

The physical map multiplies all six scalar widths by eight. It multiplies their product volume by $4^{-18}$. Replacing the normalized parameter $t$ by $s=4+32t$ multiplies target volume by $4^{-5}$. Replacing $z$ by $w=8z$ multiplies transverse volume by $4^{-15}$. To recover the source volume as a product of target and transverse volume requires the ratio
\[
               \frac{4^{-18}}{4^{-5}4^{-15}}=16.
\]
This derivation does not depend on selecting a depth or a particular measurable box. The physical slab proof checks the same number by cutting the original box at every finer precision and measuring its sections directly. Agreement of these derivations would be lost if one used three Gaussian widths in the numerator and five scalar widths in the transverse denominator without converting units.

For a measurable vector set $E$ restricted to the patch, the corresponding formula is
\[
 \sigma_s(E)=16\int_{(8O)^5}{\bf1}_E(w,y(s,w))\,d\kappa(w),\quad
 \rho_{is}(E)=8\int_{(8O)^5}{\bf1}_E(w,y(s,w))\,d\kappa(w).
\]
We order the displayed coordinates as $(w,x)$ to match the Fubini calculation. The polynomial itself has pivot $x$ first. The graph is continuous and the transverse region compact. A compactly supported continuous vector test consequently has a continuous fiber integral as $s$ varies in the patch, by dominated convergence. A change to the kernel on one point of the target could preserve all tube integrals but violate this continuity. The pointwise claim needs this regularity, or the explicit continuous-graph construction, in addition to the integral identity.

\section{How much of the earlier fiber is visible here?}
Under the scalar convention $P$ has mass $2^{-26}$ while the entire coordinate lattice has mass $257/128$. The fraction of the full lattice-fiber mass seen on $P$ is $1/(257\cdot2^{19})$. Under the split convention both the numerator and the denominator halve, so this fraction is invariant. The experiment therefore qualifies as a probe of a genuinely small part of the surface. A successful check here controls a small amount of total mass with considerable resolution. It cannot identify where the remainder lies.

By contrast, a surviving child carries one part in 1024 of the mass of $P$. This larger relative share is exactly explained by five independent transverse residue choices. A six-coordinate vector box loses six factors of four when narrowed, while a fiber locally has only five free coordinates. In the ambient space a width-sixteen child occupies $1/4096$ of the parent. At the selected target it occupies either zero or $1/1024$. Conditioning has excluded three potential pivot classes; it has left all the transverse classes. This is an observable difference between volume in six variables and volume on one regular equation.

The picture also explains why the earlier numerator 257 disappears from the local patch constants. That numerator records how many representations across all residue classes contribute to the entire lattice. Here we isolated a patch whose projection has constant metric derivative and exactly one solution per transverse vector. Its mass is a power of two times a transverse volume. A sum of many such contributions, including primitive and divisible strata with different scales, can produce the earlier number. We have not partitioned the entire lattice into this single type of chart.
\section{A compact-open measure test at every depth}
To make the numerical claim falsifiable, fix five transverse residue classes and a pivot class before observing their target values. Let the transverse set be $S$, and restrict the pivot to $D_m(a)$. If the local root for every $w\in S$ lies in this pivot cell and $D_m(a)\subseteq D_3(2)$, then for every $r\ge m$ there is a common section-volume formula. It applies at any chosen regular target $s$ in the chart image:
\[
 \mathrm{vol}\{(w,x):w\in S,\ x\in D_m(a),\ Q(x,w)\in D_{r+2}(s)\}
        =\kappa(S)4^{-r}.
\]
The root may vary with $w$. The content of the formula is precisely that this variation changes the center of the pivot ball but never its size. The Lean proof integrates an indicator over the original $x$ coordinate, with the section identified as an actual set. It does not infer the slab probability from the distribution of outputs of a finite solver.

At four, $m=3$ and $m=4$ with $a=2$ both satisfy the root condition for every $w\in(8O)^5$. Hence for any measurable $S$ in the transverse patch the same formula holds. A useful asymmetric test restricts one of the transverse coordinates modulo sixteen and leaves the remaining four modulo eight. Its transverse mass is $4^{-4}(4^{-3})^4=4^{-16}$. The local split mass is $8\cdot4^{-16}=2^{-29}$, exactly one quarter of the parent split mass. This differs from narrowing all six coordinates, which costs a factor of 1024 on an active child. Narrowing the pivot alone to its active child costs nothing at the exact target.

The corresponding empty test is informative only at an appropriate target precision. For the translated pivot child $D_4(10)$ a small slab centered at four cannot contain a vector with these restrictions once the target residue fixes $t\bmod2O$. For $s=36$ an analogous slab does contain such vectors. The two target points are separated in the skew line by $32i$. A coordinate change sending both to the same normalized target would predict the same active child and would be rejected by this pair of tests.

\section{Design of the independent finite experiment}
For each coefficient depth $d$, a scalar residue is an ordered pair in $(\mathbb Z/2^d)^2$. We square these pairs directly using
\[
       (a,b)^2=(a^2+2b^2,2ab)\bmod 2^d.
\]
A vector box is specified by six independent congruence conditions, one per scalar. For each scalar the program builds a histogram of the possible squares with their integer multiplicities. It convolves the six histograms using addition of output pairs. The coefficient at the target is the number $N$ of original vector residues representing that target in the specified box. At no point does the computation invert $h$, use a hyperbolic basis or take a numerical square root. This makes it an independent check on the algebraic parameterization.

All $4^{6d}$ residues of the ambient space are equally weighted. The target cell has scalar measure $4^{-d}$, so its observed scalar density is $N/4^{5d}$. The split density then uses twice the observed cell volume. The numerator is shared across conventions; the factor of two itself is inherited from an independent Haar split theorem. The new information in a successful test is the numerator for this specific vector set and its consistency across changes of target precision. An experiment that ran the same count twice and silently changed its normalization both times would give a weaker control.

The cross consists of five vector boxes, four targets and five coefficient precisions, giving 100 records. The five boxes are the parent, its centered child, an equally small child displaced in a transverse coordinate, the pivot child centered at ten and a box narrowed in one transverse coordinate only. The targets $4,36,4+32\pi,36+32\pi$ probe all four residue classes of $t$. Depths four through eight include both coarse aliasing and three precisions that resolve all the selected children. A separate enumeration inspects all 4096 source child classes and groups them by their depth-six target residue. The finite records contain exact, reproducible integer counts and rational densities. The permanence of the answer above these depths comes from a theorem.

\begin{center}\small\begin{tabular}{@{}lrrr@{}}\toprule
Vector restriction & Scalar at four & Split at four & Share of parent\\\midrule
Parent $P$ & $2^{-26}$ & $2^{-27}$ & $1$\\
All six centered modulo 16 & $2^{-36}$ & $2^{-37}$ & $1/1024$\\
Translated transverse child & $2^{-36}$ & $2^{-37}$ & $1/1024$\\
Pivot $10\bmod16O$, others $0\bmod16O$ & $0$ & $0$ & $0$\\
One transverse digit fixed & $2^{-28}$ & $2^{-29}$ & $1/4$\\\bottomrule
\end{tabular}\end{center}
The common ratio on positive rows does not make every row an equally discriminating test. A zero remains zero under any positive rescaling. Its purpose is to test the support and location of the chart. Comparing two nonzero boxes tests both the overall constant and how the vector restrictions alter the result. Conservation over all the active children tests a different aspect again: a chart that correctly counted one branch but covered it more than once would fail the partition identity.

\section{A direct calculation in the readable parameters}
The graph calculation predicts what a compatible skew measure must assign to a small box. We can ask a stronger question. Take the formula that the readable source actually pushes forward, fix the Gaussian and skew coordinate Haars, and measure the preimage of the vector box in those parameters. This gives an independent numerator in the readable coordinates. It works at a single target, so no choice of a representative of an almost-everywhere conditional measure enters the argument.

Write $B=O+Oi$, and keep the basis $e,f,z$ constructed above. Denote its Gaussian coefficients by $B_0=1+i$ and $D_0=c+2i$ to distinguish the coefficient $B_0$ from the order $B$. Put $b=X-p/2$, where $p$ is the readable pivot, $X$ the solved hyperbolic coordinate, and $v$ the residual Gaussian coordinate. The physical vector $V=M(X,p,v)$ has the especially simple expression
\[
 V_1=p+b,\qquad V_2=B_0 b+\overline{D_0}v,\qquad
 V_3=D_0 b-\overline{B_0}v.
\]
This substitution uses the half inside the frame before asking which variables are integral. Its lower two-by-two matrix is integral and has determinant
$-N(B_0)-N(D_0)=1$. Its inverse is its adjugate, also integral. Consequently, for every width $W=2^m$ and every real center $a$ one has the exact equivalence
\[
 V\in (a+WB)\times WB\times WB
 \quad\Longleftrightarrow\quad
 p\in a+WB,\quad b\in WB,\quad v\in WB.
\]
For the reverse inclusion substitute into the three displayed coordinates. For the forward inclusion apply the integral inverse to $(V_2,V_3)$, then recover $p=V_1-b$. This use of the lower minor is stronger than a statement about the modulus of the full determinant. It identifies the preimage of the particular coordinate box the experiment observes.

\subsection{The skew coordinate occupies a wider scalar box}
Let $m\ge3$ and $a\in 2+8O$. Then $p=2U$ with $U$ a unit in $B$ whenever $p\in a+WB$. Indeed $U$ has unit real coefficient and imaginary coefficient in $4O$. Its Gaussian norm is a unit of $O$ and $\bar U/N(U)$ belongs to $B$. For $v\in WB$ and $s$ real, the readable flat chart, with residual coefficient $-1$, gives
\[
 b=p(r+i\alpha),\qquad
 r=\frac{s+N(v)-N(p)}{2N(p)}.
\]
The sign of $N(v)$ is positive here because the residual form is negative. Multiplication by $U$ preserves the entire order in both directions. Hence $b\in WB$ exactly when both $r$ and $\alpha$ belong to $(W/2)O$.

Write $p=a+W(A+iD)$ and $v=W(L+iH)$. Squaring and adding real and imaginary coordinates gives
\[
 N(p)=a^2+2aWA+W^2(A^2+D^2),\qquad
 N(v)=W^2(L^2+H^2).
\]
Since $a\in2+8O$, the difference $N(p)-a^2$ is divisible by $4W$. The same holds for $N(v)$ because $m\ge3$. The denominator $2N(p)=8N(U)$ is eight times a unit. Therefore
\[
 r\in (W/2)O\quad\Longleftrightarrow\quad s-a^2\in4WO.
\]
The equivalence is independent of the two Gaussian parameters after their box conditions have been specified. When $a=2$, $W=4d$, and $s=4+16dt$, the expression becomes
\[
 r=2d\,\frac{t-A+d(L^2+H^2-A^2-D^2)}{N(U)} .
\]
The rational identity and the equivalence under rotation by an integral Gaussian unit are separate Lean lemmas. This keeps the cancellation and the arithmetic unit hypothesis individually inspectable.

\subsection{Box preimage and volume}
For $s\in a^2+4WO$, the readable parameter preimage of the physical vector box is precisely
\[
    (a+WB)\ \times\ WB\ \times\ i(W/2)O .
\]
For $s$ in the complement it is empty. These statements concern the full preimage of the vector box. A parameter with nonunit pivot cannot make a second contribution elsewhere in the domain: membership in the physical box already forces $p\in a+WB$. The ring inverse at nonunits in the original definition introduces no extraneous mass.

The two Gaussian parameter volumes are each $4^{-2m}$. The skew parameter carries $2j_*\lambda$, so the last factor has mass $2\cdot4^{-(m-1)}$. Multiplication gives the split mass
\[
   (4^{-2m})^2\,2\,4^{-(m-1)}=8\,4^{-5m}.
\]
The original scalar normalization gives $16\,4^{-5m}$. For $m>4$ the same scalar value follows directly by taking the centered pivot cell and its full transverse box: for $w\in(2^mO)^5$ and target four, the squared transverse sum belongs to $2^{2m}O$. In the parent the pivot squaring map shifts every precision by two. Thus the root differs from two by an element of $2^{2m-2}O\subseteq2^mO$. Every such transverse point occurs exactly once. The closed-form numbers for all $m\ge3$ are written consequences of these arguments. The named literal rational endpoints in the old Lean module treat $m=3,4$; the general rational identity on the readable side treats every $m\ge1$.

This calculation is direct for a prescribed product of parameter Haars in the defining flat-chart formula. It supplies a pointwise comparison on these boxes. It also explains which use of the factor two would be circular. If one defined the readable measure by halving the old mass, there would be nothing left to check. Here its three measured sides come separately from the two Gaussian coordinates, the resolved skew width and the independently determined skew Haar. Constructing the entire upstream ambient-place package, which carries additional global conditions on the character and kernel, remains a distinct obligation.

\section{Pushing the parameter cells forward}
To test the rectangle in the preceding derivation, we evaluated the rational chart directly and asked which original vector children it reaches. This follow-up was designed after the first six-square counts. Its purpose is to discriminate between coordinate identifications that agree on a total mass; it supplies no blind prediction of the original outcome.

Inside an active physical box of width $W$, refine both Gaussian variables by one coefficient digit and the skew scalar by one coefficient digit. Each scalar digit has four possibilities. There are five scalar parameters and therefore $4^5=1024$ input cells. The original vector has six scalar coordinates and $4096$ possible children. To compute the image, first construct $c$ modulo the working precision from $u\mapsto-1-2u^2$ and $c=1+4u$. Next invert $N(U)$ as a unit modulo the precision. For a scalar coefficient $x+y\pi$, this inversion uses $(x-y\pi)/(x^2-2y^2)$, where the scalar dyadic denominator is odd. Evaluate $b=p(r+i\alpha)$, take the three physical coordinates above, and record their residues modulo $2W$. We verify the resulting target modulo $8W$ before collecting each image. This computation uses the hyperbolic change of variables and the skew input, neither of which appears in the direct six-square census.

Why does one representative suffice for these cells? Changing $p$ or $v$ by $2WB$ or $\alpha$ by $WO$ changes $b$ by a multiple of $2W$ in $B$. For $r$ this assertion can be seen without differentiating. A change $p\mapsto p+2W\eta$ changes $N(p)$ by an element of $8WO$, hence changes $N(U)=N(p)/4$ by an element of $2WO$. A change $v\mapsto v+2W\xi$ changes $N(v)$ by an element of $4W^2O$. Writing $r=(s+N(v))/(8N(U))-1/2$ and observing that $s+N(v)\in4O$ on these active targets shows $\Delta r\in WO$. Multiplication by $p=2U$ gives the required $2W$ bound. In the other direction, changing a physical scalar by $2W$ in the specified parent changes its square by an element of $8WO$. Thus the target class and the vector class used by the two enumerators are both well defined at the declared depths.

All four targets in the first parent gave exactly 1024 distinct original children. The centered box modulo sixteen at four and the pivot-ten box modulo sixteen at thirty-six gave the same count at the next refinement. For each of the six positive cases the two independently generated \emph{sets} of children agreed, in all six physical coordinates, without repetitions. At four the pivot-ten box had zero hits. For the deliberately mis-scaled control we narrowed only the skew parameter by a further factor of two. All 1024 parameter samples still landed on the required equation, yet only 256 distinct children remained. Four inputs now shared each image at the tested precision. Checking the equation alone would have missed this failure completely.

The combined evidence has three levels. The original-direction experiment measures the mass in fixed physical coordinates. The inverse image in readable parameters gives an analytic formula for the same physical boxes and a second mass computation. The direct parameter census verifies the image set against the physical census at multiple moduli. None of these requires assigning a global readable disintegration outside the patch. Together they answer the question that a whole-lattice ratio leaves open: the normalization persists on individual regular pieces, and the size of the free skew coordinate explains exactly how the readable chart realizes it.

\section{Which alternatives the evidence rules out}
One possible explanation for the earlier whole-lattice factor was a discrepancy affecting only the divisible vectors that dominate a particular congruence class. In that scenario sufficiently small regular patches might agree after a pure change of label. Here the mismatch occurs on a chart with a single root for each transverse vector and an explicit nonzero slope. The positive child and its translated counterpart both retain the factor two. Divisibility does help determine which vector sets meet the fiber, but it cannot account for the normalization discrepancy on this regular sheet.

Another possibility was that the trace-one failure prevented any comparison with the readable construction in these coordinates. The field-level trace identity and the algebraic specialization of the readable formula test that proposition directly. Once the Hermitian form has been put in the identified frame and the pivot denominator is nonzero, the cancellation is valid. The integral ring still lacks its trace-one element. This distinction can be tested without assigning a probability to the correctness of the readable global proof.

A third possibility concerned equality at one point. Agreement of fine tube averages could persist even if one family were modified on a target singleton. The explicit graph and a continuous family of test integrals fix the representative at a regular target. In normalized variables the full Borel change of variables establishes an integral statement. The inverse graph then defines a particular surface measure rather than an arbitrary almost-everywhere version. If comparing a different independently supplied kernel, one must still show it uses the same source and target normalizations and satisfies its regularity requirements. Equality on all target-neighborhood integrals alone would identify it only almost everywhere.

\section{A source map for independent review}
Four objects should be checked in order. The inherited arithmetic statement identifies the whole lattice in the original completion and provides the scalar Haar convention. The previous Haar calculation identifies $j(s)=is$ and determines the skew normalization forced by the readable split. The present direct-coordinate module supplies a homeomorphism of actual compact coordinate boxes plus a Haar calculation on new tests. The readable-edition module compares the field formula with the pinned \id{Place.flatChart} and evaluates it with its own Gaussian residual form. These steps use separately versioned sources. Typechecking one of them cannot automatically establish a theorem in the other Lean environment.

The dependency that matters at the source boundary is the difference between \id{Place.form_flatChart}, \id{Place.isDisintegration_quadricMeasure} and \id{IntegralPlace}. The first checks the equation of a chart point at a unit pivot. The second identifies a global family with a source measure and needs the full place structure and analytic argument. The third provides a different collection of hypotheses for integral estimates. Our restricted measure is built directly on the explicit coordinate patch. The preceding preimage calculation also tests the prescribed readable parameter pushforward on these boxes. Neither construction supplies the full upstream ambient-place data.

The source pins in the companion register give immutable file identities. Reviewers interested primarily in the geometry can begin with the determinant and chart equation and then check the finite controls. Reviewers interested in measure theory can begin with the isometry, the compact-open generating family and the physical slab identity. Either route should lead to the same numerical pair and the same explicit restriction on its interpretation.

\section{What remains open at this interface}
The result identifies regular pieces of a particular fiber with more precision than a total lattice mass. It gives a determinant-one-in-modulus transition to a hyperbolic form and a chart on a proper subset of the original coordinate lattice. The remaining global problem is not to guess another factor for this subset. It is to exhibit independent readable ambient-place data, use the actual upstream disintegration with those data, and identify it on a determining family of vector tests. For a full integral theorem one would also have to replace or bypass its documented trace-one hypothesis in a way justified by the theorem's downstream uses.

An immediate local comparison beyond this paper would choose a second patch with a different pivot, examine a nonempty overlap and compute the transition density there. The strongest test is a region where both pivots are regular, the two parameterizations present different transverse variables, and the observed vector subset has a convenient description in the original coordinates. A successful overlap identity would support gluing measures across charts. It would be distinct from checking another congruence class within the same sheet.

For target four the inference is already more limited and more definite. The direct preimage calculation identifies the measure from the readable formula on positive and empty vector boxes, with the specified product Haars. The observed factor two survives positive refinements, and the exact location of the inactive children follows from a single residue of the pivot. Neither a global disintegration nor an assertion about the zeros of an $L$-function is needed to test these local claims.

\section{Numerical witnesses for the resolving digit}
The entries below are counts of residue vectors, before changing the target measure. Two centered transverse boxes have identical widths. Their pivot centers differ by eight. At coefficient depth five the targets four and thirty-six agree modulo $32O$. The zero in the next row is the first observation that can distinguish them. Independent shifts of the $\sqrt2$ coefficient exercise the second coefficient of the residue-ring calculation; here the quotient $O/2O$ has four elements even though the residue field $O/\pi O$ has only two.

\begin{center}\small\begin{tabular}{@{}rrrr@{}}\toprule
Depth & Center child at 4 & Pivot 10 at 4 & Pivot 10 at 36\\\midrule
4 & 1 & 1 & 1\\
5 & 4,096 & 4,096 & 4,096\\
6 & 16,777,216 & 0 & 16,777,216\\
7 & 17,179,869,184 & 0 & 17,179,869,184\\
8 & 17,592,186,044,416 & 0 & 17,592,186,044,416\\
\bottomrule\end{tabular}\end{center}

A reviewer can reconstruct the entire first column of counts without iterating over field elements. At depth four there is one representative of the centered box and its square sum is four modulo sixteen. Increasing the vector precision by one coefficient digit adds four possibilities per coordinate. Until the target distinguishes a child, all of these choices survive, so the count grows by $4^6$. Once the image is resolved, a new target digit removes one factor of four. The growth then becomes $4^5$. Dividing the resulting counts by $4^{5d}$ gives a stable fiber value. The translated box has the same ambient size at each precision; its mass disappears only because its target residue moves. This pair of experiments can distinguish a defect in counting the ambient cells from a defect in computing the image.

\subsection{Which parameter is free?}
There are two inversions in the paper. In the original coordinates the free data are $w_1,\dots,w_5$ and a target; the contraction computes a single scalar $x$. In the readable coordinates the free data for a fixed target are a Gaussian pivot $p$, a Gaussian residual coordinate $v$ and one skew scalar $a$. They give two plus two plus one scalar degrees of freedom. The first inversion is nonlinear but integral throughout the chosen box. The second is rational and uses $1/2$ and the inverse of $\bar p$. These formulas describe the same equation in the frame, but the five free variables are different. A determinant of modulus one for the change of basis controls ambient Haar volume. It does not identify the two five-variable parameter measures by inspection.

One can recover the last free parameter directly at a fixed invertible pivot. Write $q=\bar v T v$ and $X=\bar p^{-1}\tau(s-q)+pa$. Given a vector satisfying the Hermitian equation, subtract the particular solution and divide by $p$. For $a=p^{-1}(X-\bar p^{-1}\tau(s-q))$, multiplying $a+\bar a$ by the unit $\bar p p$ gives $\bar p X+p\bar X-(\tau+\bar\tau)(s-q)=0$. Thus $a$ is skew, and cancellation proves uniqueness. Conversely every skew $a$ solves the equation. This short argument explains the one real degree of freedom which remains after solving a Hermitian equation in a Gaussian coordinate. The readable-side theorem below checks the forward substitution for arbitrary allowed inputs. The direct-coordinate homeomorphism supplies both inverse identities on the box.

\subsection{Which statement transports between editions?}
The square root witness can be constructed twice. In the old coefficient ring, solving $u+2u^2=-1$ gives $(1+4u)^2=-7$. On the readable side the polynomial $X^2+7$ has residual value eight and derivative two at one. Their norms satisfy $1/8<1/4$, so Hensel's strict inequality applies already over $\mathbb Z_2$. Extending scalars gives a witness in the readable-side coefficient field. The frame identities hold for every such root. They do not depend on proving that two particular choices of root returned by two libraries have identical representations. This distinction avoids a spurious cross-version proof obligation while retaining the same mathematical field and equation.

\section{Annotated formal evidence}
The excerpts give line numbers in the accompanying sources. Each subsection identifies a proof obligation which could fail independently of the others. The complete files and receipts in the package are the authoritative machine inputs. For visual readability, the mathematical glyphs of Lean are typeset in place.
\subsection{Invertibility on an entire integral ball}
The isometry proof factors a difference and uses a strict inequality of norms. Surjectivity needs completeness in addition. Following these two proofs separately exposes exactly where the existence of an infinite limit enters.\par\smallskip\id{FourChart.lean}\sourceblock
\sourceline{24}{def pivot (x : D) : D := x+2*x\textasciicircum{}2}
\sourceline{26}{theorem pivot\_\allowbreak{}isometry : Isometry pivot := by}
\sourceline{27}{\hspace*{1.0em}apply Isometry.\allowbreak{}of\_\allowbreak{}dist\_\allowbreak{}eq}
\sourceline{28}{\hspace*{1.0em}intro x y}
\sourceline{29}{\hspace*{1.0em}change dist (x.\allowbreak{}val+2*x.\allowbreak{}val\textasciicircum{}2) (y.\allowbreak{}val+2*y.\allowbreak{}val\textasciicircum{}2) = dist x.\allowbreak{}val y.\allowbreak{}val}
\sourceline{30}{\hspace*{1.0em}simp only [dist\_\allowbreak{}eq\_\allowbreak{}norm]}
\sourceline{31}{\hspace*{1.0em}rw [show x.\allowbreak{}val+2*x.\allowbreak{}val\textasciicircum{}2-(y.\allowbreak{}val+2*y.\allowbreak{}val\textasciicircum{}2)=(x.\allowbreak{}val-y.\allowbreak{}val)*(1+2*(x.\allowbreak{}val+y.\allowbreak{}val)) by ring,norm\_\allowbreak{}mul]}
\sourceline{32}{\hspace*{1.0em}have hlt : \ensuremath{\Vert}(2:K)*(x.\allowbreak{}val+y.\allowbreak{}val)\ensuremath{\Vert} < 1 := by rw [norm\_\allowbreak{}mul,norm\_\allowbreak{}two]; linarith [sum\_\allowbreak{}bound x y]}
\sourceline{33}{\hspace*{1.0em}rw [IsUltrametricDist.\allowbreak{}norm\_\allowbreak{}add\_\allowbreak{}eq\_\allowbreak{}max\_\allowbreak{}of\_\allowbreak{}norm\_\allowbreak{}ne\_\allowbreak{}norm (by rw [norm\_\allowbreak{}one]; exact ne\_\allowbreak{}of\_\allowbreak{}gt hlt),norm\_\allowbreak{}one,max\_\allowbreak{}eq\_\allowbreak{}left hlt.\allowbreak{}le,mul\_\allowbreak{}one]}
\sourceline{35}{def iteration (b x : D) : D := b-2*x\textasciicircum{}2}
\sourceline{37}{theorem iteration\_\allowbreak{}contracts (b : D) : ContractingWith (1/2:\ensuremath{\mathbb R}\ensuremath{\ge}0) (iteration b) := by}
\sourceline{38}{\hspace*{1.0em}refine \ensuremath{\langle}by norm\_\allowbreak{}num, LipschitzWith.\allowbreak{}of\_\allowbreak{}dist\_\allowbreak{}le\_\allowbreak{}mul ?\_\allowbreak{}\ensuremath{\rangle}}
\sourceline{39}{\hspace*{1.0em}intro x y}
\sourceline{40}{\hspace*{1.0em}change dist (b.\allowbreak{}val-2*x.\allowbreak{}val\textasciicircum{}2) (b.\allowbreak{}val-2*y.\allowbreak{}val\textasciicircum{}2) \ensuremath{\le} ((1/2:\ensuremath{\mathbb R}\ensuremath{\ge}0):\ensuremath{\mathbb R})*dist x.\allowbreak{}val y.\allowbreak{}val}
\sourceline{41}{\hspace*{1.0em}rw [dist\_\allowbreak{}eq\_\allowbreak{}norm,dist\_\allowbreak{}eq\_\allowbreak{}norm,show b.\allowbreak{}val-2*x.\allowbreak{}val\textasciicircum{}2-(b.\allowbreak{}val-2*y.\allowbreak{}val\textasciicircum{}2)= -(2*(x.\allowbreak{}val+y.\allowbreak{}val))*(x.\allowbreak{}val-y.\allowbreak{}val) by ring,norm\_\allowbreak{}mul,norm\_\allowbreak{}neg,norm\_\allowbreak{}mul,norm\_\allowbreak{}two]}
\sourceline{42}{\hspace*{1.0em}push\_\allowbreak{}cast}
\sourceline{43}{\hspace*{1.0em}exact mul\_\allowbreak{}le\_\allowbreak{}mul\_\allowbreak{}of\_\allowbreak{}nonneg\_\allowbreak{}right (by linarith [sum\_\allowbreak{}bound x y]) (norm\_\allowbreak{}nonneg \_\allowbreak{})}
\sourceline{45}{theorem pivot\_\allowbreak{}surjective : Function.\allowbreak{}Surjective pivot := by}
\sourceline{46}{\hspace*{1.0em}let : CompleteSpace D := integers\_\allowbreak{}compact.\allowbreak{}isComplete.\allowbreak{}completeSpace\_\allowbreak{}coe}
\sourceline{47}{\hspace*{1.0em}let : Nonempty D := \ensuremath{\langle}\ensuremath{\langle}0,(mem\_\allowbreak{}D \_\allowbreak{}).\allowbreak{}mpr (by simp)\ensuremath{\rangle}\ensuremath{\rangle}}
\sourceline{48}{\hspace*{1.0em}intro b}
\sourceline{49}{\hspace*{1.0em}let c := iteration\_\allowbreak{}contracts b}
\sourceline{50}{\hspace*{1.0em}let u := ContractingWith.\allowbreak{}fixedPoint (iteration b) c}
\sourceline{51}{\hspace*{1.0em}have hu := congrArg Subtype.\allowbreak{}val (c.\allowbreak{}fixedPoint\_\allowbreak{}isFixedPt (f:=iteration b))}
\sourceline{52}{\hspace*{1.0em}change b.\allowbreak{}val-2*u.\allowbreak{}val\textasciicircum{}2=u.\allowbreak{}val at hu}
\sourceline{53}{\hspace*{1.0em}exact \ensuremath{\langle}u,Subtype.\allowbreak{}ext (by change u.\allowbreak{}val+2*u.\allowbreak{}val\textasciicircum{}2=b.\allowbreak{}val; linear\_\allowbreak{}combination -hu)\ensuremath{\rangle}}
\sourceline{55}{def pivotEquiv : D \ensuremath{\simeq}\ensuremath{_{i}} D :=}
\sourceend
\subsection{The original form and the joint inverse}
The scalar polynomial is first related to the original six-index expression. The normalized map then keeps the five variables and solves the target. Observe that the inverse depends continuously on both the target and all transverse coordinates.\par\smallskip\id{FourChart.lean}\sourceblock
\sourceline{108}{theorem original\_\allowbreak{}gram (x : K) (w : W) : TargetOne.\allowbreak{}Q (six x w) = gram x w := by}
\sourceline{109}{\hspace*{1.0em}simp [TargetOne.\allowbreak{}Q,gram,six,Fintype.\allowbreak{}sum\_\allowbreak{}prod\_\allowbreak{}type,Fin.\allowbreak{}sum\_\allowbreak{}univ\_\allowbreak{}succ]}
\sourceline{110}{\hspace*{1.0em}ring}
\sourceline{112}{def offset (w : WD) : D := 2*\ensuremath{\sum} j, w j\textasciicircum{}2}
\sourceline{114}{def normalized : D \ensuremath{\times} WD \ensuremath{\simeq}\ensuremath{_{t}} D \ensuremath{\times} WD where}
\sourceline{115}{\hspace*{1.0em}toFun p := (pivot p.\allowbreak{}1 + offset p.\allowbreak{}2,p.\allowbreak{}2)}
\sourceline{116}{\hspace*{1.0em}invFun p := (pivotEquiv.\allowbreak{}symm (p.\allowbreak{}1-offset p.\allowbreak{}2),p.\allowbreak{}2)}
\sourceline{117}{\hspace*{1.0em}left\_\allowbreak{}inv := by}
\sourceline{118}{\hspace*{2.0em}rintro \ensuremath{\langle}x,w\ensuremath{\rangle}; simp only [add\_\allowbreak{}sub\_\allowbreak{}cancel\_\allowbreak{}right]; exact Prod.\allowbreak{}ext (pivotEquiv.\allowbreak{}symm\_\allowbreak{}apply\_\allowbreak{}apply \_\allowbreak{}) rfl}
\sourceline{119}{\hspace*{1.0em}right\_\allowbreak{}inv := by}
\sourceline{120}{\hspace*{2.0em}rintro \ensuremath{\langle}t,w\ensuremath{\rangle}}
\sourceline{121}{\hspace*{2.0em}change (pivotEquiv (pivotEquiv.\allowbreak{}symm (t-offset w))+offset w,w)=(t,w)}
\sourceline{122}{\hspace*{2.0em}rw [pivotEquiv.\allowbreak{}apply\_\allowbreak{}symm\_\allowbreak{}apply,sub\_\allowbreak{}add\_\allowbreak{}cancel]}
\sourceline{123}{\hspace*{1.0em}continuous\_\allowbreak{}toFun := by}
\sourceline{124}{\hspace*{2.0em}have hp := pivot\_\allowbreak{}isometry.\allowbreak{}continuous}
\sourceline{125}{\hspace*{2.0em}unfold offset}
\sourceline{126}{\hspace*{2.0em}fun\_\allowbreak{}prop}
\sourceline{127}{\hspace*{1.0em}continuous\_\allowbreak{}invFun := by}
\sourceline{128}{\hspace*{2.0em}have hp := pivotEquiv.\allowbreak{}symm.\allowbreak{}continuous}
\sourceline{129}{\hspace*{2.0em}unfold offset}
\sourceline{130}{\hspace*{2.0em}fun\_\allowbreak{}prop}
\sourceline{132}{def physical (p : D \ensuremath{\times} WD) : K \ensuremath{\times} W := (2+8*p.\allowbreak{}1.\allowbreak{}val,fun j => 8*(p.\allowbreak{}2 j).\allowbreak{}val)}
\sourceline{133}{def scalarTarget (t : D) : K := 4+32*t.\allowbreak{}val}
\sourceline{135}{theorem chart\_\allowbreak{}equation (p : D \ensuremath{\times} WD) :}
\sourceline{136}{\hspace*{2.0em}TargetOne.\allowbreak{}Q (six (physical p).\allowbreak{}1 (physical p).\allowbreak{}2) = scalarTarget (normalized p).\allowbreak{}1 := by}
\sourceline{137}{\hspace*{1.0em}rw [original\_\allowbreak{}gram]}
\sourceline{138}{\hspace*{1.0em}rcases p with \ensuremath{\langle}u,w\ensuremath{\rangle}}
\sourceline{139}{\hspace*{1.0em}change (2+8*u.\allowbreak{}val)\textasciicircum{}2 + \ensuremath{\sum} j, (8*(w j).\allowbreak{}val)\textasciicircum{}2 = 4 + 32*((u+2*u\textasciicircum{}2+offset w:D):K)}
\sourceline{140}{\hspace*{1.0em}simp [offset, Fin.\allowbreak{}sum\_\allowbreak{}univ\_\allowbreak{}succ, show ((2:D):K) = 2 from rfl]; ring}
\sourceline{142}{def graph (t : D) (w : WD) := physical (normalized.\allowbreak{}symm (t,w))}
\sourceline{144}{theorem graph\_\allowbreak{}skew (t : D) (w : WD) :}
\sourceline{145}{\hspace*{2.0em}TargetFourHaar.\allowbreak{}skewBorel (TargetOne.\allowbreak{}Q (six (graph t w).\allowbreak{}1 (graph t w).\allowbreak{}2)) =}
\sourceline{146}{\hspace*{2.0em}TargetFourHaar.\allowbreak{}skewBorel (scalarTarget t) := by}
\sourceline{147}{\hspace*{1.0em}unfold graph}
\sourceline{148}{\hspace*{1.0em}rw [chart\_\allowbreak{}equation,normalized.\allowbreak{}apply\_\allowbreak{}symm\_\allowbreak{}apply]}
\sourceend
\subsection{The physical slope}
In normalized coordinates the pivot is an isometry. The physical slope has a different modulus. This theorem proves the metric statement with the actual distances before either Haar measure is introduced.\par\smallskip\id{FourChart.lean}\sourceblock
\sourceline{154}{theorem delta\_\allowbreak{}norm \{x y : K\} (hx : x \ensuremath{\in} cell 3 2) (hy : y \ensuremath{\in} cell 3 2) :}
\sourceline{155}{\hspace*{2.0em}\ensuremath{\Vert}x\textasciicircum{}2-y\textasciicircum{}2\ensuremath{\Vert}=(1/4:\ensuremath{\mathbb R})*\ensuremath{\Vert}x-y\ensuremath{\Vert} := by}
\sourceline{156}{\hspace*{1.0em}rw [cell\_\allowbreak{}image] at hx hy}
\sourceline{157}{\hspace*{1.0em}obtain \ensuremath{\langle}u,\_\allowbreak{},rfl\ensuremath{\rangle} := hx}
\sourceline{158}{\hspace*{1.0em}obtain \ensuremath{\langle}v,\_\allowbreak{},rfl\ensuremath{\rangle} := hy}
\sourceline{159}{\hspace*{1.0em}have h := pivot\_\allowbreak{}isometry.\allowbreak{}dist\_\allowbreak{}eq u v}
\sourceline{160}{\hspace*{1.0em}change dist (u.\allowbreak{}val+2*u.\allowbreak{}val\textasciicircum{}2) (v.\allowbreak{}val+2*v.\allowbreak{}val\textasciicircum{}2)=dist u.\allowbreak{}val v.\allowbreak{}val at h}
\sourceline{161}{\hspace*{1.0em}simp only [dist\_\allowbreak{}eq\_\allowbreak{}norm] at h}
\sourceline{162}{\hspace*{1.0em}rw [show (2+(2:K)\textasciicircum{}3*u.\allowbreak{}val)\textasciicircum{}2-(2+2\textasciicircum{}3*v.\allowbreak{}val)\textasciicircum{}2=2\textasciicircum{}5*((u.\allowbreak{}val+2*u.\allowbreak{}val\textasciicircum{}2)-(v.\allowbreak{}val+2*v.\allowbreak{}val\textasciicircum{}2)) by ring,}
\sourceline{163}{\hspace*{2.0em}show (2+(2:K)\textasciicircum{}3*u.\allowbreak{}val)-(2+2\textasciicircum{}3*v.\allowbreak{}val)=2\textasciicircum{}3*(u.\allowbreak{}val-v.\allowbreak{}val) by ring,}
\sourceline{164}{\hspace*{2.0em}norm\_\allowbreak{}mul,norm\_\allowbreak{}scale,h,norm\_\allowbreak{}mul,norm\_\allowbreak{}scale]}
\sourceline{165}{\hspace*{1.0em}ring}
\sourceend
\subsection{Locating the unique root modulo sixteen}
In the parent every pivot has the same residue modulo eight. At four the graph imposes a stronger condition. The source proof locates the root using the defining equation of the normalized inverse, so the next digit is constrained even when the transverse center is translated.\par\smallskip\id{FourChart.lean}\sourceblock
\sourceline{209}{theorem root\_\allowbreak{}formula (t : D) (z : WD) :}
\sourceline{210}{\hspace*{2.0em}gram (graph t z).\allowbreak{}1 (fun j => 8*(z j).\allowbreak{}val) = scalarTarget t := by}
\sourceline{211}{\hspace*{1.0em}have h := chart\_\allowbreak{}equation (normalized.\allowbreak{}symm (t,z))}
\sourceline{212}{\hspace*{1.0em}rw [normalized.\allowbreak{}apply\_\allowbreak{}symm\_\allowbreak{}apply,original\_\allowbreak{}gram] at h}
\sourceline{213}{\hspace*{1.0em}exact h}
\sourceline{215}{theorem root\_\allowbreak{}box (t : D) (z : WD) : (graph t z).\allowbreak{}1 \ensuremath{\in} cell 3 2 := by}
\sourceline{216}{\hspace*{1.0em}rw [cell\_\allowbreak{}image]}
\sourceline{217}{\hspace*{1.0em}exact \ensuremath{\langle}(normalized.\allowbreak{}symm (t,z)).\allowbreak{}1,mem\_\allowbreak{}univ \_\allowbreak{},by norm\_\allowbreak{}num [graph,physical]\ensuremath{\rangle}}
\sourceline{219}{theorem root\_\allowbreak{}child (z : WD) : (graph 0 z).\allowbreak{}1 \ensuremath{\in} cell 4 2 := by}
\sourceline{220}{\hspace*{1.0em}let u := pivotEquiv.\allowbreak{}symm (0-offset z)}
\sourceline{221}{\hspace*{1.0em}have hu := pivotEquiv.\allowbreak{}isometry.\allowbreak{}dist\_\allowbreak{}eq u 0}
\sourceline{222}{\hspace*{1.0em}have he : pivotEquiv u = -offset z := by rw [show pivotEquiv u = \_\allowbreak{} from pivotEquiv.\allowbreak{}apply\_\allowbreak{}symm\_\allowbreak{}apply \_\allowbreak{},zero\_\allowbreak{}sub]}
\sourceline{223}{\hspace*{1.0em}rw [he,show pivotEquiv 0 = 0 by change (0:D)+2*0\textasciicircum{}2=0; ring] at hu}
\sourceline{224}{\hspace*{1.0em}change dist (-(offset z):D).\allowbreak{}val (0:K) = dist u.\allowbreak{}val 0 at hu}
\sourceline{225}{\hspace*{1.0em}rw [dist\_\allowbreak{}zero\_\allowbreak{}right,dist\_\allowbreak{}zero\_\allowbreak{}right] at hu}
\sourceline{226}{\hspace*{1.0em}have hh : \ensuremath{\Vert}u.\allowbreak{}val\ensuremath{\Vert} \ensuremath{\le} 1/2 := by}
\sourceline{227}{\hspace*{2.0em}rw [\ensuremath{\leftarrow}hu]}
\sourceline{228}{\hspace*{2.0em}change \ensuremath{\Vert}-((2:K)*((\ensuremath{\sum} j, z j\textasciicircum{}2 : D):K))\ensuremath{\Vert} \ensuremath{\le} \_\allowbreak{}}
\sourceline{229}{\hspace*{2.0em}rw [norm\_\allowbreak{}neg,norm\_\allowbreak{}mul,norm\_\allowbreak{}two]}
\sourceline{230}{\hspace*{2.0em}linarith [(mem\_\allowbreak{}D \_\allowbreak{}).\allowbreak{}mp (\ensuremath{\sum} j, z j\textasciicircum{}2 : D).\allowbreak{}2]}
\sourceline{231}{\hspace*{1.0em}change \ensuremath{\Vert}2+8*u.\allowbreak{}val-2\ensuremath{\Vert} \ensuremath{\le} (1/2:\ensuremath{\mathbb R})\textasciicircum{}4}
\sourceline{232}{\hspace*{1.0em}rw [add\_\allowbreak{}sub\_\allowbreak{}cancel\_\allowbreak{}left,norm\_\allowbreak{}mul,show (8:K)=2\textasciicircum{}3 by norm\_\allowbreak{}num,norm\_\allowbreak{}scale]}
\sourceline{233}{\hspace*{1.0em}norm\_\allowbreak{}num at *; linarith}
\sourceend
\subsection{Disintegrating an actual source slab}
The first result identifies a section as one exact ball in the pivot. The source mass comes from integrating these sections. The transverse set need only be measurable, and the input hypothesis explicitly places the root in the selected pivot cell.\par\smallskip\id{FourChart.lean}\sourceblock
\sourceline{237}{theorem slice \{m r : \ensuremath{\mathbb N}\} \{w : W\} \{y a s : K\} (hr : m \ensuremath{\le} r)}
\sourceline{238}{\hspace*{2.0em}(hy : y \ensuremath{\in} cell m a) (hsub : cell m a \ensuremath{\subseteq} cell 3 2) (hs : gram y w=s) :}
\sourceline{239}{\hspace*{2.0em}\{x : K | x \ensuremath{\in} cell m a \ensuremath{\wedge} gram x w \ensuremath{\in} cell (r+2) s\} = cell r y := by}
\sourceline{240}{\hspace*{1.0em}ext x}
\sourceline{241}{\hspace*{1.0em}have hv (hx : x \ensuremath{\in} cell m a) : gram x w \ensuremath{\in} cell (r+2) s \ensuremath{\leftrightarrow} x \ensuremath{\in} cell r y := by}
\sourceline{242}{\hspace*{2.0em}change \ensuremath{\Vert}gram x w-s\ensuremath{\Vert} \ensuremath{\le} \_\allowbreak{} \ensuremath{\leftrightarrow} \ensuremath{\Vert}x-y\ensuremath{\Vert} \ensuremath{\le} \_\allowbreak{}}
\sourceline{243}{\hspace*{2.0em}rw [\ensuremath{\leftarrow}hs,gram,gram,add\_\allowbreak{}sub\_\allowbreak{}add\_\allowbreak{}right\_\allowbreak{}eq\_\allowbreak{}sub,delta\_\allowbreak{}norm (hsub hx) (hsub hy),pow\_\allowbreak{}add]}
\sourceline{244}{\hspace*{2.0em}norm\_\allowbreak{}num}
\sourceline{245}{\hspace*{2.0em}constructor <;> intro h <;> linarith}
\sourceline{246}{\hspace*{1.0em}constructor}
\sourceline{247}{\hspace*{1.0em}\ensuremath{\cdot} rintro \ensuremath{\langle}hx,ht\ensuremath{\rangle}; exact (hv hx).\allowbreak{}mp ht}
\sourceline{248}{\hspace*{1.0em}\ensuremath{\cdot} intro hx}
\sourceline{249}{\hspace*{2.0em}have hh := nest hr hy hx}
\sourceline{250}{\hspace*{2.0em}exact \ensuremath{\langle}hh,(hv hh).\allowbreak{}mpr hx\ensuremath{\rangle}}
\sourceline{252}{def slab (S : Set W) (m : \ensuremath{\mathbb N}) (a : K) (r : \ensuremath{\mathbb N}) (s : K) : Set (W \ensuremath{\times} K) :=}
\sourceline{253}{\hspace*{1.0em}\{p | p.\allowbreak{}1 \ensuremath{\in} S \ensuremath{\wedge} p.\allowbreak{}2 \ensuremath{\in} cell m a \ensuremath{\wedge} gram p.\allowbreak{}2 p.\allowbreak{}1 \ensuremath{\in} cell (r+2) s\}}
\sourceline{255}{theorem slab\_\allowbreak{}meas \{S : Set W\} (hS : MeasurableSet S) (m : \ensuremath{\mathbb N}) (a : K) (r : \ensuremath{\mathbb N}) (s : K) :}
\sourceline{256}{\hspace*{2.0em}MeasurableSet (slab S m a r s) := by}
\sourceline{257}{\hspace*{1.0em}have h : Continuous (fun p : W \ensuremath{\times} K => gram p.\allowbreak{}2 p.\allowbreak{}1) := by unfold gram; fun\_\allowbreak{}prop}
\sourceline{258}{\hspace*{1.0em}exact (hS.\allowbreak{}preimage measurable\_\allowbreak{}fst).\allowbreak{}inter (((cell\_\allowbreak{}closed m a).\allowbreak{}measurableSet.\allowbreak{}preimage measurable\_\allowbreak{}snd).\allowbreak{}inter}
\sourceline{259}{\hspace*{2.0em}((cell\_\allowbreak{}closed (r+2) s).\allowbreak{}measurableSet.\allowbreak{}preimage h.\allowbreak{}measurable))}
\sourceline{261}{theorem slab\_\allowbreak{}mass (lam : Measure K) [lam.\allowbreak{}IsAddHaarMeasure] (hvol : lam integers=1)}
\sourceline{262}{\hspace*{2.0em}(kap : Measure W) \{S : Set W\} (hS : MeasurableSet S) \{m r : \ensuremath{\mathbb N}\} \{a s : K\}}
\sourceline{263}{\hspace*{2.0em}(hr : m \ensuremath{\le} r) (hsub : cell m a \ensuremath{\subseteq} cell 3 2)}
\sourceline{264}{\hspace*{2.0em}(hroot : \ensuremath{\forall} w \ensuremath{\in} S, \ensuremath{\exists} y \ensuremath{\in} cell m a, gram y w=s) :}
\sourceline{265}{\hspace*{2.0em}(kap.\allowbreak{}prod lam) (slab S m a r s) = kap S * (1/(2\textasciicircum{}r:\ensuremath{\mathbb R}\ensuremath{\ge}0\ensuremath{\infty})\textasciicircum{}2) := by}
\sourceline{266}{\hspace*{1.0em}rw [Measure.\allowbreak{}prod\_\allowbreak{}apply (slab\_\allowbreak{}meas hS .\allowbreak{}.\allowbreak{})]}
\sourceline{267}{\hspace*{1.0em}calc}
\sourceline{268}{\hspace*{2.0em}\_\allowbreak{} = \ensuremath{\int}\ensuremath{^{-}} w, S.\allowbreak{}indicator (fun \_\allowbreak{} : W => (1/(2\textasciicircum{}r:\ensuremath{\mathbb R}\ensuremath{\ge}0\ensuremath{\infty})\textasciicircum{}2)) w \ensuremath{\partial}kap := by}
\sourceline{269}{\hspace*{3.0em}apply lintegral\_\allowbreak{}congr}
\sourceline{270}{\hspace*{3.0em}intro w}
\sourceline{271}{\hspace*{3.0em}by\_\allowbreak{}cases hw : w \ensuremath{\in} S}
\sourceline{272}{\hspace*{3.0em}\ensuremath{\cdot} obtain \ensuremath{\langle}y,hy,hs\ensuremath{\rangle} := hroot w hw}
\sourceline{273}{\hspace*{4.0em}simpa [slab,Set.\allowbreak{}preimage,hw,slice hr hy hsub hs] using cell\_\allowbreak{}volume lam hvol r y}
\sourceline{274}{\hspace*{3.0em}\ensuremath{\cdot} simp [slab,Set.\allowbreak{}preimage,hw]}
\sourceline{275}{\hspace*{2.0em}\_\allowbreak{} = \_\allowbreak{} := by rw [lintegral\_\allowbreak{}indicator\_\allowbreak{}const hS,mul\_\allowbreak{}comm]}
\sourceend
\subsection{Why the normalized distribution is determined}
Closed balls of coefficient precision form a generating family for Borel sets in the integral ball. Their common volumes and the intersection rule allow uniqueness of the finite measures. This is the step from a test on balls to measurable distributions.\par\smallskip\id{FourBoxMeasures.lean}\sourceblock
\sourceline{33}{def discs : Set (Set D) := \{U | \ensuremath{\exists} n a, U=disc n a\}}
\sourceline{35}{theorem disc\_\allowbreak{}basis : TopologicalSpace.\allowbreak{}IsTopologicalBasis discs := by}
\sourceline{36}{\hspace*{1.0em}apply isTopologicalBasis\_\allowbreak{}of\_\allowbreak{}isOpen\_\allowbreak{}of\_\allowbreak{}nhds}
\sourceline{37}{\hspace*{1.0em}\ensuremath{\cdot} rintro \_\allowbreak{} \ensuremath{\langle}n,a,rfl\ensuremath{\rangle}}
\sourceline{38}{\hspace*{2.0em}exact IsUltrametricDist.\allowbreak{}isOpen\_\allowbreak{}closedBall a (by positivity)}
\sourceline{39}{\hspace*{1.0em}\ensuremath{\cdot} intro a U ha ho}
\sourceline{40}{\hspace*{2.0em}obtain \ensuremath{\langle}e,he,h\ensuremath{\rangle} := Metric.\allowbreak{}isOpen\_\allowbreak{}iff.\allowbreak{}mp ho a ha}
\sourceline{41}{\hspace*{2.0em}obtain \ensuremath{\langle}n,hn\ensuremath{\rangle} := exists\_\allowbreak{}pow\_\allowbreak{}lt\_\allowbreak{}of\_\allowbreak{}lt\_\allowbreak{}one he (by norm\_\allowbreak{}num : (1/2:\ensuremath{\mathbb R})<1)}
\sourceline{42}{\hspace*{2.0em}exact \ensuremath{\langle}disc n a,\ensuremath{\langle}n,a,rfl\ensuremath{\rangle},Metric.\allowbreak{}mem\_\allowbreak{}closedBall\_\allowbreak{}self (by positivity),}
\sourceline{43}{\hspace*{3.0em}(Metric.\allowbreak{}closedBall\_\allowbreak{}subset\_\allowbreak{}ball hn).\allowbreak{}trans h\ensuremath{\rangle}}
\sourceline{45}{theorem disc\_\allowbreak{}pi : IsPiSystem discs := by}
\sourceline{46}{\hspace*{1.0em}rintro \_\allowbreak{} \ensuremath{\langle}n,a,rfl\ensuremath{\rangle} \_\allowbreak{} \ensuremath{\langle}m,b,rfl\ensuremath{\rangle} h}
\sourceline{47}{\hspace*{1.0em}rcases IsUltrametricDist.\allowbreak{}closedBall\_\allowbreak{}subset\_\allowbreak{}trichotomy (x:=a) (y:=b)}
\sourceline{48}{\hspace*{2.0em}(r:=(1/2:\ensuremath{\mathbb R})\textasciicircum{}n) (s:=(1/2:\ensuremath{\mathbb R})\textasciicircum{}m) with hab|hba|hd}
\sourceline{49}{\hspace*{1.0em}\ensuremath{\cdot} exact \ensuremath{\langle}n,a,inter\_\allowbreak{}eq\_\allowbreak{}left.\allowbreak{}mpr hab\ensuremath{\rangle}}
\sourceline{50}{\hspace*{1.0em}\ensuremath{\cdot} exact \ensuremath{\langle}m,b,inter\_\allowbreak{}eq\_\allowbreak{}right.\allowbreak{}mpr hba\ensuremath{\rangle}}
\sourceline{51}{\hspace*{1.0em}\ensuremath{\cdot} exact False.\allowbreak{}elim ((Set.\allowbreak{}disjoint\_\allowbreak{}iff\_\allowbreak{}inter\_\allowbreak{}eq\_\allowbreak{}empty.\allowbreak{}mp hd \ensuremath{\triangleright} h).\allowbreak{}ne\_\allowbreak{}empty rfl)}
\sourceline{53}{theorem iso\_\allowbreak{}preserves (lam : Measure K) [lam.\allowbreak{}IsAddHaarMeasure] (hvol : lam integers=1) (e : D \ensuremath{\simeq}\ensuremath{_{i}} D) :}
\sourceline{54}{\hspace*{2.0em}MeasurePreserving e (integralHaar lam) (integralHaar lam) := by}
\sourceline{55}{\hspace*{1.0em}have hp : IsProbabilityMeasure (integralHaar lam) := \ensuremath{\langle}integral\_\allowbreak{}univ lam hvol\ensuremath{\rangle}}
\sourceline{56}{\hspace*{1.0em}let : IsProbabilityMeasure (integralHaar lam) := hp}
\sourceline{57}{\hspace*{1.0em}let mu := integralHaar lam}
\sourceline{58}{\hspace*{1.0em}have hmap := Measure.\allowbreak{}isProbabilityMeasure\_\allowbreak{}map (\ensuremath{\mu}:=mu) e.\allowbreak{}continuous.\allowbreak{}aemeasurable}
\sourceline{59}{\hspace*{1.0em}refine \ensuremath{\langle}e.\allowbreak{}continuous.\allowbreak{}measurable, ext\_\allowbreak{}of\_\allowbreak{}generate\_\allowbreak{}finite discs}
\sourceline{60}{\hspace*{2.0em}((BorelSpace.\allowbreak{}measurable\_\allowbreak{}eq (\ensuremath{\alpha}:=D)).\allowbreak{}trans disc\_\allowbreak{}basis.\allowbreak{}borel\_\allowbreak{}eq\_\allowbreak{}generateFrom) disc\_\allowbreak{}pi ?\_\allowbreak{} ?\_\allowbreak{}\ensuremath{\rangle}}
\sourceline{61}{\hspace*{1.0em}\ensuremath{\cdot} rintro \_\allowbreak{} \ensuremath{\langle}n,a,rfl\ensuremath{\rangle}}
\sourceline{62}{\hspace*{2.0em}rw [Measure.\allowbreak{}map\_\allowbreak{}apply e.\allowbreak{}continuous.\allowbreak{}measurable (show MeasurableSet (disc n a) from Metric.\allowbreak{}isClosed\_\allowbreak{}closedBall.\allowbreak{}measurableSet)]}
\sourceline{63}{\hspace*{2.0em}have h : e \ensuremath{^{-}}\ensuremath{^{1}}\textquotesingle{} disc n a = disc n (e.\allowbreak{}symm a) := by}
\sourceline{64}{\hspace*{3.0em}ext x}
\sourceline{65}{\hspace*{3.0em}change dist (e x) a \ensuremath{\le} \_\allowbreak{} \ensuremath{\leftrightarrow} dist x (e.\allowbreak{}symm a) \ensuremath{\le} \_\allowbreak{}}
\sourceline{66}{\hspace*{3.0em}rw [\ensuremath{\leftarrow}e.\allowbreak{}dist\_\allowbreak{}eq x (e.\allowbreak{}symm a),e.\allowbreak{}apply\_\allowbreak{}symm\_\allowbreak{}apply]}
\sourceline{67}{\hspace*{2.0em}rw [h,disc\_\allowbreak{}mass lam hvol,disc\_\allowbreak{}mass lam hvol]}
\sourceline{68}{\hspace*{1.0em}\ensuremath{\cdot} simp [Measure.\allowbreak{}map\_\allowbreak{}apply\_\allowbreak{}of\_\allowbreak{}aemeasurable e.\allowbreak{}continuous.\allowbreak{}aemeasurable]}
\sourceend
\subsection{Translations depending on the transverse vector}
The translation is a function of all transverse variables. A product argument must therefore prove invariance for each parameter and retain joint measurability. Marginal invariance alone would not control a test correlating the target with those variables.\par\smallskip\id{FourBoxMeasures.lean}\sourceblock
\sourceline{79}{theorem normalized\_\allowbreak{}preserves (lam : Measure K) [lam.\allowbreak{}IsAddHaarMeasure] (hvol : lam integers=1)}
\sourceline{80}{\hspace*{2.0em}(kap : Measure WD) [SFinite kap] :}
\sourceline{81}{\hspace*{2.0em}MeasurePreserving normalized ((integralHaar lam).\allowbreak{}prod kap) ((integralHaar lam).\allowbreak{}prod kap) := by}
\sourceline{82}{\hspace*{1.0em}let mu := integralHaar lam}
\sourceline{83}{\hspace*{1.0em}let : IsProbabilityMeasure mu := \ensuremath{\langle}integral\_\allowbreak{}univ lam hvol\ensuremath{\rangle}}
\sourceline{84}{\hspace*{1.0em}have hg (w : WD) : MeasurePreserving (fun x : D => pivot x+offset w) mu mu :=}
\sourceline{85}{\hspace*{2.0em}(iso\_\allowbreak{}preserves lam hvol (pivotEquiv.\allowbreak{}trans (translateIso (offset w))))}
\sourceline{86}{\hspace*{1.0em}have hm : Measurable (Function.\allowbreak{}uncurry (fun (w : WD) (x : D) => pivot x+offset w)) := by}
\sourceline{87}{\hspace*{2.0em}have hp := pivot\_\allowbreak{}isometry.\allowbreak{}continuous}
\sourceline{88}{\hspace*{2.0em}unfold Function.\allowbreak{}uncurry offset}
\sourceline{89}{\hspace*{2.0em}fun\_\allowbreak{}prop}
\sourceline{90}{\hspace*{1.0em}have h := (MeasurePreserving.\allowbreak{}id kap).\allowbreak{}skew\_\allowbreak{}product (g:=fun w x => pivot x+offset w) hm}
\sourceline{91}{\hspace*{2.0em}(Filter.\allowbreak{}Eventually.\allowbreak{}of\_\allowbreak{}forall fun w => (hg w).\allowbreak{}map\_\allowbreak{}eq)}
\sourceend
\subsection{Measuring the graph in physical coordinates}
This definition pulls the transverse measure onto its actual domain and pushes it onto the graph. The unit graph weight comes only from this operation. The independently derived fiber factors enter as explicit scalar multipliers.\par\smallskip\id{FourBoxMeasures.lean}\sourceblock
\sourceline{110}{\hspace*{2.0em}exact ((continuous\_\allowbreak{}apply i).\allowbreak{}comp continuous\_\allowbreak{}subtype\_\allowbreak{}val).\allowbreak{}div\_\allowbreak{}const \_\allowbreak{}}
\sourceline{111}{\hspace*{1.0em}have hp : Continuous (fun w : TBox => (t,decode w)) := continuous\_\allowbreak{}const.\allowbreak{}prodMk hd}
\sourceline{112}{\hspace*{1.0em}have hg := Continuous.\allowbreak{}comp (f := fun w : TBox => (t,decode w)) graph\_\allowbreak{}continuous hp}
\sourceline{113}{\hspace*{1.0em}exact continuous\_\allowbreak{}subtype\_\allowbreak{}val.\allowbreak{}prodMk hg.\allowbreak{}fst}
\sourceline{115}{def chartMeasure (factor : \ensuremath{\mathbb R}\ensuremath{\ge}0\ensuremath{\infty}) (kap : Measure W) (t : D) : Measure (W \ensuremath{\times} K) :=}
\sourceline{116}{\hspace*{1.0em}factor \ensuremath{\bullet} (kap.\allowbreak{}comap ((↑) : TBox \ensuremath{\to} W)).\allowbreak{}map (point t)}
\sourceline{118}{theorem graph\_\allowbreak{}support (t : D) (w : TBox) :}
\sourceline{119}{\hspace*{2.0em}gram (point t w).\allowbreak{}2 (point t w).\allowbreak{}1=scalarTarget t := by}
\sourceline{120}{\hspace*{1.0em}have h (j : Fin 5) : 8*(decode w j).\allowbreak{}val = w.\allowbreak{}val j := by change 8*(w.\allowbreak{}val j/8)=\_\allowbreak{}; ring}
\sourceline{121}{\hspace*{1.0em}simpa [point,h] using root\_\allowbreak{}formula t (decode w)}
\sourceline{123}{theorem factor\_\allowbreak{}two (kap : Measure W) (t : D) :}
\sourceline{124}{\hspace*{2.0em}chartMeasure 16 kap t = (2:\ensuremath{\mathbb R}\ensuremath{\ge}0\ensuremath{\infty}) \ensuremath{\bullet} chartMeasure 8 kap t := by}
\sourceline{125}{\hspace*{1.0em}unfold chartMeasure; rw [smul\_\allowbreak{}smul]; norm\_\allowbreak{}num}
\sourceline{127}{theorem measured\_\allowbreak{}box (c : \ensuremath{\mathbb R}\ensuremath{\ge}0\ensuremath{\infty}) (kap : Measure W) \{S : Set W\} (hS : MeasurableSet S)}
\sourceline{128}{\hspace*{2.0em}(hsub : S \ensuremath{\subseteq} transverse 3 0) \{m : \ensuremath{\mathbb N}\} (hm : m \ensuremath{\le} 4) :}
\sourceline{129}{\hspace*{2.0em}chartMeasure c kap 0 (S \ensuremath{\times}\ensuremath{^{s}} cell m 2) = c*kap S := by}
\sourceline{130}{\hspace*{1.0em}have emb : MeasurableEmbedding ((↑) : TBox \ensuremath{\to} W) :=}
\sourceline{131}{\hspace*{2.0em}MeasurableEmbedding.\allowbreak{}subtype\_\allowbreak{}coe (transverse\_\allowbreak{}closed \_\allowbreak{} \_\allowbreak{}).\allowbreak{}measurableSet}
\sourceline{132}{\hspace*{1.0em}have hp : point 0 \ensuremath{^{-}}\ensuremath{^{1}}\textquotesingle{} (S \ensuremath{\times}\ensuremath{^{s}} cell m 2) = ((↑) : TBox \ensuremath{\to} W) \ensuremath{^{-}}\ensuremath{^{1}}\textquotesingle{} S := by}
\sourceline{133}{\hspace*{2.0em}ext w}
\sourceline{134}{\hspace*{2.0em}exact and\_\allowbreak{}iff\_\allowbreak{}left (cell\_\allowbreak{}mono hm 2 (root\_\allowbreak{}child (decode w)))}
\sourceline{135}{\hspace*{1.0em}rw [chartMeasure,Measure.\allowbreak{}smul\_\allowbreak{}apply,Measure.\allowbreak{}map\_\allowbreak{}apply (point\_\allowbreak{}continuous \_\allowbreak{}).\allowbreak{}measurable}
\sourceline{136}{\hspace*{2.0em}(hS.\allowbreak{}prod (cell\_\allowbreak{}closed .\allowbreak{}.\allowbreak{}).\allowbreak{}measurableSet),hp,emb.\allowbreak{}comap\_\allowbreak{}preimage,Subtype.\allowbreak{}range\_\allowbreak{}val\_\allowbreak{}subtype,smul\_\allowbreak{}eq\_\allowbreak{}mul]}
\sourceline{137}{\hspace*{1.0em}change c*kap (S \ensuremath{\cap} transverse 3 0) = \_\allowbreak{}}
\sourceline{138}{\hspace*{1.0em}rw [inter\_\allowbreak{}eq\_\allowbreak{}left.\allowbreak{}mpr hsub]}
\sourceline{140}{end FourChart}
\sourceline{141}{\#print axioms FourChart.\allowbreak{}iso\_\allowbreak{}preserves}
\sourceline{142}{\#print axioms FourChart.\allowbreak{}normalized\_\allowbreak{}preserves}
\sourceline{143}{\#print axioms FourChart.\allowbreak{}graph\_\allowbreak{}support}
\sourceline{144}{\#print axioms FourChart.\allowbreak{}factor\_\allowbreak{}two}
\sourceend
\subsection{Zero and positive children in the same interface}
Disjoint pivot cells imply genuinely empty preimage. The positive endpoints follow by evaluating a product of five one-dimensional masses. The slab law below checks the product with the doubled target Haar on each finer target cell.\par\smallskip\id{FourBoxEndpoints.lean}\sourceblock
\sourceline{20}{theorem negative\_\allowbreak{}child : Disjoint (cell 4 (2:K)) (cell 4 10) := by}
\sourceline{21}{\hspace*{1.0em}apply Set.\allowbreak{}disjoint\_\allowbreak{}left.\allowbreak{}mpr; intro x hx hy}
\sourceline{22}{\hspace*{1.0em}have hx\textquotesingle{} : (2:K) \ensuremath{\in} cell 4 x := by}
\sourceline{23}{\hspace*{2.0em}change \ensuremath{\Vert}(2:K)-x\ensuremath{\Vert} \ensuremath{\le} \_\allowbreak{}}
\sourceline{24}{\hspace*{2.0em}rw [norm\_\allowbreak{}sub\_\allowbreak{}rev]; exact hx}
\sourceline{25}{\hspace*{1.0em}have h := nest (le\_\allowbreak{}refl 4) hy hx\textquotesingle{}}
\sourceline{26}{\hspace*{1.0em}change \ensuremath{\Vert}(2:K)-10\ensuremath{\Vert} \ensuremath{\le} \_\allowbreak{} at h}
\sourceline{27}{\hspace*{1.0em}rw [show (2:K)-10= -(2\textasciicircum{}3) by norm\_\allowbreak{}num,norm\_\allowbreak{}neg,norm\_\allowbreak{}scale] at h}
\sourceline{28}{\hspace*{1.0em}norm\_\allowbreak{}num at h}
\sourceline{30}{theorem missing\_\allowbreak{}mass (c : \ensuremath{\mathbb R}\ensuremath{\ge}0\ensuremath{\infty}) (kap : Measure W) \{S : Set W\} (hs : MeasurableSet S) :}
\sourceline{31}{\hspace*{2.0em}chartMeasure c kap 0 (S \ensuremath{\times}\ensuremath{^{s}} cell 4 10) = 0 := by}
\sourceline{32}{\hspace*{1.0em}have he : point 0 \ensuremath{^{-}}\ensuremath{^{1}}\textquotesingle{} (S \ensuremath{\times}\ensuremath{^{s}} cell 4 10) = \ensuremath{\varnothing} := by}
\sourceline{33}{\hspace*{2.0em}apply Set.\allowbreak{}eq\_\allowbreak{}empty\_\allowbreak{}iff\_\allowbreak{}forall\_\allowbreak{}notMem.\allowbreak{}mpr}
\sourceline{34}{\hspace*{2.0em}intro w hw}
\sourceline{35}{\hspace*{2.0em}exact Set.\allowbreak{}disjoint\_\allowbreak{}left.\allowbreak{}mp negative\_\allowbreak{}child (root\_\allowbreak{}child (decode w)) hw.\allowbreak{}2}
\sourceline{36}{\hspace*{1.0em}rw [chartMeasure,Measure.\allowbreak{}smul\_\allowbreak{}apply,Measure.\allowbreak{}map\_\allowbreak{}apply (point\_\allowbreak{}continuous \_\allowbreak{}).\allowbreak{}measurable}
\sourceline{37}{\hspace*{2.0em}(hs.\allowbreak{}prod (cell\_\allowbreak{}closed .\allowbreak{}.\allowbreak{}).\allowbreak{}measurableSet),he,measure\_\allowbreak{}empty,smul\_\allowbreak{}zero]}
\sourceline{39}{theorem exact\_\allowbreak{}boxes (lam : Measure K) [lam.\allowbreak{}IsAddHaarMeasure] (hvol : lam integers=1) :}
\sourceline{40}{\hspace*{2.0em}let kap := Measure.\allowbreak{}pi fun \_\allowbreak{} : Fin 5 => lam}
\sourceline{41}{\hspace*{2.0em}chartMeasure 16 kap 0 (transverse 3 0 \ensuremath{\times}\ensuremath{^{s}} cell 3 2) = 1/2\textasciicircum{}26 \ensuremath{\wedge}}
\sourceline{42}{\hspace*{2.0em}chartMeasure 8 kap 0 (transverse 3 0 \ensuremath{\times}\ensuremath{^{s}} cell 3 2) = 1/2\textasciicircum{}27 \ensuremath{\wedge}}
\sourceline{43}{\hspace*{2.0em}chartMeasure 16 kap 0 (transverse 4 0 \ensuremath{\times}\ensuremath{^{s}} cell 4 2) = 1/2\textasciicircum{}36 \ensuremath{\wedge}}
\sourceline{44}{\hspace*{2.0em}chartMeasure 8 kap 0 (transverse 4 0 \ensuremath{\times}\ensuremath{^{s}} cell 4 2) = 1/2\textasciicircum{}37 := by}
\sourceline{45}{\hspace*{1.0em}dsimp only}
\sourceline{46}{\hspace*{1.0em}simp only [box\_\allowbreak{}mass lam hvol 3 (by decide) (by decide),box\_\allowbreak{}mass lam hvol 4 (by decide) (by decide)]}
\sourceline{47}{\hspace*{1.0em}repeat\textquotesingle{} constructor}
\sourceline{48}{\hspace*{1.0em}all\_\allowbreak{}goals}
\sourceline{49}{\hspace*{2.0em}apply (ENNReal.\allowbreak{}toReal\_\allowbreak{}eq\_\allowbreak{}toReal\_\allowbreak{}iff\textquotesingle{} (by finiteness) (by finiteness)).\allowbreak{}mp}
\sourceline{50}{\hspace*{2.0em}norm\_\allowbreak{}num [ENNReal.\allowbreak{}toReal\_\allowbreak{}mul,ENNReal.\allowbreak{}toReal\_\allowbreak{}pow,ENNReal.\allowbreak{}toReal\_\allowbreak{}div,ENNReal.\allowbreak{}toReal\_\allowbreak{}ofNat]}
\sourceline{52}{theorem target\_\allowbreak{}slab\_\allowbreak{}law (lam : Measure K) [lam.\allowbreak{}IsAddHaarMeasure] (hvol : lam integers=1)}
\sourceline{53}{\hspace*{2.0em}(kap : Measure W) \{S : Set W\} (hs : MeasurableSet S) (hsub : S \ensuremath{\subseteq} transverse 3 0)}
\sourceline{54}{\hspace*{2.0em}(m r : \ensuremath{\mathbb N}) (hm : 3 \ensuremath{\le} m) (hm\textquotesingle{} : m \ensuremath{\le} 4) (hr : m \ensuremath{\le} r) :}
\sourceline{55}{\hspace*{2.0em}(kap.\allowbreak{}prod lam) (slab S m 2 r 4) =}
\sourceline{56}{\hspace*{2.0em}chartMeasure 8 kap 0 (S \ensuremath{\times}\ensuremath{^{s}} cell m 2) * ((2:\ensuremath{\mathbb R}\ensuremath{\ge}0\ensuremath{\infty}) \ensuremath{\bullet} lam) (cell (r+2) 4) := by}
\sourceline{57}{\hspace*{1.0em}rw [centered\_\allowbreak{}slabs lam hvol kap hs hsub hm hm\textquotesingle{} hr,measured\_\allowbreak{}box \_\allowbreak{} \_\allowbreak{} hs hsub hm\textquotesingle{},}
\sourceline{58}{\hspace*{2.0em}Measure.\allowbreak{}smul\_\allowbreak{}apply,smul\_\allowbreak{}eq\_\allowbreak{}mul,cell\_\allowbreak{}volume lam hvol]}
\sourceline{59}{\hspace*{1.0em}have h : (1/(2\textasciicircum{}r:\ensuremath{\mathbb R}\ensuremath{\ge}0\ensuremath{\infty})\textasciicircum{}2) = 8*(2*(1/(2\textasciicircum{}(r+2):\ensuremath{\mathbb R}\ensuremath{\ge}0\ensuremath{\infty})\textasciicircum{}2)) := by}
\sourceline{60}{\hspace*{2.0em}apply (ENNReal.\allowbreak{}toReal\_\allowbreak{}eq\_\allowbreak{}toReal\_\allowbreak{}iff\textquotesingle{} (by finiteness) (by finiteness)).\allowbreak{}mp}
\sourceline{61}{\hspace*{2.0em}simp only [ENNReal.\allowbreak{}toReal\_\allowbreak{}div,ENNReal.\allowbreak{}toReal\_\allowbreak{}mul,ENNReal.\allowbreak{}toReal\_\allowbreak{}one,ENNReal.\allowbreak{}toReal\_\allowbreak{}ofNat,}
\sourceline{62}{\hspace*{3.0em}ENNReal.\allowbreak{}toReal\_\allowbreak{}pow,pow\_\allowbreak{}add]}
\sourceline{63}{\hspace*{2.0em}ring}
\sourceline{64}{\hspace*{1.0em}rw [h]; ac\_\allowbreak{}rfl}
\sourceend
\subsection{The readable expression and its field hypotheses}
Definitional equality compares with the actual upstream constant. Erasing the ambient-place data for the subsequent algebraic calculation makes the minimal hypotheses visible. The cancellation uses invertibility of the pivot and the trace of the chosen field element.\par\smallskip\id{ReadableFourChart.lean}\sourceblock
\sourceline{17}{def algebraChart (\ensuremath{\tau} s : A) (T : Matrix I I A)}
\sourceline{18}{\hspace*{2.0em}(p : (A \ensuremath{\times} (I \ensuremath{\to} A)) \ensuremath{\times} skewAdjoint A) : Fin 2 \ensuremath{\oplus} I \ensuremath{\to} A :=}
\sourceline{19}{\hspace*{1.0em}let ((t,w),a) := p}
\sourceline{20}{\hspace*{1.0em}![(star t)\ensuremath{^{-}}\ensuremath{^{1}}\ensuremath{^{r}} * \ensuremath{\tau} * (s - star w \ensuremath{\cdot}\ensuremath{_{v}} T *\ensuremath{_{v}} w) + t * a, t] \ensuremath{\oplus}\ensuremath{_{v}} w}
\sourceline{22}{theorem literal [TopologicalSpace A] [MeasureSpace A] [P : Place A] (T : Matrix I I A)}
\sourceline{23}{\hspace*{2.0em}(b : skewAdjoint A) :}
\sourceline{24}{\hspace*{2.0em}Place.\allowbreak{}flatChart T b = algebraChart P.\allowbreak{}element (P.\allowbreak{}skewUnit\ensuremath{^{-}}\ensuremath{^{1}}.\allowbreak{}val * b.\allowbreak{}val) T := rfl}
\sourceline{27}{theorem form (\ensuremath{\tau} s : A) (h\ensuremath{\tau} : \ensuremath{\tau} + star \ensuremath{\tau} = 1) (hs : star s=s) (T : Matrix I I A)}
\sourceline{28}{\hspace*{2.0em}(hT : T\ensuremath{^{H}}=T) (p : (A \ensuremath{\times} (I \ensuremath{\to} A)) \ensuremath{\times} skewAdjoint A) (hp : IsUnit p.\allowbreak{}1.\allowbreak{}1) :}
\sourceline{29}{\hspace*{2.0em}star (algebraChart \ensuremath{\tau} s T p) \ensuremath{\cdot}\ensuremath{_{v}} planeSumGram T *\ensuremath{_{v}} algebraChart \ensuremath{\tau} s T p = s := by}
\sourceline{30}{\hspace*{1.0em}obtain \ensuremath{\langle}\ensuremath{\langle}t,w\ensuremath{\rangle},a\ensuremath{\rangle} := p}
\sourceline{31}{\hspace*{1.0em}let q : A := star w \ensuremath{\cdot}\ensuremath{_{v}} T *\ensuremath{_{v}} w}
\sourceline{32}{\hspace*{1.0em}suffices ((star t)\ensuremath{^{-}}\ensuremath{^{1}}\ensuremath{^{r}} * \ensuremath{\tau} * (s-q)+t*a)*star t +}
\sourceline{33}{\hspace*{3.0em}t*(t\ensuremath{^{-}}\ensuremath{^{1}}\ensuremath{^{r}}*star \ensuremath{\tau}*(s-q)-star t*a)+q=s by}
\sourceline{34}{\hspace*{2.0em}simpa [algebraChart,fromBlocks\_\allowbreak{}mulVec,star\_\allowbreak{}sumElim,inverse\_\allowbreak{}star,}
\sourceline{35}{\hspace*{3.0em}star\_\allowbreak{}dotProduct\_\allowbreak{}mulVec\_\allowbreak{}comm hT w w,hs,\ensuremath{\leftarrow}sub\_\allowbreak{}eq\_\allowbreak{}add\_\allowbreak{}neg]}
\sourceline{36}{\hspace*{1.0em}linear\_\allowbreak{}combination (s-q)*(\ensuremath{\tau}*inverse\_\allowbreak{}mul\_\allowbreak{}cancel (star t) hp.\allowbreak{}star +}
\sourceline{37}{\hspace*{2.0em}star \ensuremath{\tau}*inverse\_\allowbreak{}mul\_\allowbreak{}cancel t hp+h\ensuremath{\tau})}
\sourceend
\subsection{A determinant and a pairing matrix}
The coordinate matrix contains a half, which makes integral-lattice preservation an unjustified assumption. Its Gram matrix and determinant have direct algebraic proofs. Evaluating the matrix on the central parameter identifies the center without appealing to transitivity of an abstract isometry group.\par\smallskip\id{GaussianFrame.lean}\sourceblock
\sourceline{10}{def frame (c : K) : Matrix (Fin 3) (Fin 3) (G K) :=}
\sourceline{11}{\hspace*{1.0em}fun}
\sourceline{12}{\hspace*{2.0em}| 0,0 => \ensuremath{\langle}1,0\ensuremath{\rangle}}
\sourceline{13}{\hspace*{2.0em}| 0,1 => \ensuremath{\langle}1/2,0\ensuremath{\rangle}}
\sourceline{14}{\hspace*{2.0em}| 0,2 => \ensuremath{\langle}0,0\ensuremath{\rangle}}
\sourceline{15}{\hspace*{2.0em}| 1,0 => \ensuremath{\langle}1,1\ensuremath{\rangle}}
\sourceline{16}{\hspace*{2.0em}| 1,1 => \ensuremath{\langle}-1/2,-1/2\ensuremath{\rangle}}
\sourceline{17}{\hspace*{2.0em}| 1,2 => \ensuremath{\langle}c,-2\ensuremath{\rangle}}
\sourceline{18}{\hspace*{2.0em}| 2,0 => \ensuremath{\langle}c,2\ensuremath{\rangle}}
\sourceline{19}{\hspace*{2.0em}| 2,1 => \ensuremath{\langle}-c/2,-1\ensuremath{\rangle}}
\sourceline{20}{\hspace*{2.0em}| 2,2 => \ensuremath{\langle}-1,1\ensuremath{\rangle}}
\sourceline{22}{def plane (K : Type) [Field K] : Matrix (Fin 3) (Fin 3) (G K) :=}
\sourceline{23}{\hspace*{1.0em}!![0,1,0;1,0,0;0,0,-1]}
\sourceline{25}{local macro "frame\_\allowbreak{}entry" : tactic => `(tactic|}
\sourceline{26}{\hspace*{1.0em}(apply QuadraticAlgebra.\allowbreak{}ext <;>}
\sourceline{27}{\hspace*{1.5em}norm\_\allowbreak{}num [frame,plane,Matrix.\allowbreak{}mul\_\allowbreak{}apply,Matrix.\allowbreak{}conjTranspose\_\allowbreak{}apply,Fin.\allowbreak{}sum\_\allowbreak{}univ\_\allowbreak{}succ,}
\sourceline{28}{\hspace*{2.5em}QuadraticAlgebra.\allowbreak{}re\_\allowbreak{}star,QuadraticAlgebra.\allowbreak{}im\_\allowbreak{}star,QuadraticAlgebra.\allowbreak{}re\_\allowbreak{}mul,}
\sourceline{29}{\hspace*{2.5em}QuadraticAlgebra.\allowbreak{}im\_\allowbreak{}mul,QuadraticAlgebra.\allowbreak{}re\_\allowbreak{}one,QuadraticAlgebra.\allowbreak{}im\_\allowbreak{}one,}
\sourceline{30}{\hspace*{2.5em}QuadraticAlgebra.\allowbreak{}re\_\allowbreak{}zero,QuadraticAlgebra.\allowbreak{}im\_\allowbreak{}zero] <;> try ring\_\allowbreak{}nf))}
\sourceline{32}{theorem gram\_\allowbreak{}frame (c : K) (hc : c\textasciicircum{}2 = -7) : (frame c)\ensuremath{^{H}} * frame c = plane K := by}
\sourceline{33}{\hspace*{1.0em}ext i j : 1}
\sourceline{34}{\hspace*{1.0em}fin\_\allowbreak{}cases i <;> fin\_\allowbreak{}cases j <;> frame\_\allowbreak{}entry <;>}
\sourceline{35}{\hspace*{2.0em}try simp only [hc,QuadraticAlgebra.\allowbreak{}re\_\allowbreak{}ofNat,QuadraticAlgebra.\allowbreak{}im\_\allowbreak{}ofNat] <;> ring}
\sourceline{37}{theorem det\_\allowbreak{}frame (c : K) (hc : c\textasciicircum{}2 = -7) : (frame c).\allowbreak{}det = -1 := by}
\sourceline{38}{\hspace*{1.0em}rw [Matrix.\allowbreak{}det\_\allowbreak{}fin\_\allowbreak{}three]}
\sourceline{39}{\hspace*{1.0em}frame\_\allowbreak{}entry <;> try simp only [hc] <;> ring}
\sourceline{41}{theorem central\_\allowbreak{}vector (c : K) :}
\sourceline{42}{\hspace*{2.0em}frame c *\ensuremath{_{v}} (![1,2,0] : Fin 3 \ensuremath{\to} G K) = ![2,0,0] := by}
\sourceline{43}{\hspace*{1.0em}ext i : 1}
\sourceline{44}{\hspace*{1.0em}fin\_\allowbreak{}cases i <;> simp [Matrix.\allowbreak{}mulVec,dotProduct,Fin.\allowbreak{}sum\_\allowbreak{}univ\_\allowbreak{}succ] <;> frame\_\allowbreak{}entry <;>}
\sourceline{45}{\hspace*{2.0em}simp only [QuadraticAlgebra.\allowbreak{}re\_\allowbreak{}ofNat,QuadraticAlgebra.\allowbreak{}im\_\allowbreak{}ofNat] <;> ring}
\sourceline{47}{theorem pullback \{I J : Type\} [Fintype I] [Fintype J]}
\sourceline{48}{\hspace*{2.0em}(M : Matrix I J (G K)) (v : J \ensuremath{\to} G K) :}
\sourceline{49}{\hspace*{2.0em}star (M *\ensuremath{_{v}} v) \ensuremath{\cdot}\ensuremath{_{v}} (M *\ensuremath{_{v}} v) = star v \ensuremath{\cdot}\ensuremath{_{v}} ((M\ensuremath{^{H}} * M) *\ensuremath{_{v}} v) := by}
\sourceline{50}{\hspace*{1.0em}rw [\ensuremath{\leftarrow}Matrix.\allowbreak{}mulVec\_\allowbreak{}mulVec,Matrix.\allowbreak{}dotProduct\_\allowbreak{}mulVec (star v) M\ensuremath{^{H}},Matrix.\allowbreak{}star\_\allowbreak{}mulVec]}
\sourceend
\subsection{Keeping the readable pivot invertible}
In the chosen frame the pivot of a vector in the parent has real part of norm one half and imaginary part of norm at most one eighth. Their squares have unequal norms. The denominator of the inverse consequently survives throughout the parent box.\par\smallskip\id{FourFramePivot.lean}\sourceblock
\sourceline{10}{def framePivot (c : D) (a b d : GInt) : GaussianFrame.\allowbreak{}G K :=}
\sourceline{11}{\hspace*{1.0em}2+8*includeGaussian a + star (\ensuremath{\langle}1,1\ensuremath{\rangle} : GaussianFrame.\allowbreak{}G K)*(8*includeGaussian b) +}
\sourceline{12}{\hspace*{2.0em}star (\ensuremath{\langle}c.\allowbreak{}val,2\ensuremath{\rangle} : GaussianFrame.\allowbreak{}G K)*(8*includeGaussian d)}
\sourceline{14}{theorem small\_\allowbreak{}pivot\_\allowbreak{}unit (u v : D) :}
\sourceline{15}{\hspace*{2.0em}IsUnit (\ensuremath{\langle}2+8*u.\allowbreak{}val,8*v.\allowbreak{}val\ensuremath{\rangle} : GaussianFrame.\allowbreak{}G K) := by}
\sourceline{16}{\hspace*{1.0em}have h8 (z : D) : \ensuremath{\Vert}(8:K)*z.\allowbreak{}val\ensuremath{\Vert} \ensuremath{\le} 1/8 := by}
\sourceline{17}{\hspace*{2.0em}rw [show (8:K)=2\textasciicircum{}3 by norm\_\allowbreak{}num,norm\_\allowbreak{}mul,norm\_\allowbreak{}scale]}
\sourceline{18}{\hspace*{2.0em}norm\_\allowbreak{}num}
\sourceline{19}{\hspace*{2.0em}linarith [(mem\_\allowbreak{}D \_\allowbreak{}).\allowbreak{}mp z.\allowbreak{}2]}
\sourceline{20}{\hspace*{1.0em}have hn : \ensuremath{\Vert}(2:K)+8*u.\allowbreak{}val\ensuremath{\Vert} = 1/2 := by}
\sourceline{21}{\hspace*{2.0em}rw [IsUltrametricDist.\allowbreak{}norm\_\allowbreak{}add\_\allowbreak{}eq\_\allowbreak{}max\_\allowbreak{}of\_\allowbreak{}norm\_\allowbreak{}ne\_\allowbreak{}norm (by rw [norm\_\allowbreak{}two]; linarith [h8 u]),}
\sourceline{22}{\hspace*{3.0em}norm\_\allowbreak{}two,max\_\allowbreak{}eq\_\allowbreak{}left (by linarith [h8 u])]}
\sourceline{23}{\hspace*{1.0em}apply QuadraticAlgebra.\allowbreak{}isUnit\_\allowbreak{}iff\_\allowbreak{}norm\_\allowbreak{}isUnit.\allowbreak{}mpr}
\sourceline{24}{\hspace*{1.0em}rw [isUnit\_\allowbreak{}iff\_\allowbreak{}ne\_\allowbreak{}zero]}
\sourceline{25}{\hspace*{1.0em}change (2+8*u.\allowbreak{}val)*(2+8*u.\allowbreak{}val)+0*(2+8*u.\allowbreak{}val)*(8*v.\allowbreak{}val)-(-1)*(8*v.\allowbreak{}val)*(8*v.\allowbreak{}val) \ensuremath{\ne} 0}
\sourceline{26}{\hspace*{1.0em}have h : (2+8*u.\allowbreak{}val)\textasciicircum{}2 + (8*v.\allowbreak{}val)\textasciicircum{}2 \ensuremath{\ne} 0 := by}
\sourceline{27}{\hspace*{2.0em}intro h}
\sourceline{28}{\hspace*{2.0em}have he : (2+8*u.\allowbreak{}val)\textasciicircum{}2= -(8*v.\allowbreak{}val)\textasciicircum{}2 := eq\_\allowbreak{}neg\_\allowbreak{}of\_\allowbreak{}add\_\allowbreak{}eq\_\allowbreak{}zero\_\allowbreak{}left h}
\sourceline{29}{\hspace*{2.0em}have hh := congrArg (fun x : K => \ensuremath{\Vert}x\ensuremath{\Vert}) he}
\sourceline{30}{\hspace*{2.0em}simp only [norm\_\allowbreak{}neg,norm\_\allowbreak{}pow,hn] at hh}
\sourceline{31}{\hspace*{2.0em}nlinarith [h8 v,norm\_\allowbreak{}nonneg ((8:K)*v.\allowbreak{}val),sq\_\allowbreak{}nonneg (1/8-\ensuremath{\Vert}(8:K)*v.\allowbreak{}val\ensuremath{\Vert}),}
\sourceline{32}{\hspace*{3.0em}mul\_\allowbreak{}self\_\allowbreak{}le\_\allowbreak{}mul\_\allowbreak{}self (norm\_\allowbreak{}nonneg ((8:K)*v.\allowbreak{}val)) (h8 v)]}
\sourceline{33}{\hspace*{1.0em}convert h using 1 <;> ring}
\sourceline{35}{theorem original\_\allowbreak{}patch\_\allowbreak{}unit (c : D) (a b d : GInt) : IsUnit (framePivot c a b d) := by}
\sourceline{36}{\hspace*{1.0em}let u := a.\allowbreak{}re+b.\allowbreak{}re+b.\allowbreak{}im+c*d.\allowbreak{}re+2*d.\allowbreak{}im}
\sourceline{37}{\hspace*{1.0em}let v := a.\allowbreak{}im+b.\allowbreak{}im-b.\allowbreak{}re+c*d.\allowbreak{}im-2*d.\allowbreak{}re}
\sourceline{38}{\hspace*{1.0em}have h : framePivot c a b d = (\ensuremath{\langle}2+8*u.\allowbreak{}val,8*v.\allowbreak{}val\ensuremath{\rangle} : GaussianFrame.\allowbreak{}G K) := by}
\sourceline{39}{\hspace*{2.0em}ext <;> simp [framePivot,includeGaussian,u,v,QuadraticAlgebra.\allowbreak{}re\_\allowbreak{}star,QuadraticAlgebra.\allowbreak{}im\_\allowbreak{}star,}
\sourceline{40}{\hspace*{3.0em}QuadraticAlgebra.\allowbreak{}re\_\allowbreak{}ofNat,QuadraticAlgebra.\allowbreak{}im\_\allowbreak{}ofNat,QuadraticAlgebra.\allowbreak{}re\_\allowbreak{}mul,}
\sourceline{41}{\hspace*{3.0em}QuadraticAlgebra.\allowbreak{}im\_\allowbreak{}mul,show ((2:D):K)=2 by rfl] <;> ring}
\sourceline{42}{\hspace*{1.0em}rw [h]; exact small\_\allowbreak{}pivot\_\allowbreak{}unit u v}
\sourceend
\subsection{Putting the readable output back into the identity space}
The readable source indexes its plane and residual line by a disjoint sum. The index conversion in this proof must preserve the pairing matrix. The final part supplies the local square root needed for the frame directly in this edition.\par\smallskip\id{ReadableIdentityChart.lean}\sourceblock
\sourceline{9}{def idx : Fin 2 \ensuremath{\oplus} Fin 1 \ensuremath{\to} Fin 3 := Sum.\allowbreak{}elim Fin.\allowbreak{}castSucc (fun \_\allowbreak{} => 2)}
\sourceline{11}{variable \{K : Type\} [Field K] [CharZero K]}
\sourceline{13}{def M (c : K) := (GaussianFrame.\allowbreak{}frame c).\allowbreak{}submatrix id idx}
\sourceline{14}{def residual (K : Type) [Field K] : Matrix (Fin 1) (Fin 1) (G K) := fun \_\allowbreak{} \_\allowbreak{} => -1}
\sourceline{16}{theorem frame\_\allowbreak{}pullback (c : K) (hc : c\textasciicircum{}2 = -7) :}
\sourceline{17}{\hspace*{2.0em}(M c)\ensuremath{^{H}} * M c = planeSumGram (residual K) := by}
\sourceline{18}{\hspace*{1.0em}ext a b : 1}
\sourceline{19}{\hspace*{1.0em}change ((GaussianFrame.\allowbreak{}frame c)\ensuremath{^{H}} * GaussianFrame.\allowbreak{}frame c) (idx a) (idx b) = \_\allowbreak{}}
\sourceline{20}{\hspace*{1.0em}rw [GaussianFrame.\allowbreak{}gram\_\allowbreak{}frame c hc]}
\sourceline{21}{\hspace*{1.0em}rcases a with a|a <;> rcases b with b|b <;> fin\_\allowbreak{}cases a <;> fin\_\allowbreak{}cases b <;> rfl}
\sourceline{23}{theorem gram\_\allowbreak{}original (c z : K) (hc : c\textasciicircum{}2 = -7)}
\sourceline{24}{\hspace*{2.0em}(p : (G K \ensuremath{\times} (Fin 1 \ensuremath{\to} G K)) \ensuremath{\times} skewAdjoint (G K)) (hp : IsUnit p.\allowbreak{}1.\allowbreak{}1) :}
\sourceline{25}{\hspace*{2.0em}let s := -(imaginary K 1).\allowbreak{}val*(imaginary K (4+32*z)).\allowbreak{}val}
\sourceline{26}{\hspace*{2.0em}let v := algebraChart (halfTrace K) s (residual K) p}
\sourceline{27}{\hspace*{2.0em}star ((M c) *\ensuremath{_{v}} v) \ensuremath{\cdot}\ensuremath{_{v}} ((M c) *\ensuremath{_{v}} v) = hermitian K (4+32*z) := by}
\sourceline{28}{\hspace*{1.0em}dsimp only}
\sourceline{29}{\hspace*{1.0em}rw [GaussianFrame.\allowbreak{}pullback,frame\_\allowbreak{}pullback c hc]}
\sourceline{30}{\hspace*{1.0em}exact four\_\allowbreak{}chart K z (residual K) (by ext i j : 1; simp [residual,conjTranspose\_\allowbreak{}apply]) p hp}
\sourceline{32}{abbrev LocalK := QuadraticAlgebra \ensuremath{\mathbb Q}\_\allowbreak{}[2] 2 0}
\sourceline{33}{instance : Fact (\ensuremath{\forall} x : \ensuremath{\mathbb Q}\_\allowbreak{}[2], x\textasciicircum{}2 \ensuremath{\ne} 2+0*x) := by}
\sourceline{34}{\hspace*{1.0em}constructor; intro x hx}
\sourceline{35}{\hspace*{1.0em}simp only [zero\_\allowbreak{}mul,add\_\allowbreak{}zero] at hx}
\sourceline{36}{\hspace*{1.0em}have hv := congrArg Padic.\allowbreak{}valuation hx}
\sourceline{37}{\hspace*{1.0em}rw [Padic.\allowbreak{}valuation\_\allowbreak{}pow,show (2:\ensuremath{\mathbb Q}\_\allowbreak{}[2]).\allowbreak{}valuation = 1 from Padic.\allowbreak{}valuation\_\allowbreak{}p] at hv}
\sourceline{38}{\hspace*{1.0em}omega}
\sourceline{40}{theorem root\_\allowbreak{}exists : \ensuremath{\exists} c : LocalK, c\textasciicircum{}2 = -7 := by}
\sourceline{41}{\hspace*{1.0em}let F : Polynomial \ensuremath{\mathbb Z}\_\allowbreak{}[2] := Polynomial.\allowbreak{}X\textasciicircum{}2+Polynomial.\allowbreak{}C 7}
\sourceline{42}{\hspace*{1.0em}have hnorm : \ensuremath{\Vert}F.\allowbreak{}aeval (1:\ensuremath{\mathbb Z}\_\allowbreak{}[2])\ensuremath{\Vert} < \ensuremath{\Vert}F.\allowbreak{}derivative.\allowbreak{}aeval (1:\ensuremath{\mathbb Z}\_\allowbreak{}[2])\ensuremath{\Vert}\textasciicircum{}2 := by}
\sourceline{43}{\hspace*{2.0em}have hn : \ensuremath{\Vert}(2:\ensuremath{\mathbb Z}\_\allowbreak{}[2])\ensuremath{\Vert} = 1/2 := by simpa using (PadicInt.\allowbreak{}norm\_\allowbreak{}p (p:=2))}
\sourceline{44}{\hspace*{2.0em}have h : \ensuremath{\Vert}(8:\ensuremath{\mathbb Z}\_\allowbreak{}[2])\ensuremath{\Vert} < \ensuremath{\Vert}(2:\ensuremath{\mathbb Z}\_\allowbreak{}[2])\ensuremath{\Vert}\textasciicircum{}2 := by}
\sourceline{45}{\hspace*{3.0em}rw [show (8:\ensuremath{\mathbb Z}\_\allowbreak{}[2])=2\textasciicircum{}3 by norm\_\allowbreak{}num,norm\_\allowbreak{}pow,hn]; norm\_\allowbreak{}num}
\sourceline{46}{\hspace*{2.0em}convert h using 2 <;> norm\_\allowbreak{}num [F,Polynomial.\allowbreak{}derivative\_\allowbreak{}pow]}
\sourceline{47}{\hspace*{1.0em}obtain \ensuremath{\langle}a,ha,-\ensuremath{\rangle} := hensels\_\allowbreak{}lemma hnorm}
\sourceline{48}{\hspace*{1.0em}have hh : a\textasciicircum{}2+7=0 := by simpa [F] using ha}
\sourceline{49}{\hspace*{1.0em}let f : \ensuremath{\mathbb Z}\_\allowbreak{}[2] \ensuremath{\to}+* LocalK := (algebraMap \ensuremath{\mathbb Q}\_\allowbreak{}[2] LocalK).\allowbreak{}comp PadicInt.\allowbreak{}Coe.\allowbreak{}ringHom}
\sourceline{50}{\hspace*{1.0em}refine \ensuremath{\langle}f a,?\_\allowbreak{}\ensuremath{\rangle}}
\sourceline{51}{\hspace*{1.0em}have hz := congrArg f hh}
\sourceline{52}{\hspace*{1.0em}simp only [map\_\allowbreak{}add,map\_\allowbreak{}pow,map\_\allowbreak{}ofNat,map\_\allowbreak{}zero] at hz}
\sourceline{53}{\hspace*{1.0em}exact eq\_\allowbreak{}neg\_\allowbreak{}of\_\allowbreak{}add\_\allowbreak{}eq\_\allowbreak{}zero\_\allowbreak{}left hz}
\sourceend
\subsection{Why multiplication by the pivot unit preserves the box}
Both directions matter. The inverse rotation uses the inverse norm as a unit of the integral subring. Treating it as merely a nonzero field element would leave a denominator unaccounted for in the width of the skew parameter.\par\smallskip\id{ReadableBoxParameters.lean}\sourceblock
\sourceline{19}{theorem rotate (S : Subring K) (u v : S) (N : S\ensuremath{^{\times}})}
\sourceline{20}{\hspace*{2.0em}(hN : (u:K)\textasciicircum{}2+(v:K)\textasciicircum{}2=(N.\allowbreak{}val:K)) (w x y : K) :}
\sourceline{21}{\hspace*{2.0em}box S w ((u:K)*x-(v:K)*y) ((v:K)*x+(u:K)*y) \ensuremath{\leftrightarrow} box S w x y := by}
\sourceline{22}{\hspace*{1.0em}have hinv : ((N\ensuremath{^{-}}\ensuremath{^{1}}).\allowbreak{}val:K)*(N.\allowbreak{}val:K)=1 := by exact\_\allowbreak{}mod\_\allowbreak{}cast N.\allowbreak{}inv\_\allowbreak{}mul}
\sourceline{23}{\hspace*{1.0em}constructor}
\sourceline{24}{\hspace*{1.0em}\ensuremath{\cdot} rintro \ensuremath{\langle}f,g,hf,hg\ensuremath{\rangle}}
\sourceline{25}{\hspace*{2.0em}refine \ensuremath{\langle}(N\ensuremath{^{-}}\ensuremath{^{1}}).\allowbreak{}val*(u*f+v*g),(N\ensuremath{^{-}}\ensuremath{^{1}}).\allowbreak{}val*(-v*f+u*g),?\_\allowbreak{},?\_\allowbreak{}\ensuremath{\rangle}}
\sourceline{26}{\hspace*{2.0em}all\_\allowbreak{}goals push\_\allowbreak{}cast}
\sourceline{27}{\hspace*{2.0em}\ensuremath{\cdot} linear\_\allowbreak{}combination ((N\ensuremath{^{-}}\ensuremath{^{1}}).\allowbreak{}val:K)*(u:K)*hf + ((N\ensuremath{^{-}}\ensuremath{^{1}}).\allowbreak{}val:K)*(v:K)*hg -}
\sourceline{28}{\hspace*{4.0em}((N\ensuremath{^{-}}\ensuremath{^{1}}).\allowbreak{}val:K)*x*hN - x*hinv}
\sourceline{29}{\hspace*{2.0em}\ensuremath{\cdot} linear\_\allowbreak{}combination -((N\ensuremath{^{-}}\ensuremath{^{1}}).\allowbreak{}val:K)*(v:K)*hf + ((N\ensuremath{^{-}}\ensuremath{^{1}}).\allowbreak{}val:K)*(u:K)*hg -}
\sourceline{30}{\hspace*{4.0em}((N\ensuremath{^{-}}\ensuremath{^{1}}).\allowbreak{}val:K)*y*hN - y*hinv}
\sourceline{31}{\hspace*{1.0em}\ensuremath{\cdot} rintro \ensuremath{\langle}f,g,rfl,rfl\ensuremath{\rangle}}
\sourceline{32}{\hspace*{2.0em}refine \ensuremath{\langle}u*f-v*g,v*f+u*g,?\_\allowbreak{},?\_\allowbreak{}\ensuremath{\rangle} <;> push\_\allowbreak{}cast <;> ring}
\sourceline{36}{theorem offset\_\allowbreak{}box (S : Subring K) (u v : S) (N : S\ensuremath{^{\times}})}
\sourceline{37}{\hspace*{2.0em}(hN : (u:K)\textasciicircum{}2+(v:K)\textasciicircum{}2=(N.\allowbreak{}val:K)) (h a : K) (k : S) :}
\sourceline{38}{\hspace*{2.0em}box S (2*h) (2*((u:K)*(h*k)-(v:K)*a)) (2*((v:K)*(h*k)+(u:K)*a)) \ensuremath{\leftrightarrow}}
\sourceline{39}{\hspace*{2.0em}\ensuremath{\exists} t : S, a=h*(t:K) := by}
\sourceline{40}{\hspace*{1.0em}have scale (x y : K) : box S (2*h) (2*x) (2*y) \ensuremath{\leftrightarrow} box S h x y := by}
\sourceline{41}{\hspace*{2.0em}constructor <;> rintro \ensuremath{\langle}f,g,hf,hg\ensuremath{\rangle} <;> refine \ensuremath{\langle}f,g,?\_\allowbreak{},?\_\allowbreak{}\ensuremath{\rangle}}
\sourceline{42}{\hspace*{2.0em}\ensuremath{\cdot} linear\_\allowbreak{}combination hf/2}
\sourceline{43}{\hspace*{2.0em}\ensuremath{\cdot} linear\_\allowbreak{}combination hg/2}
\sourceline{44}{\hspace*{2.0em}\ensuremath{\cdot} linear\_\allowbreak{}combination 2*hf}
\sourceline{45}{\hspace*{2.0em}\ensuremath{\cdot} linear\_\allowbreak{}combination 2*hg}
\sourceline{46}{\hspace*{1.0em}rw [scale,rotate S u v N hN,axis S h a k]}
\sourceend
\subsection{The real correction and the lower minor}
The first identity isolates the entire dependence on the target and on the two Gaussian parameters in one integral correction. The next identity inverts the lower two physical coordinates, which is what permits an actual box preimage computation.\par\smallskip\id{ReadableBoxParameters.lean}\sourceblock
\sourceline{51}{\hspace*{2.0em}(hN : N = (1+2*d*A)\textasciicircum{}2+(2*d*B)\textasciicircum{}2) (hne : N \ensuremath{\ne} 0) :}
\sourceline{52}{\hspace*{2.0em}((4+16*d*t)+((4*d*L)\textasciicircum{}2+(4*d*M)\textasciicircum{}2)-}
\sourceline{53}{\hspace*{3.0em}((2+4*d*A)\textasciicircum{}2+(4*d*B)\textasciicircum{}2)) / (2*((2+4*d*A)\textasciicircum{}2+(4*d*B)\textasciicircum{}2)) =}
\sourceline{54}{\hspace*{3.0em}2*d*((t-A+d*(L\textasciicircum{}2+M\textasciicircum{}2-A\textasciicircum{}2-B\textasciicircum{}2))/N) := by}
\sourceline{55}{\hspace*{1.0em}have he : (2+4*d*A)\textasciicircum{}2+(4*d*B)\textasciicircum{}2 = 4*N := by rw [hN]; ring}
\sourceline{56}{\hspace*{1.0em}rw [he]; field\_\allowbreak{}simp; rw [hN]; ring}
\sourceline{60}{theorem bottom\_\allowbreak{}inverse \{R : Type*\} [CommRing R] (B D E F b v : R)}
\sourceline{61}{\hspace*{2.0em}(h : B*E+D*F = -1) :}
\sourceline{62}{\hspace*{2.0em}(-E)*(B*b+F*v)-F*(D*b-E*v)=b \ensuremath{\wedge}}
\sourceline{63}{\hspace*{2.0em}(-D)*(B*b+F*v)+B*(D*b-E*v)=v := by}
\sourceline{64}{\hspace*{1.0em}constructor <;> first | linear\_\allowbreak{}combination -b*h | linear\_\allowbreak{}combination -v*h}
\sourceline{66}{theorem width\_\allowbreak{}ratio (m : \ensuremath{\mathbb N}) (hm : 1 \ensuremath{\le} m) :}
\sourceline{67}{\hspace*{2.0em}2*(1/4:\ensuremath{\mathbb Q})\textasciicircum{}(2*m)*(1/4)\textasciicircum{}(2*m)*(1/4)\textasciicircum{}(m-1) = 8*(1/4)\textasciicircum{}(5*m) := by}
\sourceline{68}{\hspace*{1.0em}rw [show 5*m = (2*m+2*m+(m-1))+1 by omega, pow\_\allowbreak{}succ,}
\sourceline{69}{\hspace*{2.0em}pow\_\allowbreak{}add (1/4:\ensuremath{\mathbb Q}) (2*m+2*m) (m-1), pow\_\allowbreak{}add (1/4:\ensuremath{\mathbb Q}) (2*m) (2*m)]}
\sourceline{70}{\hspace*{1.0em}ring}
\sourceline{72}{end ReadableBoxParameters}
\sourceend

\section{Source register and reproduction}
The primary sources identify the exact editions used in the comparison. The earlier regular-family witness, the readable chart definition and the new finite data play different evidential roles.

\begin{enumerate}[leftmargin=1.5em]
\item Original regular fibers. The published revision specifies the normalization of coordinate Haar and the regular-family contract. The previous study realizes the relevant completion and proves its square-dilate mass formula. See the \href{https://github.com/nousentllc/riemann-hypothesis/blob/40e56bab026552e089d7afcfeff08d2defb9c005/src/Formalism/SharpRestriction/GlobalSchurBounds/GeneralLinearActions/GenusMassVolumes.lean}{original family definition} and the \href{https://github.com/nousentllc/riemann-hypothesis/tree/17263859a4d27c62cfc67571737355f61a080806/research/uriel/experiments/ramified-square-dilates}{ramified square-dilate supplement}.
\item Target Haar. The earlier paper proves the split comparison for a fixed vector source. That calculation fixes the denominator before the present vector boxes are measured. See the \href{https://github.com/nousentllc/riemann-hypothesis/tree/4be022986be0e29b4d38f9d04b17956131fe4d56/research/uriel/experiments/target-four-haar}{target-four Haar release}.
\item Readable chart. The chart formula begins at line 239, its form computation at line 386 and the global disintegration theorem at line 616. The source replay in this release recompiles this module, with 27 project modules supplied as cached imports. See \href{https://github.com/nousentllc/riemann-hypothesis/blob/e063a923620ee7787b4a1270421478645bf66bdb/src/Lib/General/Measure/QuadricDisintegration.lean#L239}{the pinned readable chart} and \href{https://github.com/nousentllc/riemann-hypothesis/blob/e063a923620ee7787b4a1270421478645bf66bdb/src/Lib/SiegelWeil/OneVector/FinitePlaceDensity.lean}{the distinct integral-place specialization}.
\end{enumerate}

For the finite control, execute \id{python3 -B replay.py} from the extracted archive. It recreates every record from the six-variable polynomial and compares the output with the archived data. For the two formal environments the README gives separate commands. The additional parameter replay checks the scalar box identities, the lower minor and the rational volume formula. Each rebuilds the indicated source modules in a fresh directory, then runs Lean's standalone checker relative to the imported environment. The three JSON receipts specify the queried declarations and list their foundational dependencies. A source digest describes the input bytes. Kernel replay checks the proof terms arising from those bytes and their imports. Neither operation certifies the unexamined remainder of an upstream development.

The package uses one editable LuaLaTeX source; the original Nous PNG is embedded in its payload. The checksum manifest covers the releases of the paper, data, proof sources and replay scripts. The archive is meant to permit independent mathematical criticism of both the input normalization and the claimed scope of the result.
\end{document}

BEGIN FIGURE nouslogo
89504e470d0a1a0a0000000d4948445200000296000004960806000001e725dca6000000097048597300002e2300002e230178a53f76
0000200049444154789cecbd3f8c16c913f737fdd3459665c0720e64ce00a796059b387002a44e0ed23739909c1fc889250707927320
b0e4c0d241e0c801acf4869660035b7a238ed8b2616559b21cb8ad5e6a96d9e93fd37faabbab67be1f6975dcccf3cc9f7e7abe53dd55
5da5b4d6532a4aa9575aeb67c95f148c52eaadd6fa49c9156635e6e59795325f3ed35adfb5760e8069c0699a7ed75a2b8eab2d6acccb
83fc6cd40bb82eac1635aff55fd6960ccc45cd17662e76f1f7b647832d3192b4bca66975bdace7e2e899d64117bfbe8313adf5277b33
cb798ddc7cb67610b59f9a2a8d79e504e186f5f1729aa675833f98a6e94fcfe7bdb4949daa8da994faa1b5bebef8ffbabf1cb16e4073
de168d5aad31636f80a1816f6badffb1b63ace53bb41595e40252c5e064fe77fbb5e10ae7df4efcd86243e585b98a9f602cae905ebef
f9cc18aee373d3bd67ee89e2c6544a55317346a4a831e931bcef7a89d436d895520fac8d69dfbf62c8b35c13c3d8fc0297d6a5e893eb
3b358fef3b76096cc3496b6382d993db3b0a8f7fdbda524876632aa51e591b09572ff291d343388e9f60524593fd98fb2ed2f79908ee
69adbff8beef3a47caf137be7faab52ed2e089a3315d17e9fa5c08df31b87eb0c077cd8f7727f49914aadb997491f7ac1d3f79577a13
f3e8c9daf193e7a1e3734f6a57ef99b9b87a5c953948c6fbc8ea990b1bef86b5f3c0e43ee61fa79fbfe60f6b0f03ae5e3905b697b0b0
3b8b0719b9dec9e68ff88ce447bde40574666d6120d49053c309e61cb27b668f5e3953ab7776eb99dc9304335b3754f3312f25ab31d7
ae5deb03e5bcf61ce1b1b5a580458738e1f891588310387b8deb47aa70fcf3a5c3aff8989c6e0b6e2dad314d56d312611d4e9a0b3406
bd52eabdb55300f38f5dede5d9dbd51b82ab677244b8459da74144c7f5dc911247632aa56ed598bb7451bd314b1829ba6e82ab971734
2623684c46d0988ca0311941633252ad319552ff9bb571e7541d014de52317d899ae199f23c0de98eb863c52c3e205c4086b63fa7aa1
6ffbde606b4ca5d40b6be3c1607b9bd7f02a1ef26d2efd316eb586937d556f08f4cc0d527ae5003dd838b1aeaffe5fc7ae28295d6dd1
6509b4ef47312e0a6be3f677d6ab2ebe98359fabcf3b57945898c73cf7cf7c3df52fe55ccb63bbb6fb3eef3b966bbbeb5c8bbff7abff
7fe2faeee5315c1b536f34f1efbdeb78aebf8cc6fc14d8676d731dd3f5b9d0bee51fcb028154625f26be17906ffb149860d9723dc71c
d3b56f496ee470ef17c966cc912b10a276845d51b06b2ea53d73b5efdc9720c0d59b7de70eedf71d734d72cf1468de5c5bfdff37eb13
8d486acca50d5602c74a605f0fd15a7bcda3da243de69cbdd2d718bef3b93eef7b34e947ff2ef63137f981ac8d9df1dd586a6cd3b2b1
1c93db3f56ffef4ddb16dd33b9b5d2d710be737245d4397ab1f3be283cd258040f5dfbac2fc436a652ca4491ddb47614e0bb20df8dd6
6a4c4e621bb3f94280f5795b34a6ef3e43df59b2a999be13ec8df50fe77a791537e6c11af27cb5c03fa941838d7934fff7d6ca8b2d8b
26a899351b53e2dbdc778e623b93dee020016f63729b423d692557cec63c980ffc7cfe476874138353335bfc9221ed715d075334dd07
adb5954228770e608db36736e0a4c7495d0d49bc9bffe1189bcfffef4bdaf2ebb3eb9e29a557aeafa56051d545faf0ad639067f3abb5
63e37b578e7180c64c3ec6fcced05a27bd3bae3466abb79ec4c6a4746c7ffb3e13b316f4376bcb8e588e587c0d11ea40a17d2e2e7ba6
e45e99f23ddf3142df5fce5bba3e27fd6d9ecae69b7483d3f0ee2938e68e6588c65c672f8c61d52b8b3310c670d198d21ef151f997c4
a163e98fdb2baef35fa9b6542eb9375512201b48055905719a59d22bd741125aebd8f851af6b38e57afe35b57914b29225675cd7f79c
ef2e5f701b3ef4e08bac899dd9c2b62cd549dffd93ff3ceabafeb5f58152121af2caa3d6b2217ddf736da35a454e7e5b7fd9f70bd584
2a4cada3d936e11825c57cdfb77d8df502e2eca19113c04fd6a5ba52278e63bf531b6b0ace77b1a9e434484ca338bef38d338cb0e447
f2ce1a153ef2c1acafbe8958df85fbaec3f7f95c5ce7498951f2f6ccad1b09e13b31931edfa891387a716d57ca88a534a6a5996be820
efac1d01d4d5da93dca562bed3f1aa24ca2bc91b971a39fc23e7ad5b1b96acac577f6ca717738ba4e1a489c5a9997fd2d870da2e3cb7
e9c964eef9868739c7ac9e3dc647a11ffcfa72e8b886c35f9473cc1a094fd68b3897dce0e8d9e605143a4ee0fc4142c78ca1468e0e4b
53178f2dfb5bd8d700f4985aabd47c78825dcf16c7db9caae45abcef7cec5c37591b57af8c1d40b83eb738de66c1e6e2c6745dfcd479
78e7baa6d891952f7eb3ba66fa2ebaf73899ce7fa5bc8eeb5a5d786ce3283b9b7589b44487d9ea3a838faa6b989b3415c8b5aa57aae7
71ade735af336755af15102ad985bbb6205c1d818be49e29711e318616d79dba447ac8869c1a5d6bcaaade611b72c165cc518dc73d7b
55efa8a12e35a33da27ae65e1a92b8f4e173f7cee455bd7b08beaa753f875fd5cb9577644af501ed2924b0c67d61893423c1c63cda12
e952828db964cf51bf5cba79e4c77ce9a8b326b67308be80f6faf299e1be3fbc8018f136a6526ab95cc41b9308166de67bccf7fe884f
78ccf9a0005bde631eb567d698bcc10be8275b6b2ba30835e6615e3a5c6b2bbd8ff974f55138d35ab36b4c4f6ac8d8a18d766e428f39
48048dc9c856635ea6e9e69c91de2b41cd9ca09b49e0316724a6312fed4dccbc8739a4abb7167ba86d2186a8c6446f8c23a567423b37
28291472be953dfa681405bbe2f1bf4a8e9df97afe071ef7abb02c10400ffd49d608c8b7f0e8e8640f275d0d5ac34935125596fb1df5
b12f9ee870351cf5d22e757c7bc2592cde9b72c2d5e07b843d1382ebb13f4aa3b2cf670656f53eb7b6ec8c6a393a6652d66b8f4ef5c6
3c12d98fb971b085f47104c8ea60cb94cd61672e6b52de2ec958559bda6bcf394d239361e0fe729b88348cf6d3536deab08a663a6ee0
929a0d1caadc2a6a89740a0bf3e8c3fa6bae847babbf4f2669deac65e6bf26fb026d777d7e999c64dd90ff6fc05463c79b3f93891fcb
1b09f5d805f767b9502aab0dd6291b6fb5d2f1aab9e0227209593a9bc8694c6ca519eab670b154d3cc9cccffae69bdf9dfeb1ece717c
6e101ec3488d2c85de244c1b9cc6686a61b986aa838ca2c7dcccacaf6bf5e48ec55ddf733de6aecfe51e9f9bd2c6b42ed0b5adf45853
85c6acd1b8256373dfaf708ff6b394d272f0928e9f33a60ea63c2f8525b19eef2d9c50d7c2f979df397c9ff7e1fabcefd82570bc80bc
c5e26864e2cde64a239ba2174aa874b8f143b59cd92aee99be5f35f5265cc709f51ecee3bbf6e550cdce4cb9c09c9ba97dfc1cb2c6e6
8b5ef1cddae9b8095f2f0adda4ef3baeeffb3e1b3a7e0d8a628d2a4fa7358967e2bc1791c349574f736d9306c6e68cb0056e7111ea81
a17d055813d8b9e4440e5f84c1d4684c5ff2fc2595cecb323d971d865deba6ac8d0ed8bd8a4cf794fc986f992325c4dc4cad86643916
f7ac11cb45796e506a8f9c297a9bd345444dea661c77735b2eb33793db73c91984a0b91dfca1b179c9316b59226caedea5960aad7161
66973ed5bc3676bff9a2511f69ada38b1bd5a4d50f5c73d6a87b439a1f746a38e1213a3e73b4a585189b3382c664048dc9081a931134
2623684c46aa35e611974d5769cca3ae3faffa981fad516bc4671e76c95b8d2ad29bdbf60adee68cb03566a80786f6ed0996c644f698
9f347bcce7b9c53d53ad58bc8bdc38f4a9703ed30437ac0b20d7a0b458fc2889f582512213c5e02f7f3cb34c505d2d18bf8de999b97f
e6eb897f8f52ceb53cf66afb5bdfe753b6bbeec1b58dfe5eb9bebffc63adbc1ff9e3c557b2f73ce63e0759aebbd9732fdfd6ab845de7
5c527b556f357c0deacbad14f8bc856f55c616599a29dd142ab93e47839f581ff2d07c046496ab581bf38f25ea47ed910dfb2f6b0b13
8e5ed59421c7e6318da6947a616dac4c52630e326c9c97d3fcb9dea194aa1a6512dd98d28683bedea9b5be656dfcc5436b0b23293df3
6f6b4b67020daa626bc0ade5c0b1d6f3e3625ff8c97459f28e51c227c78820fbcf750ed79f6f0494f3e73b8ee71e9e05f659c798ff92
0b8530716f9d41c1856f0494438b55759b8fb9a3db7320723ace9578ca4c78581ff410a399cee159216c990139301d66d97357bdf76b
ec93196ccc8aa650f4afcd4184cd79d161968de858dfbe2d4b21cdacd898b199b2b82687bdc709ed4bbd066fcfaca49533b9b98f4a61
a9c9ebc3db336b8e76521395a47cc745e84d5ebd671664cdda23d1bdd9d9336b8fc10fd33339131a4b20255955a983d0ea992d66865a
f6cc9863f83eb32c31e1fbee12b13ea0964f8869a8b9415d9d29f687bcf298bb0ed4918f9ca7de6a9079e433cf38adfe3f0a4b336b93
f9b8464da7ad49e91c6a5176476bbd9e967b1f95db73d6cc56bd52925ebaee39e4e6751d63c921e2337d8db0f51887f6b9b868ccdabe
9105cdcad8b87a5df5735ecc100b7ec453bee73bc6d6f7439f0ded5b23fe31afdd909cfc4b5aaf2c8533622415454ea2eab47ac4737a
65e83ba17d6bfed5a8c7dcb6b644d02bd23897269a195ba8a334fff0e27fcfad0f941f73938bc6acf92be61e3b372896be9b3afb73a5
10be2ffc9ab63fb1be4d2c7be6636b6f435c171f83a32173f2175b0eb79c6c5c57a6e05aa4217391db20ad52e9c6724533992f266aba
3fa7411c8fe1a984543dd67ce6726eaf8454572ef1cefa50f8f3ecbdb1a4b73bdfe6a51718337c7334cc3badb525ee731500d74dd6c8
5ce8da6e7dd8872b9a6bb9dec6150516fa731d6723e28cedf8dc51721413151fb5e7da687dc87153bdff5cd759d0902f5c8db9b8f717
aeefadff9c8ff91aee478a83853dc85194fe324ec0232796e9e4c4d5c211bfe4930a3dd378026f39ce151b68fbc575ada98f7989ac58
ae5ea65ef360cb21c61588d5e2b8d1762f7763fa2e28e5a238ce57e0377ae48adf8fb27fb91a73e9b05f51add0f0eafcd68d248ea46e
acd7a4a7ce4271a695b01a92c4bbc99a74bad9d78eeb8ac5b5263d2904b14a2684de6f7ed71bd9fa90e3f32513c35369634a6cc89994
0675dd47ec779770964a7c27c9160df532d7673dd77ec3b3dd4951a52adf854ba2e5b5162f91969ea5dad143abbd104b57f57a6b4e4a
62d5a0aeb7360b29ab7aadf9c998257b1271740a16a23573b4c7db456d577054cfdc43431297c94b6af4ce1cd3283a9b8a34b4d65597
e4c4acea5df7cad1d7085dfa99b87be7a666f60e39a941ad7b0af6cc9895ac20b231a769bab3f8f7536befb82c1f7536d90a3ee67b7c
c4676adc1b123833e26dcc75f889f50160e17dccf7fc884f78cce583c664048dc9081a9311342623684c460ed998b5b2be1ed2cea48a
d29739d8b9eeefb0467b0d428ff9376b0b08e2ed99137a6732a19e0912416332126ccc50261560136c4c90464c635eaedd46ef0c137c
9bcfe0ad1e47ec638ede194156e0167aa79be81710deecdb70c6b41f9ea4c6748434a3411724f74c34a89facc71c0deaa66404746571
001ab4a0316971001a7401cbaa5e57231ed116e55c22ed3cd0911a9575f1beaf41a783342aeb145c6041e721a8329fb968d4d7d6ce1d
5325e1490892820f5a6b9145e98e08a533fa5d92aa34ef98574efeb3bacc1f8b4d4d72a31c155a89b45c40f36da31a6f37ba76cc35eb
08a01551f5008f0ed582f096771ec5d612e5e7354fefaa08cf7249dc7d65d77d5dfefda8992f420a46f51cf7becc38bdec94af97ed39
d400c095694ff21f6376c4b7ae8c882dffa669320fd22bc7b565ddcf68bf65e84fd4ab3c1552c8ef946f7a6dafee952b03475a90fc4a
6bddaa9c531386ec98f4ca3ad35adfb5767ad8b05fa5713bb612c8a24d4c5b7cdecb7cdd50b184264395e994642f4577cac961bffafe
165f39d9cabe1f8159537eb29e8bf49d7bf197d429e9985fe64a05f4100ecd308a392b5e0b45f02525757d668def3b31c7e4a2e5b96a
319262367f0d732594514a75992bec5993af942114b375786dccf9521453e2f54ba79b623ae6e0926c4650d4f63f7c458fa46095fa6a
81a7c4e0e7e96755d620f3e027f499126aff58bdaf7fbdbff6f5e4d2e555ee6abcd82a09b19fe7b8aed0395cf710fa8e94ebf75c77d2
d45b0b8619fcb822e73d8d9c8ce7582f190eddf3fa9de5a55d9d7515d82182e61dd3d180bec6b2707dcee12b8eb986b7a1ef24559d4b
20e2fa7fc41c6ded2ff75c7f682ef3b9b54518bde23197984083e4690dc77138b8175b52c2777e57e773e1fb7e2156a9321faef34bb2
359b2aa6a7b648d65cdbc29b525a86fddbc2e312db298be73717d75f6a32bc5b5c7f54a79c0698466aaa982dcb8bd0e4b20976b86f8c
fb699ade73bc9e0305202f60cb22f733d5df03ba7ee3dafcc46d5e385453ccca82dd76cc1a387e481fc9411852ee49ca6fd2ec557ea0
4e69f85a2bb969054a4da12a3453ccd50ffb546bcd51cfbd0905039521968348148d261d938218beceff3f925a728c9e47b85f69fef5
56aff2c3764acee3b4825603f4bd86468ab93ec9302b1e1dd79ecc208a7965b6a1f73537514cc72abde58a47d1318325ab0ba5af4c34
4b7d17dea38b4e29e59a9bbb24579dd42c5ff86bd149452ebe9fd215e44c6a8734915daba5beef1cc2d11d69090f8608700dbdde4586
900d38f89496f0e0ca93eb0b52e8cd2803b845fb5d1f2de981e8b0b755277d7ff4149c312ca39452630024316cc283de91d7be87a4c7
35498d422f61d81a35ab920562074db558ce35ee3129e9d029627a821cca7541552f2012744c2012744c2012744c2012744c2012744c
2092e13a26721d1d83613aa6c36ffe992a34801d325c1a42075d6a03f9ae0913ec3c88574c5f0758702d36b50a1807b11dd3a4b64e88
26bac695fd17c840e4abbc20bced71abb222be6bc4ab9c07711dd3f783c7d2aa63f8ae131d930731af72534ccaf763271e07e1523b40
44c724fbf0bbb523ff78e89c83d3bd635227ba6fed283feeb0a54472d9d303d9b563566ec8bfac2d1df114446027359adf513dc4fa5b
7cf6966bffea8f657d51d7722ad6c601cf91c01bfa01a39c01ab6404b13c5f2ed38dc4e46aff10f8e8c9e23a628e7d8765ac20a56a45
2d6a8d927df7e03b5f6a46b5dc0c6cf4bdac04acbe7bdaba06dff77c9f8fa175aaeb5bbe9ba8784e917657c4753d4dfcfc9287d69602
22d6a39f5a5b0a6999b8f56dc66b66d7843a9b2b7f68e8f38ecf72a5ac76966559c19ea4b64965b4deaa2561ddf55cb2d9da9148c2bd
fc1d53692ee2babba4ecaeae989853bcc2ebf58650fb1cd98b54b563861abd3512922250e998736b470057e74c28b435440e5217d53a
a640a51461df52ece83b6b470057e78ce463ddbba947958e89d77718adf593d4ce469f3fb376ec14f68e894e194f6a5a405321973b95
60a826e57aff325fd2b29ea54f99536b7c5ef92ee704bbf44ec9fd835a1b050c5846afa734c3a698be1f4a12adfcd57b62991a7bf567
cdb372c2a29803bdbed9aa65f8ee7974c5a48a6e7f5a3b22e0bcf762c5f4fd40602c16199b979df2de3245f62a5db63510e34c4d5ee4
f941a79445eeefe1fade96fa998198efbb1c9eb66cc5745d10188f5a6bac4a8f9bd53107ee94a354c4e5a2b4487f3792073f232ba5ef
e9cec1d70ebd063f942ee74ab85bccb570dc87c901304dd3cd9263ac49524cdf4d0011b0c66026c21e8114ad987be8943b574ceb7a12
a3de2d4abedf64f0e33af1802445f58c84ebf74974753a3feb3a6ecc677cc74b61b363ba4e3c223d32c28dc4d6dc24fd3da0bf2f9e39
cba75c6f8ce03ca6e3c4e057820611787ea3e7d69608e6b9c9e9d75298df57df72056b3c752d0329c56b63fa465a83c29a6ccbd3192e
686963d2ba9ebf7b5e432d9c1d93821dde583b0685fb87f275ccd61dc2751d89d33c3fe602fa3e7cc7739d3bf73a5cf86cccdd74cac8
557ed15090437738cc2cb2bb4fac1d71df553527f02dc5dc9b5dd94a2d6b9ccb87ef1a4ace4feb832c1b3277929e55315d2718990a9d
d25b2903793979b9ec98be1b1e987b152efdb3b5a521bedf688fcb7c2f3ae61e53f669ad59b28ecdf83ac5d4a863f8cebfd7b5e7b362
8a4ad9574a85577828ddca0d6b0b33474cab3d77ccac9199442afd58d65c21616a0c552be56226f28f9aebfda2636aad3f8d1cbb3753
41298379e16bba39039996efd5e894d40792985d93dcd732ada78b288dca9019d92a744a2bbeb1e6f916e7ed66cb86ce9d43d1f495c7
f3636f144ccbb9caa95207099db3b19bd37b1d6be6eb8af94eea3d383d3f03d92f1f98632c9f6d34f209f3f91e78a2742ee0ceba1183
6335e44bfa3b71ecbb729da1bfd4eb702ae6950f781aad37cc1de4eed61c25d7f928077bb0740c8a58791473498fa7760bc64ef2831e
bc50a72c56c9f93c742e67a77429d191895ef323a4d16e3074924f8b4e128aac39a37b4e1aadba52aaf8ce23b9333a52c26cfdf1c6a8
9a5779ce9ff96acbbf82ebfcaf12aff381eb38ab63be777c6ff3cf752c697fd334bd75dcc7dd9836e0bc17f1b98b1216441917e41d6b
8700467a3d17967d39c9990f75115c5a11cbf2e2b93be91e824b02f770a3a6e7281557ace956ba172a7a603aa3990337a1732c1db37a
01aabd45c357e6bcf7a2b9c043b4a47a5df8aef5ca7dc1a90d314b7a6fd5522d06f382edd59842c65b8afd81125748df47c92b3df42a
ea41a6eb37ab0c5f09f48a76f9eb7df0e51f95dc31733ae3a8f38089ab52bbd9a65ba619db1cb3b48e99da19f73a211ddb0eb5d2de14
24e47aa9b5b60651a9742fa43f2dc2a73c37bae6f111bc24b1f7c899c57775dc9857b2957687a3534e02063fb127af3e0a1c8198198e
c29592aedfe39bd6fad6ea734e7f3fab7a4bae570ebfb19f8d36b43a13d3717db0dbbc4d3b66e44ddfe35e48b667426d3af283ddc4c6
5c44f17859d854e8940950e77326d1aa657fb6a085e7277802bcaef908b5f568ed5cb5fa6ea8a14c055a744a5eb6d4d3da28982a8a09
95ec4fe037e8e2e64ca575f5dd2a4b4f814da09d3f8e5053934d315bb9aa401a01a1103dfbc115288c4e29185fe794fcbb14774c9f17
60069d5206a3754e0e1b139d720c9c99957d1db637458a19ba29744a79f87e2f89bf5549f55d6f50053aa54c7cbf8bafc3f6245b3103
37735b6bcd5e5b10f0e0fbdd7c9db617b965a1bd13b4e894e239755d60e837ed4196621e3599e85e18e1f71311c10ec09ae48ee95a14
0f0037c9af72bcc6c707af72003241c73c18a1f96749e0557e307cbfdfb48357f999b505ec01a72fbd17391dd3b9107ea37a1810c086
5a8a728c2477ccc0fa615ff530208455768f6f8bab72ae13ea093c3f402459a3725f07544abdb5360290017b7491afd3029042f63ca6
af03fa3a2c004d3a66088cd041291c8bd19c07f0292a0031142ba6af03fa3a2c004d3ae684ce092ac06663a273024e58073fa1ce296d
4d09904df36c6fbece0bc0922ad345d4f9ac8a06d32ff5bc6bed0060d94f7a6614867a021f4d8a03843ae7840e0a1c88aa5a810e0a66
9aaef9a18e77c3da418c5c6501f0d27c319a0934a60e7acfda49cc1d149df4b8742f72ba95f875095ef5c7a1fbf2dd594163abbd2efe
5e591f00bb4164bdf214155dd1bcd83ca883e842fa33a8617e3c86e8986b421d159d721f0c992266b90c75619f3eb53e0886a599621a
ff382aebcac6d8f681bc01cd68a696a6432aa59ed188fabaf501d08d39245142879c7ad9962bdb1045043ab0fc0d765110a094d55ce5
d7c57c24a6772a61cca7b5172d76ceb8355d073b8e46f97bd17043e46d948cc984b2e8849f17977a7bd7752339094df98c5253bb274a
a95be6cd13b884214c25a91e9e988b3abc2d9ae009bb2165101383c8b94ac76bddc5d7957f7cf7153316b317f32b39d8211773b9c374
c849fa04faa273beb676dafcb9eaa443077298876c7d3fd334fd657dd0e654ea002696e1dc8d4cf1972621e90bad75d7948834db6014
fe8eb533913db95987f481cf340a10363515e701d6e5406b1e7429a59629bdcdbf8d9d675677deb78ec4cfbd3d7ac986ef94b3421c28
82fdcaa065d9067b61d8ba3deb1f6315a0f1d2fac2b87c5805a05c19b4986d7b73dd0ed929b7d4416bfd62d549c525b20ff07ad50937
bd5c748fdff7e2111b72a013ea9009c779324dd31b6b475b5e9a0788eb8c64c20c3527e96228a5e4b41bcdc8db1197f978faa53c2f7d
85e32339a563dc5e4dd19c2fcec93db7fa2e7359892886514acaf676bfb651bf0c56b076fefa8ce94c7f5a3b02df8b392e07adce5393
9194b2c514cbb42aaa043a3044a75cb810abd79fd45adf9a984d05a5d4eca36f312b707b1a7c8a6cb4958d4d268bb75e81a9afefd6af
d4d15fe1430d74b0c6e718a0e87d009a36028de9da294d7479ca0f6f42b7ac8d7561ed94adbd2ea33e54dd3a25d93d0fcd047682511e
13ba55c42ac082fb78d5839257e17ab7ac0f5ccdd3e4dcdf9b2e9d52f8c8f0e3e2dfc5cb2f564b38ae591fe0e78fd0f5afda3eb474a2
1b626ccad864fe2d5f49153c2e4d5fe1eb354dae72331205424ca74c185957f35737fa81aab90123aedfea941219364ac8da587ecc66
0aac94620f9830cb69ad8d2b6a287f0d9a774ad72b840209bcb8268195526c7396f44a5d2b305bd22cc7f55f8be944b190e9f3fbeae3
dedcf3d2e9a194960f5b6b9da35277381493a699ac576a83f53bbf732826256df8bcde1e085fb37cfbd2a68e7eb3b608658eb05e5f5d
497ca5eb7884b780402e9eebbf56e3fa43c733be7dc7f7cc5ba2eb22ba25126c4a67793d0fce8086d45c445b55292aba339d01258beb
df9c8130a6c6c6f5a7b4a7489a0764ac1b33f454bba0575eb5f9bec82201c9f1948bef566df0d8220739df6b4553a5e4485aa5b5be5e
4b0d5afc3035cf91706cd183a0d6afef87d6960ca86372a6943e4fec2c450bb4e85cc1198744ce52aedf350892b4e8ac77fdc6e22455
8e632691a35ca173669823de63c5c0394892f20aef9d9fb2384061b1282be9959e9b6f671145ce42ee7530e40b125bbc6037f194e695
bef8a1cc92800f8bdda794df5231fc9837ad2d0cacaeedde6a2525e7f55fd03b8f528896d521deaebc0ecfb5d643654573bdf21c0cb3
eeba7426a4162d3ba5c8064821d6fe1be5dee665cb8b4d22b225a35346923820192615b6c4dfa557a7fc362f651d85bdd68a94d8297b
952cd975879c2243c9809b269df2a0f5bad7a1642291a8e84d5edf23db938593dbe7e47d12cdea1ebb67076ede29773ec0b118e17ea5
8906921104600a221e21a7cf89b5a523504a0f11d5bba2194d2d77af94aba24b5628be5438ab99a194741a2d5edf775a8c6400002000
494441544b97dbcdc152d471bcd65e8f96eeb977e68c56af6fd7abd0fc58ad7303255318e92e7ef44dab42af1497eafdfa6e1d4f692d
231871741a8bf47b73dc97084f5bd3d1f75c4a6419fb488ba058934a71931bef686d14827166b8a681a478daba66f275358cf5214138
94c589f00e29becd7b479e5f09581da072d6e67a70a91dd2bc8d561df286d46b153179be88169fa87296c8113ab9df9c6bb72543ed39
a7387cea2aa727097189f8978321c1aae36d3449d74c6f9d3925cd30e182e2dc8c8bc1d0377a9d4b4cbeefccd421099acabae8909206
313188cd25b4ac67b3c7f2c13541c992cacc8321ea9c43e457ec855973b3281a3aec433c5ad6b5eb504d37d419cd2046f49c6f0cc374
cae957ba11458af91ec59e2e07860ff6f4a00ed52967e634c9265bc568eb7d38d9eb5b63e820dfa376c8398fe55ecd18449e0fc8decd
16744a200e744a200e744a200e744a200e744a200e744a200e744a200e744a200e744a200e744a200e744a200e744a208ee13a259245
ed9f1195d22a180ff6c5509d725e10155ae20ac667984ea9947ab6fa7f2c22db292329e55fabff77168107e33344a7f4bdaeb92bca02
1988ef941bafe92a1565415f4650cae06b9ad293801d21ba53fa5edb2b72533f03a188ed944aa927d6460fb02df78564a57c636df103
db724788ec949955638f5894749788eb9405beed3fac2d6048242a65b66f1bc11afb4054a76498de41b0c60e90a69498de01723a2557
e48fd01ce92001119d5229f5c8da98cf1dc663810e4851cabfad2d20893d4d8975ef943502760f1a04bc9b29312c1cdb0f1ff672275d
3b259635f0a1b57ec431c833b5d9e7da45bedf4729f57ef919fa634bf5dd2d117fed72caa327a93701295aebb7d68e3049833c5fa7f3
b1f1f9af4a29963cec3d95f2a3b565a7382ac8c6f0462995da29ab117bfd1c6fbf2e9df28081b9738995d41fecf7d4d762ca289c54ed
b1b5e3eaf1ae67286a5128612fa56ce2b911344d727ffe47c60ff6d5da12266914aeb57e3f1768b576fec4e5ba3da53f1f45a184cd3b
65e3c18dc46992e41fac519b6d0555cf75c2cddf03faab62b30f59716c404e576a39d4202c74ad7331576b47014d95f2c053409fac2d
8944b6ddc52b957980147a4d57a159a74c5973c37cdeee23d8b996e4924a0fe87bfaefefd69e4c7a54c56da994296b6e3861fb817ab3
d591b5d6bbf07f37e994586de86643c54fac2d07a1955262b5a11baf8a6bad8bedd051a9de29e1dffe4968049b428267c5b26347e110
51424aa95155e7b9b5259e60ba1bc954ed948254f2beb5a50f96d724a4687b19b8a48278ca8668ad5d7eec2d453bb7b6ec9c6a9d12b6
a497a44ea6b53edc5af62a9d522975d7da082ec8ec64afad2d3ba696527eb6b6802537525a436bfdccdab863d83b65c870ef09f332de
22b4d63f5c839e102347d1a7524329a54e4588529bc5a0e7a5b5d343ad8e49f6bf732500adbf71c62d98edbeb143688dcf16aca16b1b
6eb3de489916ba24a793a57cc72cc3087986123a8d599af16679eed4e51129d7cdad945eb719e8826b0a4a3c6c4a89a00b91043b6589
3950d3c6e5544a045d0016583a25ca86c8a0c65a7a93e0c091788035f9c01a2ea5445e49193847c9a9cca36a1aa4f8121c7c2d196187
28b629e14e1445f14033b30882f9ce3dad354b6e500464ec9ba4c1a7a7433e5e2cad558179d5cf5c5382459dd27313400ed19dd2f55b
52477cbfdaf62230f2fe9d23ee014ab9634213e74b7c1dd2fa60dcfee2b887ec4ee9ba1100a6817309012178c4e5b6b5c541402ddbe7
12f2dc08e80867612bad7557ef1c94723fbcef7427ceb42e2523f1e44e09ef8d5892a3a038522506d2ba64cf99e62825bc37fbc1d7a1
ba92d4291109b43b44bef552951291403f11a9301938e7317d91e6ad88ee94c2a3ca0f8d2762c7e70ebcc495a290e0c890979da02b45
29478f2a8f5e0b332096e2053a5c134a127419dfa6b5d1fad04f39ef955f9285c0446f32a179da1eab0e5dd7137b1daeefa67cdf778c
927688ed94c34f96a353fa717d3ff618aeef96b6c1e6ebdb63af8003d06b1c116353a6d67139324eef464d38923f0494cd84a279d7cb
d750c9096e46767ae4c174257f0856117311e84c7fd1b2871754ceef856f1944e01849046d4ad78907e54c6bcd96742bd02e27add342
bbaea56890e1385e04e79cd9e10ed129b9071fbe76197190e3c3755c1735eed9bb700ce169f2a9f91bad52b43c5878b13ed57e1b7895
724f9d72af4ae9b98ec7eb7535a3e11ce8c0a5e847524a4117a51dd2142df00d6442cc35cde9af68a9ad532977f6ea66558e50dbb454
4adf75704c9a6f1d8b626a83218cbeefc66029256758bd04467f9525f2a1d179821db214ab537a8a8e836d9ae525f779d9b4d649a605
a959f275d3f74e6a55aeb05edf3b7b75b3ce9f4d81f619edd51d73cc48dfb7f1f8fcc5792d5794d275612353a1434aae7af1d4dad280
1a05a8ae28e50e3b6593a9a01ae7f2e1bb86d2f3bb8e5b12fec6a294d2a73ac071b8544a576f1f991aca156823b63478217ce7afe556
ecae94204c2884ab45873c12179d5229f5ffedec9edf595bcab146982d71a9d1b4d3a24fbb544aad75cb25a267d61666f6e6d0d8e237
df13087ea194f27a8538e33403381d1a7b2d8db73ba5acf4433db4b6342290bba9ba42f7e25f472a4499c3c6abb3aa6b91ceedf43337
52e84d6a64d3b89812dad12bfc0655886523d436b51f68dfb91b4e776d4e7579bec73225d4c545c50d7787ec89efc7ae1504e181dd85
18c39e26cf6f7367a00db5494d950cc52b562eaf6c113a9fef3ba5d7b91ce80cad96153a6497cc6a14f5dfb4438670ad4230db421db2
18a394f39ff9df41ffae2fef83e32fd40edce75ab4bf193458e79b2adde3eadc8f1ce78cf97bebfa8ceb1cd1d7626d709c40fa9febc6
0a7fa0e03dbbbe53b943be757da7f57d3bfe5ed5682fd712db73dfeb432215c2d39c51dd339546be6f03a916cf1a7aa86ef826ead7d4
3425ac4e6902630feee509e54e62cf15141ad450e47cb3f9489abd501bbfff6920f9bef5d99c5436d67288cb1de10b1341cb20dea9fd
f9d89772a44231b6170f45cb24aca14e692ee86f6b871c586318373a086b87244f4de835f9ae71508928bc9d72a2aaf88122e45d6995
049560eb2411e7629f6f1d8d60a79c84bec61b774896f3c5bc791087f0136ba0b3c63494b08e79c3da92498b0e19f1aa9e24d88f9288
0a5d13f4049f73f8b74d7c64ed0e39e7d6d9ea9054e85d6487344b4016f9815c7f96b787e5bc5bafefd54586a62faa3342dec584b7ca
cbde654542a4be1d39852b29c8979ee81aeb5f62ce5dbcaeb9468734afe7a582581fb07947eab89b0e99fb1d1f494ab9b880183b898d
96e992b7ce45811a1fad1ddb344f3d9d83abbd5c6d5233e779d6720863d735b43393cba9cd512cae868bf8aecb765a2a614a87fc40aa
a806e990968de8fb9d7ddb597039c4139df85f5c4e78a6bf4fae733aaee199e0a09117ae6b96f8e7b8f66030c53ab2c8f5999cbfacd7
b70bac8a8c43f25ca4eb378c306774ec6763615bcd481774dbda01aec09582b915110bc3ee7107aab029e59583fe0cff0a45db800552
d4d3a594334df36fd6e894574e10b85160c1be1a3385c8dfaabaf7a97a3282c5e853355e8957c2294de1a8d01fb93c5f32ded7f79e0f
71a41a5e5b9820552299aa2b65888d88eb1654f3aa7084fe49a960b6057b9c69cf4e990a838a7473ed153c80afb5d6de348435c87427
b3991e43744a5370689aa6fbd68e38c455e0524a9978c99bd68e0023a8a6b829a15a50c3a476c86f0b5b505c1d1dadf5add42934b2e1
9a2ae6622c1045283b5d0a629532f31532847fd9458a2af59c426ab16c4464a7ccb01dbb4ea5709132bf5b73115b64fd1ce76fc4715d
a25edf19e940eed12b661789adccda1cfa51373d24bde77f6baab5984e498d1c3b3a9da36f769900dfacab8e55ab8dfc99295c3e08ae
68a19688e894a9f6546a0dc251a18eb95504f47b854a683de78efb77ca840e79ef88abfdcc031871df9f193ae69501e2d6efa29472e5
ae6459d4d76da0438df8d9dae1004b4f7f12f100670ff80211f5d6317df3ac5cbf53974e493ed337d60e9b0f477955c752734aa664f0
d46de1180764446f76c823d98e29a404dd66f0b2c635a5d2b45392426e1ad1785d87a9d5314d5cc0c28bb3b56af543aac7279666afef
581b121d329ed659e25ad15229d121998950cc2133b73551ca5649a48e4aa87d476cd7ea4a196ab01974c83242ed37e27294aa9d121d
b229f77c27e30a296b45b5d737cdf8ff61ed58800ec94b4804466aeb9a9d32786074c83a84da7d9436aff2fa0e35cc840e599550db8e
52cc9ebd53925f3484d7f601d5699629af841a4a6939ea179cef350652121b6a297eee92b55346bcb691d7bb1dcf3d678a0984e90a5b
a7ec513e0ef8d15abbe21d2f90ae96ad5201a6944e034c8416a249160916a5dc7af2d021fb306a9128aed7b7d74ec16bbb3b4f5d17b0
f166eb4af1eb3b10467f013a657f7c1d50ea6fc3a194e890f271a62aecbd94d60787527a0f804e2907dfef24f1372a524adf8d4ee890
a000f159d7001bdf5c07f2acdfee4af6eb7b63cd8db5561808f8b13d6f36696fb5924ee9fd225edd32f1fd66d27eaf1aaf6fe76b0280
58d83ba5c9526b6d045238735d8734bb32b78aad3397cc8457b778467885e72aa5b343fa266901482157299d5f824aca67cf4a094035
923b2555d202a01ac9af6f9ffc4f787d0f41e0f71353ee85f3f57d626d01232126109bb308fd904595c025623a25dbeb1bafee31f0fd
7e93a0df10a36f200e744a200e744a200e744a30b3590fb215e8946046ccec093a259811d329312574301090014006e89407629484fc
787d1f08df6f37e1f50d4018744a30e556afad055edf0761a404aa50ca836012a82e4a218baed091d3299d09de47a9d1022e3ae89745
077d2cad493857339e69adef5a5b0148ed5f58620ba4019b1288039d128823b753deb0b6fc7cada3c41d28863d3f25ec4a500a5edf40
1c259dd2375f89573828a2a864095ee1a006787d03719476cadbd6968082021043b58a637885835c385edfcef5c2bece0ac0162c45e8
7d1d106a0972a83ad0f175560042b0744a2822e0a4fa9410d412a4c299c9d7ab964a29315962817c58063a33e462bc63edd8e8b4005c
e9479c9d72da785da36382186ad8944e2fcff4ab46380041d89572825a8242aa8cbe431d2fd46101982a4f0979d713a3638210555edf
331b9def546b8da922605175f23cf41a9fa6e9be52ea89b5151c9eaa4a39b3a198f74c1a116b2b382cad22cf9f5a5b7ef11953456049
13a59c7eaae58f699aae593b7e01c50417345ba3a3b5decaca6614f3adb5151c8e664a39b3615f1ace233a30d831cd3be514394fb931
72073ba6cb12db980e8709f6e3d26ddd776cc7a4011238105d931144bea2af51e7849d7910ba67c8a08e796eedb0f98e57fa311091b6
8546dba109f64b4835d139774c97d177888c0ef7586b3d44cd411087b84e39958dbc6f9b7a31d656301422b3ae2d6abca4f2152b27c7
47742a40ea98efac1d01b4d6622afc833c44bebe5d28a5cc6bf9a663d715e0091a9f6192a66aad6f51877b6ded04bb6218a574e11810
bdd65a3f737c140cc4d09d7289714722ba681feca6532e2197e47753c4749aa647982602201e8a6bfe7d9aa6134c1cfc629762b9c691
7bebc3344d4fb4d60842028765218a4b6ee0b97073886a3aa644fec25764f2773c9c638f567f5fb0be12ec0933ca524abd77f4754d42
f96e7e36e80f42e9e110966508a5941966dc0f7c64e6a5d6fa85b515000150da993711576282986f4114d339bc58ae89f5b7af78aeb5
7e656d0580b76f3e9aa6e95546ff4478270310cb0d32c5738989c47b81ce0ab620c7a419bdfcb1f1d12d208e15805826a2947ac5d099
971867d32b781df70fcd873f7338554a403ae246402c0b216be04ba1f5b98509817a8ba1be5c68cef049e4fc770908e7e904c4b2024a
a95bd3347daa2ca03efe1f636d4cd3f4dfe0a1ca8756cb3e2001ecf13b1a9e6aad91cb550810cb8678e2daa4714a426fbca55f46145c
12ba5bf4f7a081b5570a164f0c00c4b2334a293387f5d7a11be1389cd162882f34faf88a39c771384450ba348c938802838de5f66a15
14ac5601f4df8ede5e03f2dcf57bd2e2888bba3dc68aa46d0f4c4810d25acb0762d990c5ca89ebf3c3133abbc9cd39a7b47108e90da4
b9e9c629395aacdf85fe921c71f43b9bdff4d662854db06f80f66018de80656aa496c9e368a8f784c25542950ac12fba8772adc2d350
45520810cb8aac1278885f2eb9f0008fe014496538d1e9f592056e2096955825da1c360464751f4582af9432dffdd3da112047241629
fa6686cda4b36a7f58991dc19c6505561dfcf5c8b1722bb1fa93ac65b1d05cdf52284f464e1ab16affcff4c2011d805832b34eddbe93
b4ed278b7fdfa17851719045f979715de73b09cc5f3af2fea4d542a031104b461c352e9e0f7d438443707e27e79134beafae47e23526
e378e1c6a46203cc402c9970bded77be96fbabb5a52326c1edfaec7bced98898ccf6402cf9d8fbdbfe74bd8186bd527828e85a5a8050
b0c6402c2bb2b3c062572811ca687604a5f1db7238b1a49a24563d12fa7b647da18cbd678c11f1b0fa44632f9ee38005efbc6f508743
8925cdf3ac9d004bfe7638694ab8c378ac6e480f177238a06692623a05e3ebb3befbbec438e21c46c1f2cf9a6b076e8e665946cdf3f8
2c951c98c5b739f430f9445fbce5b683f6f78e4e022f89f9bb4f221c716fa8681fd8e0306299d2213263f33e585b7e9d7bc80796a625
bc8eab51621847f51c33e43f8db51a5df3d160c5912ccbd80e91151ba9b50ece77d2b0c737f7240e1a7aff1db82e51198f3696455e1b
ed85b5a8ebede3c4b3fd92943c9994571504388458a6cccb14c6463eb6b65ce5bb742b6776800586de173802a525f02e740df4c2126d
0d9bbe1a21ecdf12ac7aef88670512506f7014cbd23b945c71c3da9280c94b68924d6c7ce39ac444af26cc891e529f33e1920d2bae1b
5aeb2711c992ef53fb5b41ec3d5988e4665f35394ead8d1eb6463c201ec459fee28c63c50765e58919ca5fab18b2148d116d7a483fc7
7c477aaa3012122b80dec1c345fb775b16b948081df3423fcf6cffad11cf7c2d70f404d8bd58c67680adace529d0503ec54afd7bf1e0
562d5a355b308b87347625c8d928391569ae6e734e6fc1d7459b54b5384dece7aafd6331e9f1b2e6bc69c413031c3d018e6059c67480
a7d696428c954ae2729678a49b8e58384d757ba22d20aaebf2657d9ccc659937385f262d30737a99e2fed0d1f69a442e5aac8cc3c4bc
f81ced9f1cfb49a52a8ac2b462dba2e728473abb4efe4b3f7cc8a37b410b8b69d0f021d664b3ad92ffba18b4fd599316c78622212bbb
9bbd5b96228472fa651de4589a3db841d7bb9bacdc8bf63fb7760a6351f88cd509480e3090c9d11d3cb161156c50395425f0edfdbcd6
43baa0fbf24333ef476d7fdbdad9979346fde29eb565051c3d6e763b0ca789fa60da2e49824541c12d63ddbea584a094923b0c6ef51b
31ac9649e53425689c9398df0243719b3d5b9641a18c0da76885f1a03b6a833f8e881b8ce1b5a3b67513a15c04b98bc60c511dedff94
69d8fedcd1fedd3206c50821d2bfd9ecd2b2a41ffaa3b56301de9cf5715459cca1a9057c1456b5c99de019b9ca5ec5327853e804f5a1
3027aed213e7b93186c00f9e93348ee8e009ae1f06e59065cf59a367b844188310747285d2c31d91dd59965bc30bbc2deb42494baad5
23c2efc70be528f0aee2427bff628f96a5572863c226403e14745eb5701b2c4c5eb6a637e0e8f9c5aec472ab40d89e82aca54143b626
7194104c764299b2828ed223b1ab6178e821c270a21e14c4dc3c09037e533ef0ec6c7314070f9c3a95a08cea5db2d58c967d5e38de2c
59e407383cbbb12c43491af066ac03a593bb29e0524e46a9072499d0ef8967685f62e9bb91db5aebaa39228fc89617b5034fb5d60875
2924f01cb166a01a915d0cc303791ecf2194fc0814ca894aba62b8588e2fcbffe12df75d5896110e86c3bf15b9615acac80986e24cf8
accba30fc577615952528250c283cf8b6cd5700830b0c804dfdb79764689292094059878caad52178111dc21d8ebda703119d28f42e8
21ab057ebf3212d6ef7f4095c89d9795988918a623510303b42820aa4a6421ace5168e44c2f409da78c521e22ccd307d91a3d0953c60
599616739b999879616a638e1c9c2ede55cee4be5b16436c9f50ae736ea28d571cc2b20c118acfc4f0a30cc6a1392cff0c36a2169027
3491a3d7e031d6d08b5576eca5a3685916b5a814e911a1f60cad3b8ef97d1484329e65f9638750de6b9d297f4f1cdeb20c11c8b88e00
e844a4d7e0199940bda977a8e8c8c7e12dcb1073a1fe85d5794a1f7fb3b03883998ec04fa8fd5ea3397830239d8505390be5f9a28cb1
8250f202b14c60e5289a130f2c633831b40903a7410126246e2190f33cfbd385385e8763a61ebfedf5c66a439df2728848739a5f95ba
dc84d00b508c231ccb04e16334d30158964c381c45f3443b04132431970f260bf233258399fb1684b21310cb4a186fe33c345a2c2543
520fe06521902f56318fe83702c030bc01b46e793964379ef45b3d0bed0319cc230f9a6f84e75f30b02c3b60bc94b350d2900b614807
c28c30a80ae62c9288231d00886567287b0f423c76ce7251034dd1e0053918104b001a601c8068e7b1815802004004104b0000880062
09000011402c01002002882500004400b104008008209600001001c412000022805802004004104b000088006209000011402c010020
0288650328a9eb7fbffb1b0560c7402c2b62eaa72c4ac0fe974aa9f32e170200280699d22ba194fabfa669faf75747ff0f8c78222336
00e301b1acc0c29a7402c104603c300c67642e7c1f73c4d8cf01006400cb92891cf1838509c038c0b264a0c44a244f390a5601201c88
650129c3ee0dbe433001900d86e19950bde76b8c87fcbeac2d0e009005c432835ace19cc610220170cc313504a3da8edc586971c0099
402c23514a7d9aa6e963a373413001100686e111f4102f0cc90190052ccb0d7a5a794aa97fac8d00802e402c3db4989f8ce0a652ea59
e76b00e0f04c104b374aa9f7ade62723f80b319800f40773962b843a5710837910cc8b516bfde3e8ed201158960b247ba1e1213f0c0f
8ede00528158d2db7c0431c2fce5fed15abfc78b512687174ba5d4131ae68ec05f835c27284429f5056d288b438b25059abfb1760806
56c721783c4dd31da5d4ada33784240eebe0195974cc705c6bfdcada0176010dc5cdad7cade9d83359b3569be6f9d2fbd687fd9c6bad
bdd11a3472cb3548ee69adc558d8669988b571efecc13a1b71750f3d9c7f5a3b02a4de27fdb6efb4d64fac9d4c343c87e183d6fa91f5
01de7314b1fe8d94526fa769fa9df152cfb4d677adad8d39dc307c2fc358aceeb159584abf53526576315b38d9aa9d83f840ff7d68ed
e1e38ce148afe77f507b6866a19c684aa2fb737b18b134f33f3b9befbb696d016b0bec4d85dfbcc5390c97d32cb5faadb1d68c55b8fc
9ba6e9b6f5c130af16225995decfef21c492ac81afd68ec181b3c7e28eb5e597c5c315bfe89ccfa373b00d97b5d69f56c76f320cd55a
9b11cb536b879fd073f5742dc685c2dc7544b57bb1248ff76e436e5a3d443be023f585524e03dfff9be91c2e3e3bb6d5225b905682f8
d6fac0d5cffe43a2f9d2dae9a7db886ad76249a51f9c96c08e68f91049e7c3c6f5dd67b04cb6c490e31c4e84bf186fe43a1db5d66baf
bc48762b963444e5ac912316589797c48453dd2cb4fe82d6d2e21c1c82b9b6b844be18c98a2c5dcf9e625d7661976279c0b93c58978e
79be00f773978ed29c5e0c46304b63616b0de939e11239f1f7ba3bb184d3e3f084e6149794cc636f0df767fe2849afe712ff8a73a259
8c3284e6603762394a328c5a602df14fb4d6d15eefdcfe921824ce9d7760ef73f062d98558d29cdd28c9306ae10c9b3928ef626fbb20
93d36b6b8b07c7b2c2223047dd87e1c592565060ce0e29dc2e495c8698351cd75aa7b475d212cf089017a003438b25ad411d2a6b5065
90c28d480963c9b5d412cfc199d41743f10e0c2b9634d1cdbd0615ec8b1b9177936da92508e6611c217b6548b13c48b07916dcf36323
43b17f318259da975a9ce30af89ddb339c581e29d83c13eef9b1a13182593b9d5d8b7338a8961e0eb8194a2c114309726921668d0513
59a71a338c584228e3419d713719491b72cff1dcda01866708b18450268310220f66c589c3028c5df5137b8e57748ef3c5e66fd607e5
f231f6ca181743bcb7b678e895a64dbc584228b3c0bce5068b9c8ae78e84be5ce7b8be38077b3df0d2118489519e13f72effac0f86b9
937a0ccfe753fc10375dc7a81dac2f5a2c2194a0362468a5197382d0396a5843a5e2502d465929e5b4142b3fd35517a788adee08a104
60930785d97abe2d1c45f354c4a7d57f2f30493da834efba3cef83c57fafcfcb6e03ebe74f16c3fce039a7693251065f1c01fd0f56ff
9dc3b2a297b9e620522c219400d4476b9d54979cace3b5859c24d6944929296ac0917da94be62551c3f0a3670ee2c4f13606e3fd8649
6206ea22462c69b2fae899833881588e8fc8c073f32276385772fefe1929cc4dc4309cdea0a12a71201da4f11a1f112fbc8a868c992f
fdaed4e5a8fcdc38c3ac4f09a1bb6509a1ac0602d3c7a76bfe03134349d362a942991b537a6d619dc6d43a6a4a57b184500250441547
8749544322b99550fab9ab1eb8711c39b6ad83f4b7f89d44534ceece6e6209a104c04f6480356b291173cec800f11312c024215b05e9
c7f28714a76f17b1845002b0c9a6107106d3537ed8cda06e12c9228b96e6254fac1d0124086673b18450021045b3f94a12ca98f345d7
1dda8204f7f6c6c7d6d7d955309b8a2584120059500dab28614eac3b1473bc7f523334f514cc666249e10710ca768c50a01f3820016b
45d71a56a9f39e53c705174dc41201e70024112360c5eba09552bef5dbad499abf4c49e7c64975b18450760396e5bee188431491f734
c361d4a5ac4c0bcb1242d981528f25e843ec1093e9f745d1bf04aa8a259262801558afbe4d749672065833c41722e95a9c54134b0825
00d5e01296a472ba8c25244aa99ab7d24715b184507647fc5b1ad824d40267a9199e3194df5afed804ad75976c4cec6209a114010af0
8f4954ed24e6f9e8a7d69600159fefd8f9d36e9533cdd2256b63f6c120942270542f14414eff907a2fdca42cd8e06e9384153c33df52
b3ac87a064197f043e3273aab5ee36efcd6659422801282276c1466a4ce226244029533737b99e770ac042e52ae20000200049444154
8f11ca973d8572e24afe6b523a591b412fd8d6ef0a60a45adb4da81512668488c296526a861bc17caab5ce8af98c155c29a38b62cb92
0a9e7709120536dceb773b738858d1042badaae3ce083109d399b5d3cf1bca3b1955ea7759abdcda69f354d2344c916549d98c6f5a3b
00b0fb4ace1c97b86cd9dca4d4a069350cd55a5fe4d2a4e7fb77eb036ed886e6261b51a53aeb45645b964aa967090d09dad0cd531841
72b8c7415621c5ae704b4998cb8209d1596439af1ddb78b2c8aa2e4e28a75c6f78eadc06688364cf313ce136e4dc88cafa23ad2d68a4
f03673c9e4bb5eb19225248b25126388a56b58c516104b9b8436610dd50179e48825428404225d5832fa8d68f12f852248a21ca34789
35954e928307422916d121364aa99cfc837bf2ea5f81f248c6469054992ba4a9b4e5cb68fe77f4b03a55c47b9c939368cb124229971d
5a95bbb6a652daa3563b703ccf1962d9fc9c9c4479c36363a840173e486ef694d09805cd3dbfad48148c1bd6163e92d68433f1b86be3
17b2695926c65a81c60c6055fe93118b7bb2c7b0a1c4b6e8e231ee61f54ab0b463085a9634c700a1944b0feb2095e4450b3b15caa405
1c1d436bd8d79e4750d3826623289688a5944dee9adc569040a4b2a7b5ed17503c65b4d17134efb7d67a88dc12de61381c3ab219e181
8263277d0147effb4fb95ece6b1d21a986d3b284508a47fcf03b335c685719de3356baddb3b60031586299397402ed38973efc261e5a
5b36d853107a86509a0c3b526adc000757c45229f55fc0a1231bad754e284e53324726bb99aba45a3a2942f9729017e0a159afe0f99f
8ede209219649e326be5cd5ef270669468c84e9e0bda72299698a714cf28f3597f595bb6b9ddfa226b90f10ced329e74affccbccad40
28c5f37284f9accc7e7426357f610a19f77e03423916ffa21f0c75a6e562b2ef882f6d9bfbc29db3728f8a997648bd774a708bba5583
71310ca76245b02ee5f16d040f716eee80d1632a339e99f3111c74ad514a0d110570e90d47ce3c798c90f03573edf734f23ca549b196
2194cff72694b43289e31877ac1d02b156f0c0c294c1209eef5ca11cd6037c845549ad125ba42440e6385f2956503a2cccfe0c22943f
3285f2dd884269868a1942f9e100a53192a7604cdfa1b64c2ea14dc1fe5db0c4728260f6e47ca035df39b5e2872b546556b4d1fd260d
15c989f3c8da31062999f76fce75c0d731b64aa9bb26407fb15f17f49d996e6d6a0dc3afeccc1f668174c417a52a2c56f77204affe4c
4ae5c5158fb5d639ebe245d1703aeee2059a71beb3699a5e4dd3f44fab10aca0584ebf5664e4041a8378c45b5c94182379bd37318c80
1488e4d9e861506b2a0be6bd65ecf00825275083a73fe285a4f0b7bf31424c2185afe47865771d0e6486d2d3347db676e4e11d5dec4a
2c270ccb6b205a4832d6395f6180921725d30a878b99544a9969a22f29952963474c819795593063fae1dbde2bbd72ea869774304048
1612066b42f490b4f025b0bbe13688c3e90d0f612c217ad0abd4333e0067528592bc97ba50286f4b1493a5573653289f92871b427950
922dcb353981a507e69ec484184cf352e286a5854ea90b104607668ac5f2f24070000591f8d095ce492e1031f74a736a5fad1de92075
1ab0481e86fba0210adec2362f25b58b1194c2e1e892d73d33e8d0bdfc58dc4f89503e9ffb3084321d5a2fff691d804e7f3f68df33f2
790c099b65b90696a61c6bb28253ae8b9383ea4371963d11392d221d2a9bf127f365be969e2dbf9a585e9ee098e146c64a79656d6d48
41707588262259e9daa711975b4aa175148cc4512adb30dc8759c247372ebe7c2b07348c6b2e94cbe13559f59c62f392db13ec5a335c
e1da5f2e86d60a429907fd26b142f97cd5e6f3b39fb2de5c24d52d4b173bf5a0377572303a67425cdc13591546281fd0dfd4e0dca9c0
6aac40cc745aaa1518137d21d221da432c2f4fcee7bdec49b5b9164a47f5cab3b2e1a8fc9fd334fdd7bda7398e400da15cb2717c7111
095dc57209f31ad41664cddf9100ce7fd2acb33d60867b4fe0d12e83fae966edf3520b302098e2c4725d37bc1be495bc6cf80186ea77
023f34e88771267e54ca7a86cf4944874f9fd688d835ddef0bf3769e7844f901ad09178318cb32c44e86eb4026c396b8a8494678d0ee
d7cc0f21962e1a3938c0f1109f84b9158523a7dd59f2d543876a614ac43a42149e8f7a3f400c3757a14c5dd38275e6ace0f4660aed6f
4768d83f34521c8e61c5d285f190ae057421a42f1d5f0136ef6872dd6ac31e7f14a3776a5d653bd6e27998213b0dabb9dbdecc297f5d
b5e910b58a861d86d782de7a771771857b19ea1b2ff18b3dcecfd1c3f6aac34ab1c324006ee8cc149b101b6259818e9e7c787c575458
4f1e6288121a25b4aac985a0f49d121b93c64cf7f5e7a3d2e881dfbd704e4c39437d48134c8865268dab5e7a0b3d81726839e73f9546
0387f3aed3b4c85b8ef61495de1062194f4381bcddbb38d391a9289e878de92c187d8959c9b32b6f780de6ba3434c15d4b286faf3ccd
10ca8e509da9ebabac391cbc9993e14abe7f47b8cfdcffb33182b768cf9490a407d6964e402c3dcc19b82bad577f0a711c07630d2e1e
f41b0c177e6d2142c378d3297f4331a3aef4c1307c41c561f6e16a4c1f01e684b862b2b68756c771cd2186ce51e37c1cc0b2bc6a4572
0ae5b785e508a1dc217359687aa0ef15dee167b23425585dde3942c678cb98f8e5921544ec1c562c8d55b0988be19ac85f0a24d6171f
0863152e84b364d58b04d10c4e0b954c2198922191827b2e6db87eb861788dfa2ea86a095c70e468edd1b70aa717d62f8a9c157022a3
060e2396154412950141340cdee4a6a2d92957abe8d0aadd0fc317663f87509e2f86d9104a104d695dfdd6493c564968cc7cecb9f521
1e96495b44c7a0eed6b264b624912016b05262b9895c37fd33e8fc16fdadf9b487321fbb134be6708e43acef05fd2810cddd672697c6
aec4926b9e050e1bd092124710fa6a3b763167a994fac22194a5f34a0064f6bb2f19cb002fa0b94ca4e46bc0d0962557213308249044
ee8b1ffdb82ec35a96948ca054286fa0830169509ffc907a59a3ad351f8de1c4d258930cab6ee6441670de0091985adc992ff2ef94e3
003033d430dccc4d4ed374c7da110f3c88603832cb94a0af33338c588eb60202004e94522653fe9fa98744bfe743fc307c4ebe6bed88
e7353a0c181d2a2b929c4bb3d3b2c55d22dab22cadcc0791047b244700f12c9423562c0bcbc99e6aadc5a4a307801b08667b448a65e1
d001c5bec0218060b6459c58ee2dc10000358160b6439483a74028cfd101c011c9e9f7148207121123960542798a1a37e0c86408e61d
4a61081210310c2f10ca97145201c0e1497d8e301a4ba3bb65a994ca75c6402801b8ca89b525006230d3e82a964aa957d334ddb4766c
03a104600565234f4af386f46ef1741b861764347fa7b5c67c0b001e301caf434fcb324728bf412801d8e4714a1361381e4717cb92e6
299387df78030210478600a2b4f306cdc53237bb3984128078729e333c63617a0cc373b29bdfb6b60000bce42cf92581051e9a8a6566
20ec19d67a0390c5f3c42f15d7b3da334d87e158c70a405be019e7a39965a9947a646ddce669abeb03005c3ca7a8dfe3a1996509ab12
80f6e494a3c073e746b2589ed08a040040c3670f62e9a6c9303c67491584120020895673960fad2d615e07f70200aa91e95fd83d4d86
e1180600d08f8c2930d4b07250ddb2544aa1d101180b24d376d062189e9a4aedd4da020068c91db4b64df56138bce000f405617b3c88
134bfc4800f002b1e4415475470000900ac41200b0067e0307104b00764c66a62ff80c1c60ce12801da394fa314dd3b5943bc433e806
962500fb264928811f8825003b050b4278815802b05f3e66dcd94b6b0bb8a0c59ca5992cbe6fedf080f9120078407c252fe2963b2aa5
ee5a1b01004964d6023fb7b6804baa8b65c6d245a48702a080dc2a8d5a6b24d00880146d00ec181aa999e4db37b7ee12cf5d98560e1e
ac0800a0035aeb2f5aeb5b4608977fd334dd5b0ebb2194dbb4b22c8d79ffdddae1e7b9d6fa95772f000034466cc132bce90000926819
679914bf953b490d0000356866594eb02e010003d37a054f52d546c45c0200a4d0d4b29c605d020006a5c7daf0c7d696004aa9f7febd
0000d086e696e594b714eb86d6fa87b51500001ad1452c270cc7010083d13345db536b4b80ccc4000000c04237b1d45abf9da6e99bb5
2300a57b030080e6744dfe6bd6ac5a1bc3dc574a25a57c0300000ebacd592ec918623f25cb1400009a20422ca73cc13cc9c895090000
5948aac173dbda12e663664d64000048468c586aadffc910cc3798c30400b440cc307c86b20d7db57684f9a0b546390a004035c489e5
942f98e7a8210200a88548b19cf2b2ab5f80953e00801a4872f05cc1ac05cf113ee35547e26000003762c57226d352fc8ad53e00004e
c40ec3d728a5fe8929e7b906c372000007e22dcb195a1a99940b73fa352cc76a1f004011c358964b723310c1ca0400e4328c65b98444
2f2963d1f4cbca441261004032438ae5f46b589ebae2c7708d44132b7f0000d10c390c5f43d6e2356b471cf7b4d65f3a5e3e00600076
2196534110fb02882600c0cbb0c3f0358b20f677d6ce383ed3f01cb5ca010016bbb12cd730d4ec41826100c025bb15cb1906d13cd35a
c3da04e0e0ec5e2ca79f8269c4eeb3b52311c46902705c76336719c2386e48e8ee053eb609cd699abff75d6f0800d09c4358966bb82c
4de2b1d61ae209c0ce39a458ce30841bad79aeb57e656d05000ccfa1c57209832368cda9d6fa81b5150030248798b38cc1cc69d2bce6
19d321ef2fe638b1bc1280c18165e9a1c2107d0d2c4f00060296a587794510599b2fdd9f2a626d799abf2f24d2000061c0b24c442965
1c387f343e2dac50003a03b12c4029656a95ff5dfb3c088607a03f104b46cc307a9aa63bcc87fd46b93b01001d815856842a4cde2f39
03ac4a006400b16c8852cacc3b7e4c3923c4120019c01bde10adf5a7d9c3befc9ba6e9b5e72a72ca6600002ab02bcb92c26e5e68ad9f
593b0100a9cfd393699a7e20f7c14f76390c5f84f760ad3600095084c75b1249381617ec7ace7255c80c99cf0170b010c8f959b9adb5
fec7fee4b1d9f59ca5d6da0ccb4fe87fdf2c56ca3cb13e0cc0813002698c094a20f33709e54b9a4b87503a388c375c2965de9cbf5b3b
30540707c16141ce20963782c3850e29a5cc5bf3a6b5e327afe11c027b8246516f02b774c3e441b0b6028bc3850ed11bf486b5e3277f
2c86ea6f91d4028c887170cefd3820942734e486504672c838cb458df1136be72fcc90fdfb221b102a3c029128a56e99ba500b810c25
7a99e7253f597b4010ace0891baaacc13c27e80acd3fbe0a4c29ad79a7b58663b30088e5820cd1347c98a6e9193c88a0263425f4cae3
a40c0191640262e920533467cc300725244031d40f5f39bcd73140249981580628144dc3375a7e896078b0090dad5f14a6f983485602
6219414eb6200f104f70099338ce60445319886502c6eb384dd397cc61918b7333cc42273f064a2913c3fb2cc129130396f136026299
0947625f0f1f484011da3130146af62cc321138379c9de8553b12d10cb42a81ef89f954ff3ce2c538380ca8484f1c9467c23071fb4d6
8f0edaccdd815832516188bec5e9344def4944b10aa311e4f47b526954e103436d01402c2bd0c8daf4f1bf4ed3f43f1a21d55a7ff17c
0604a017df23fa6b298a4b60450a036259110a24fe54a1e2632ea7743d9f8e3ea4a70887f9af9720ae3173918f8efedb480562d9087a
38df371ca6e760429bfe9905d5fc7b342702bda0ee92089affde12f4b2f281e5b30300b1ec0043b0bb244e17d7f2c9f3ef2597024cc3
5d5f1ec5079e7f4bb1024bf99fb5d6fff9d8b7702c20969d514afddfd334fd7b876e84e370992fd564099aa6e921a54af3bd58802050
0ab71373cec1699afedbb92cee344df7a6693a3b6483ec97a78bb2c79789a5c97963f2aa7ea45a514038b02c3b4022396d15865a54a9
04e3704a59a8a223111685f560650a069665434cf0f22c943185a18c25b2b03a4f60758ae49c1c34b3f5f82035648b0aeb9d919509b1
140ac4b211147bf97922a14c3dabb138b4d67717e2f918e2d985f345b671f3779dc3936d7e5bfa3def63582e130cc31bb0aa2c59a540
1485263d23a701e0e38cd6ea375941b318929f93c5098400b1accc6adeb1e99c14a5007bb6a3709bda9c512cecab9e4b481773da104c
41402c2bb2ca83794643added77477b1944f7ab0764dc42627a1f8d3aff4bf22fa0d80585665612164cd5376b8de478b2580bb12d211
da7f89a928baf80db0c2470010cb4a2c8572d4ac31ab7b98cb062f97135eef20aae794dd69b6089debdc57d75e659eb8367bb8873d01
b1acc03a3e7234ab66665d4ea3f43e560f7f0c59a5129452ff2cb2910f3befb7720c0edb8ff6024287eab00c247f3aea4dac2db61162
0069be6f59b6c1b7f65c3cebc263945300740262c90c593597ec2069ebebc5bf47f0aa7f5dfecf0e86aedf16ffde4bf295218158f2b3
b46abeb53a692d96eb9927e1d62559954b5e5b1f1a8f2b0980697103e800c49211ca24b3e4598bf33646b275b9b62a876f7fc7d2c95e
19f80f0fc492972bab67b4d66bf11c9565cecab9a42be884c382060d80583241c1de7b653deffa97c0f6df731ce2e9eafff7f2121e0a
88251f7baebe37c2c3b94e65b716989159cf131f79e5553720967cecb603bb3ccab4da4732c1f4778381b46d028058d6634f968d0b31
f3969e39bc3d89a505e68ddb03b16460e7f3953e2479c58f18ac8d9ae28d8158f2e0eab8488bd68e078e33b9accd3d81fed51888250f
bbb62ca916b7645cc2b127b174bd0c40630e2b96662584598d42ff2d1583bd27687559ced27109e8a8402c05f0dbd16e9896eb2d1f24
f3ef3f9552c84aed07091cfa8221b7000e65595242555fc7bb9691422cc800e135b1f8da0c0c008d9ef4ea0f81ed891c6d18be190bc9
9c060bd9ad3bb287553d9eb0a894effff0ac277fc86d1cec9dc3882525528d21270d962f68f8a6b5653006119c736bcb4fd6ab7a46c4
d7b736a10c45d7429f5ba714047e8e6459fe6e6de1c35b547f07095b47109c50fb8f3e0f5df2c27559949cc73f14081d626023bbd0b0
095ba9acc40884460dc35a4e1ba3a177d696fcf3c0db1e01c4d286ad13ce0c6c5d7eb4b6fc424c62dd8d6cf4d706b62e43a3a1d03da7
826887080e219689cb116bacb91dceba8cc8c83d92f3e4bbb5453814b9e1c555d1b280902803e2289665f49bd395612792a0453ae044
7a70be4b6b2ded7e7c4e9e0b46f28cd3cb9d238b55b04d401a4711cb58274549cd9c2d8bf4e628c3f1889012890fe156fbff315086f1
cfd696ab045fcc0b30bc6604739657c90e228fb448df48cf504471795b8873086ccc5bce7cb5b6082326f6715d22d7c786e3f10a70f2
6c03b15ce0280e95caf388cf7f96ea70a079b2605cdec4d34eb5d8cc212a39103bf245550ba47cdb60f762d9d292d35ac7ce8b7d9766
619250c6cc93bdb4b608416b1d651dd1b23f512f2c12f1cd17d5344db7ad2d61627faf3d04f057e5089665ec1bf3a9b5258fd8cef959
4a0d687a50a31c0a5a6be975ab63e79dbf4b187a1ad14eb176531d6b03fc5ec37004b18c0a058a9cf38a394e4ae7fcb3e7d0cb8845e2
b094eb85520dad758a13e7e356884e4de8659912d674c3da029a7104b18c19da707392727d342c6c6a0190488782ce2db85e280d4809
98bf43eddfccca5c5893c1f0ac156705616d1fac2d0e0e5a1e251a38787ec2ba6a87028653c390fea487b66ab807253c8e9d1fbb446b
adac8d42d15ae72c2cf8585b3449247fe404c96bad4b842cb64f21d428c0aec53236ae2e36142385c4e1e092378a39dfa0c9ab39e732
cccc4d297ef8bda640dc67d1640b62372f406afbef99239d54a7ce15122c52387902ecddb2ecfda62c99637ab848d6fa252591b07949
98877d21907f5b1f8ae76ca0e1f79a94e990357f2cdaff538ac56986b32609c6a2fd4b96bbbe645a2d55b2e0e2f01894d6fbcdff4943
9ead37f959e1102708cd036dadc8900a6ba98d8c18c797a5de5cb21047b598d8fa2689fde61cf548d32dadd9bb651933e4a91a8c4b01
dcc30d63b985b217347fc99e49aa01ac2ff1d8c41b032d096dcee11d3c2d1242d030f69eb5432ebb2ade4673d2230966d5d1ce0670f2
78d8ad5846aed0689610822ccc11e2e4cef658e59204f3b1b5431e2f2b0a65cc72dc1a290a77c19e2dcb981fbd69c7305e499a13929a
3aeb71478ba63a945842f20beb76cd153791cb717bc4250fc19ec5727338d1cbcb4b969ba435d666d8ad52b2d4a442f5dabbb3786145
056a37e21bb53f8a870966cf6229ba1093b120e8a1ed1dd271d268d82daaf6b8d6fa115999bdadfcdb0531b9396cceddc2c9e3e6c80e
1e113564e841b9d141349f923553dde2939a259eaccceba541df999cf4b026231760c0c9e360976219e3dcc95c1257057a686f91a559
d36b6baca87bf490b69c82906ee5ff436d527b786e5e88375abda40a80583ad8ab6539ac47cfbcf9170f2e87f57b468e1b73ccebad13
f7764e689b8c199e53dbdf6612ced3c50bea5641320c4ece368e855ae20e76b982272291ed079ab31a0a5a8531ff5d5fdde329d5c836
f7fe5e82b3802cfc92ca8adf1acfe70559b5ffdd95e7f894dafe8bf4e5a131abcab092c766af62b97553b7e179ac4fe472d3207868eb
10f18cdc1062058be1900e1e08657dc8a35a1cb3376009e151d88a024070fa8a238a25b2afb481ab92e24da905de06674b0ce1e459b1
3bb1544a6d7502bc312b5321816ec9bc27701031af0a27cf8addcd596e39773007569f4ae566ef092ec13b245bbf139e95abec7118ee
154a509f8a653182de5b9085888519a370b439cb98ac2ba08c92ace0416ad7273a1a5b0b3322a6b40ec51e87e1de1bc2b0a22ef470fd
55f324f80d79093d2f9df36a8a635796252c8fee5415cae957ad6dc047689512a6b416ec6d181e12cb114b0b0c4343114ba9b50db6c1
503b925d0dc3378614589150918db6e7e6546b5dadbef7d108fd7698f6f8c5611c3c10ca7a98b2af8d4f292a37e60ef02ed48093e717
4711cbad2c2ba08cdf5bb71f9641b21212c4e112ced4623762b9e1dc09ed0365edde2b2f23964132b1514e04563cb19b394b7a689d3f
2ce65dea119aef6ac0ae4af6f624f43be2f9f9c99e86e14ea1145c497178040c85afa15e0c1b58b0b1c111e62c43f33120131a024b48
b6c095dde8d084cae46e4c711d86dd8ba5f4acd50323c6c15221cb11b8cae1c572da8b582aa5e0b16b085995928af17fb4b6801c7c0b
370eefe499766459fa86dac8aa520771f925b10c9205df737478a6bd78c37d9e3c78f1eae06befdee0f72ec7f7dba26d773e674955ec
003f123da7de55282009e7020e4c75eddcb2249066aa0211e5869b02cb87874099dcc3afc73f8258cea0fc2d335286e3104a1ec87afc
db77b0a3b7f35e86e1f7ac2d365f9552a1655d2011210fcf0d6b0b48861618788512ec442c4d212b7a707da10f330f8d358435c5acf4
14ab7bc82655861976d308616b81c1636bcbc1d89583476bfd8444736bb2ff3b424d7820b17adae1d44f51edb10ccaa7e09a9f5c63a6
b00e3f2adb5d0d9e19b21eb7e20191888189c60e9f77e6c5686d0551443e1b13bd90b0028ed86de890b178c8ca3cb176fee21a0dcb0f
1f16514ac388836f10ca7c2851f39650be34cf0e84f22abbb52cd744541e340f2132d81452db437e748f6c0911bfcde1c383421ca9ac
c4ab0d27d04db2321193594635870f84320f63286c08e539599210ca0087b12cd750a884cf0388376c0194d2eb0df36151702e03a5d4
8f8da427883f8ee43096e51a1a72dff02407be8f10a37c68aeeb94f19008114ac4ccc39335e913cac7644d42282339ac584ebf9c40d7
0341ed26c4c89b1415f821cbdcf5224a052142896c0498cfce9bc38702a572d861b88bd0f011f365796ccc956df1efb4d6fff1c66700
1158d73d616aa99c435b966bccf0d1e704a26139f2fda553e2f0f91fac2dc04920c01cce1b2620960e162b81d6e9aafea2097310df96
3fb054ae1e665e9dac775736f3db5874c107c43200055aaf9d4008644f84e6c72c6b1d9441f3e9ae00f313386ff881586eb07002dd5e
7df26f01a5608701ab6e78216bf28fd54167e7cda7d1ef4f2210cb48cc5b9a86e6cba4110864074df104989f924822394c4520968978
9c409f69821d806ad07cf972c9ee37386fda01b1ccc4e10442203ba882a98bee0830bf815c066d815816e2700221901db041a9ef9675
d167e70da2321a03b164c0e104fac331af0440344aa95bd487e61ca170de740662c9c8da098440769003d58afa4a5f85f3460810cb0a
2c9c40af11c80e625904983f84f3461e10cb8a68ad9f91687e212b131d1f38590598c3792390df8ede002d98ad0363619a40f6862518
c0002ce6b74f3027291758960d2127d023b23261391c1c93e58a8412ce9b0180583666e104ba4b13f9e080d03cf623386fc60162d909
935c426b6daccc1708643f0e14606e5e98d7cdef7ff4f61809cc59760656c571a000f30770de8c09c41280cad0fcf42d38f6c6066209
406528af24d2f90d0ee62c01002002882500004400b104008008209600001001c412000022805802004004104b000088006209000011
402c01002002882500004400b104008008209600001001c412000022805802004004104b000088006209000011402c01002002882500
004400b104008008209600001001c412000022805802004004104b000088006209000011402c010020028865039452ffbb52eafaee6f
14801d03b1ac8c52eadf4cd3f41f69ad7fecfa4601d83910cbfafc77d34fd1fc5f767e9f00ec1a8865459452dafc87cef09f60280ec0
b8402c2ba1947ab53ab211cdefbbb941000e06c4b2026441fee13ab252ea91b51100201e88651d4216e4dfd6160080782096cc28a53e
6d1dd131440700080762c98852eaee344df7238ee81ca20300e402b1e4e573ecd1622c5000801c20964c28a5be241e29c60205000801
62c9000dbfefa41e4929f58fb51100201288250fd1c3ef153711a80ec018402c0bc9187eaf817509c000402c0bc81d7eafb806eb1200
f9402ccbc81d7eaf81750980702096993087fec0ba04403810cb0c9452b72a84fec0ba04403010cb3cbe5638e6356b0b00400c10cb44
9452ef2b1e1bd625004281582640f38a0f2b9ee2a6b50500200288651aa1d46b2c28a5de0abb67000ecf04b18c4729f5a2d1a97eb7b6
0000ba03b18ce7cf5627524a3db3360200ba02b18c4029d5ba8ced5fd61600405720961b28a559dc64f100002000494441541ef408eb
a1584e0080102096db7cec74ded2041d0000462096013a673347903a008280587aa098caaed9cc114604801c20967eaac75446803022
008400b174d030a632e65a1e581b0100cd8158ba69165319412f0713006001c47205925900005c402c17506ca3b864164aa957d64600
4053209657a991a792833f845e17600619f3e502b124a45b6f78880e43cfd85e100062f90be9d61b622e8f4169b550500988659f4419
39d44c3a0ce4f04e29f508bf873c0e2f9654fb7b88a58548ae7108cc08e2efa33782440e2f968cb5bf5b50adfe0f9081d61a73964239
b4580e189283f9ac838004d0f238ba65395c480e86e2870109a0857158b11c78a50e86e2a029c6e164d2152aa5f4ea2f3ae7aab1941d
df77fd89b5a895d6dadab877c83a931a80be89d65a09bf4427e6c1706d0ff0526b2d26a9492b288faa490ff85c6b5d65aa88e27697c2
64fedf383b9356b16df545324a7256c59d6aad452591f9cdda720c8615ca893abad67a8470a7a6180b486b5dd5f23659a01a386166b1
3443f15af3ea2f6a4e4329a5de1786bbdda797eb0d297dfd70c37049e9d70a4080ba9b1719d6abc473b4f0883fb1b630605ee4d43e5c
71c1dfa5ac5e3be29ca5a4f46bb92040ddcd45b4406531bb5ffb1c4bcb55295545d4a66962b7d62898be46d26c0989b88f25969d6bea
808634b0fe9a9ca3e230fcaeb5a5002a81522d983ec599548bc388a5849a3a9c200e6f9b16cb581b9ca3caea32330f48ce9913e3489b
a6e99bf5a14848c86a9740e91e637c24cb528429cf08e2f0b6b9d6608efa5aedc50d35e7eccc90df441c68ad6f917826892609651321
23a751370e2196a8637368fe6ce020f8a3f2399aad3433a2696d0c132b94a7d3347da0ffe6d275aefe2896e52eebd860354f342d4615
dce7588aca88553ebf51d88fa23f1372f588feabc88a3db3be2598dd8be5ce4b321c2e603b97162163cc7d6d5467e43989e1adadf848
adf5dd54c1ec394a3c8265b9e7920ca82b1e4f8b9031cebe76452c2995a074de69ad93a6234830538058d600951ac19216e1278c7d6e
7dadd2471146287363423f585bfc402cb9a1097771951ab919c4e290420baf2d4b9f730c61252f44382f10ca69b5465d2c7bb62cf716
2ae403f39609b4186d1c6d44933af4767c7f88f6daa5581e2c54084b1fd36831daa8720ea1153e4b42818662af96e52e4385c026e731
1fa2a579b944796f2b2dad95385c3d4c71b5dd89e54eb20a258178cb4b6245b0248a207615498da5b5e2c4d231b7ba5bf66859ee21ab
502a5827fe93688bb160aa26e51cdc56d7105548f7caaec4f2c05985f61c4b1a8dd63a2534286b9d71a23302256d77c4de2ccbdd6415
02d519c14ab39c27c2a65cacebdb33bb11cb16e9b8c010443fc005c3e4e820ea0ac97b6b2503061bec422ce96d7be8f91c6456ba2465
fe36772d778b734c9ef5e187f13e4b632f96e5d005c8988093277dde322b1e3271de92fb25de3d09ee51195e2c61515d82e0f45f44c5
5b16929d591c8cc91e2c4b04a08335d1f37a05738ad1c3e1821220ae6138e8c4d06279c40074b04d62edf02cb14c1ceeb34e91601142
1f46b72c8f1880ee050fd11562bde225e166b15e71eeb5e270f2746058b1ec5dbc4828082b214cf9026b23ff396a8b96cf7a85587660
64cb120e0d1b3c445769e1e8a9768ec0ba6b2cbee8c09062890ce85e1056729516d31298fa3808c389e55132a08372c832dbb4fc4ac2
cf5a9c03c86044cb12562548a1bae51799295c62e25e90c0506289658d2015b2fcb602c839acbe2def3b6a250dce68962596356e80e1
9e8da9616d6de43f07da7de70c2396108168e01177f3d8b9959793563783e7a13d23599658d618071e2207b4aac7e788619907d75a7f
6ab8661cbf73638610cb0a3901f70cc2873c041c316c4ec3c0901febbc076714cbf28db505803c6eacbf45162127d6392a447148b02c
7d2f9fbd9df302f1628964198013f28edfabd9a8ae7324e6c08cc167c116919821e98e52aad8cb9f98b1fe4eaff9da112c4b24cb00ac
50c6a0a7355b95e31c1ba2c0b630c39c4729f54529a5a769facbfa4098cfe67ba6163b2d18893de72d93df81ce995ad8ed239df33d87
58c7225a2c0b8be103e0456bfdb6b6879cce51d58a2d81044793f3b474aedbd462ff4ec7f4fa18485435850196e67778b8106bdf3a7a
36a45b9625c5f00fcb86450208e321d75aab9aed612cccdae7c8a12021710c211f43ad67fa5a6d2b53ac581eb806380748eeb00f6abe
f4f6b8fcb2ea3d49b62c91862a1f8825d86277c64885a8862bfc666d1180996c96785d0381613808628445a9cbd9816f14da348b8d79
feaecc01d2e7d7fdeac1eabfa906ce299d677edeaf889de79cb7e8ef3aadb7bfdb2a5f8448b144607531c870b30faabef452e7521d96
5bb225c770ce6e881b8623b12f0b78d900c08c28b144625f00ae80797b4148b32c615502004422462cc9aa44625f0006c13cb3739079
e4df9750c0ba74245996b02a011880c532c5ef8941e6662efdcd42586b06c6b323422c615502201fe37c2591e42a43fdd7bcaedcda23
102996e5776b0b0040044aa9572492b59cafbf93688acef2df3dce928a900100044222d98abf9552e78124cd5d916059a20819002b1c
2b579a4229d45a0ae5ccb54ee7dda4ab58c2aa04c04bb7d2b994bd27d6883935ab727c7f9435de57fbc88b44c1ec6d59c2aa04c04d17
b12403e6b3b5c3010964d0023659e369589d5cf9529a6076134b58950004e9e5ec883160ce33d778bb6a13056991d437969e9625ac4a
00fc340fa58bb5e4721d30aeda44119839cc5739e7e3a68b58c2aa04401609850193adc325549be8a5b523cc1fc1bd8de86559c2aa04
a08c6fcced175318f09cacc322b4d6c9155b2564236b2e96b02a0160814d3c12926d079d3989dc4efc7cf76c643d2c4b58950004e810
631995ff9486d02ce4d451ef5d97aba958c2aa04208a18b164118eceebb253c389bae6f76c6d59c2aa04609b6662999035e8d4da5248
4ec988dae56e4334134b589500442331437aad21f099b5254cb730a2969625ac4a0098e028e49518bf584b2c53735a767b9134114bca
5709009045f7f84549d51bb7686559220b3a006373f85afed5c51259d00188a7776ab600357d0e4959897af93f5a5896b02a0188476a
41af9a5369a9e14bbb154b589500c49352006c2fbc1fe13eaa8aa584f59c072539d92a180af698c708aa4d0f8ce2e4a96d59765fcf79
500e3f19bf737a888bd4b9d4665413cbdeeb380f0eda7edff478194a0c946f4a4dcbf2f08ddb1131d9a5413c4aa958e7cea15f86bd86
ed55c412566577300c1f9328b1e4c82909d2a96559c2aaecc848ab22c0155a3f37a9f18d879eb76417cbce299f0000f1a426a590500b
a74724c005352ccb23c68901301c19e51da29204a792983ba25bd03eab5826143d022086c34c27240c713f585b1a5269a961b400e664
58e782dbb28c297a0400b0893534b8a7b9de595bc2d4701ec6de7beab5b2a2b48e2a15bc09853dbce97933e08253adb5c889f8d8bad4
0b4e8ee2ac4aa8d9adac8d8dce5deb1a7ade7b0a9c9625845206f084835492caea722e634e882ded6a554e5c9625cd632013ba0c445a
633427f7d1da11a0b725d10a72707c8f395dad36c9b0faef71547b1cc5aa9c182d4b08a51010633924b1a5156a2648796d6d09f339b8
3702a5546cb6a1d42a905528164b948c00911c3e11438058c768b51866ad756a2d9c1c6b74f95d53a5f1a1b5c3e6548a01c061596269
1d88012fd572aa2ef8c819eae608a652ea51a4657a2ec959c9219648c32607c9792cbbd57bde0b1c738411dc48fd8211cc58470d3987
feb67638d05a8b7ac116892512668843f25253e40b7090588eb63a94a42359304d340c89e62763391a87dee2ef2dedd391c6d5b944e7
5e9137bc64ce0254416c5c624e5f3982373ca15dbe69ad9bd69ee9f47c9f69ad458e42b22d4b696f44004ff8ce69febcd1cbaa65e28a
c75285722ab12c6155ca43b22506cbd2863cc25121383ddb22250e3413b1abce966459960951f700e4d22d15574386486768e63149ac
6f333b118d48aa1184722a18866369a33cba2f0763e608954163539e89887230197f8c879a84f3b9f58138ccbd3c1d4924677eb3b66c
502945132847ec1c7266866dc4effe425cea43adf5ab659fa32985078e785ae35dffd428eca92ac962894e2c13e19d1162b922a5a200
099368a8ffedfb374b75f0c0b12313e1ce9d4fa971960770ee443f47474928229da4394b04a18be54cf8f521203d1fe9bfed614875f0
a0d3cb04e53c0622b1a85fb7c813e39f30a562e6d537cb3febc37ce7bc4ee7fce1386fd7e58fd1c370aaaf83b21102913e4c438ce555
240cc11df5b21e241a431fb4d68facad0118ce79de73bd788a5862ae52283b14cbe64bfb5a911ae02d28d9ef9ae440720e0de9d9d7a3
86e114160064227a4e8bd271a5b2e7daf329f3fe4fad2d0c64867215d1e39cdcc4ce5916674506d5484edada989c39b73d3b12a36b6f
6badf7fcd2188ea8613886e072c17ce53890951d95cbb1763b60189ecea66589702100d88816ca5a43f019129d7bb46cb1c93afc1ee7
e464d3b28455299ad739b5535a02cbf217d203d129a4e9776b871b964c419494273ad78458cb12d98564338050e63c4cbb0cc24e1ca1
25d5f1e6426bddfc791f695e362896c82e040ac90996dfab5323259e10953005e2154b94b815cf8701ae3179c5d708492352514a258d
004c2a346b23e88e572c915d483cd24386c02ffe4a688bd7d61620829058a2c4ad60a45b1fc87bfa93d4795be9f3d047c629968e359c
401623384172fad008530ba97c4cf87c17c70e88c3193a847021f1dc1ec0b2cce943624bf9e640d6f5d784afdea0badd3daf39f67763
2b32364a6e4fcbb2c4f0493e7b7500ecb0946f8a504ebd851284b1c472e7eb72f780f8a16a66f28c5d911163fad8da0244610dc33104
97cd08ab5b9452ff64380893f3234a26f53992f2bb6218eee78a65991a0f0680879c488addf4bd8c956fb965654143ae5896b02ac5f3
5c7ad0768653e3823dad071fd5aa9c60590671cd5902a10cb2bae5bdb5e54028a5527fa397d61620924bb1544afd5bfc44a21925062f
3ab9ed82aae9c81af347cae9b4d688691e84a565f99f1ebd318423beb4476ef991bd640427c7560a7b7a49ec9e0bb1544afd67476f08
e90c128377d82138cdd52639b65036622c7e8353670846b14072bce07bf104a73ab56e5b5b8068e0e01980112c90dcb0b33da464534a
a55ad4e748c3361eff82374e3cef06b9ce943464bb8186df0f53ee476b8d5cb103f22f78e364d323d57f2a05f904f6b0c42f75f88d7c
9583721194ae9432ce836b476f0c817cd35a8b4f6c92db7f460f44574a7d490d95da53e962ae7b192d285d7c58ca11194128899c17ed
d085c92851466a4c299c3a03732196986c16c9f9081799115b38337a51ae94a4be132d0fc473b662a444e3976bc353ebf782ea884ff0
3b15e413187908bee75ae8ad87c423ada3bf0c1d4280ac2c0611cadc6b1cd6b143f394a960f8ed20631d7d57d67196082392c1280f57
56513badf5902b7d68c8983a4ff961afc3efdce5add3af52db49ebe87b83e4bf021924c16f6e04c5cb11c3d548183e5b3b36186dba21
f1f9cf8ad620a1fc6eed88405a8a362422ed8b78ab92e22ab342cd0615caeb47104a22a56cc94d6a9b68c837922594bdb12ccb09d665
5706b12a73fbc7bb1182ecd764deef3dad75cefc6657a87ed2df29d710d367739342e79cab162ecbd270626d012db821bd954b4a8f1c
4828df8d289453e67cb26923a5d4a7f51ca6115edaae3984b2374ecb728275d983f311d60c1fc9aaccbcd721565d8510feec779bf30e
8965f6242c4867e74e9dbd3b3a2ed9432d21f2faff69eda88069afccb6366becdfb7ac35ef1b86cfc966c5d7a8de09a7d26f8396f7e5
e60f182a24edc842393574c215b697093bfa4853004d2c61af6579f9010cc7abb373a7ce502252709f3706c9661f458ea327812b534e
4c1a53bdf2a9d7b25c80d50775116f75d1f03b977b7daf3e1e08e52fc8d15323d9c9cb4a73f3d5e7fb372dcbe9672732f302f7ad1da0
98015276bd2a5869318ad3aa647e7ec810a15872d2d079f0f60526cbb2faef1025961386e3b5109d2c2377d5cacc20d30b25c3cd5d0b
e50c83c3e724e488c934c68cd56bacf98be3b698674d114b78c779111f6252f88214bfac51296592c7fc6eed88e31042b924b1bd4c8a
c14721915c1ddbd7d79a7bbd7d448be5f42b20f990b556b8d953c66c1703dc5f497580c309e51a8a8e58e724fd87846d57f3b7334962
39f1ce611c99ea9ebb1218a65cc43a3c1896ddedce9903e2482e85abb546098a42840b65a910bc142c94af4a84d258cb10cae3926c59
cec0e19387e4e12943e13aafc7b3377b9f5600f549b62c17884ffa2010b131951c153e250aa5b1260b85f20c4209a612cb7262082d39
1a521f3aa652c8a29c1e4c29c144cf2d83b6fc567236f37028a59ea2d0d936828592633ae5b530a1fc27b7e4c5023872c0154a86e117
50a1b377d60eb0e4a9b4d63071b34c426986a9d9392e39317180744f2542f90d8e1ce0a2582ca75f495d21986ebe49ab9c49ab565816
1848888e302b4c48247303cc671e8f9e8b12d4a368ce724de18a885d226df8cd1927dbfbde0ad7ad2f19620dbb54687ef80905a9bb96
2d9ed2b2c44f1256e2e4c22a961304738d98b5dfdccb557b0a25736297c7a396e6ed09f5a72f05531edfa6697a3652dbb30cc397d090
5cdc1c5d075e0b12ca57ccebfa9b878d99c80be3b5a7e13687509ed2dc248432019aebfe41fda9646ed87cf76f4adefb83ac53d1b05b
963354f2f2a85e7211c33aae8a7a2b9a7989c97a31239587d6ce7c30e4cea451b9896086a29e5413cbe9c0719812e62999c267d65417
ca4a0239233a259e641ae784102998ecc3f025147b77b4953e5def77517a7418a1a492a9ffd0757faf2094b769c80da1cca043f29c75
3623115415cb890a9f91a5756eeddc1f27bde2f3163186ec19edb9e30e4d7aaf59d4e99affae20eee710c97228cc2c5628cfe81950f3
9ff9ff8cb04291625975186e9d6cdfe9ddbad4c5564afd1fd334fd87d60e1e8ae7f76858fd8cfe4a9754c6f0416bfda8c1790e41e4c2
85b39878db8439749189b19b8ae5b45fc74fd31fb7511b5e8a0e257abd4b45a1ccbf6f55b0044b4108103391215ac946428cd1247179
7073b19c7658a2a2c50f4bc3a1b78dacb351782a6d75d49e88b12a73fbfed6b1258a65f5394b178b79cc53c7eea1a8f9a3cecbf816f3
7a10ca69fa48198e1484b21e347ad9e2f9c6fe10c395d8ee62595eb98081c38b38859286ba4fb0fa2919f3c27d31f2323a8944869e15
85f884ac4b58960e4c78d1a05666728810ad4279429eeb1fb3d5489de62384320b33a7f671d196c83fc9438b39e90fd616c174b72c97
a0dc2e60c63c8c4f906e2d9d90d5b7a028f28046531fad1d70f0c473f0a592a00e5d42bb462536d769a9a8f9ce8361782466e29e1a0b
39320117bf2f86ea104d267c62170b3de72f691a4ef4549c48cb720dd2be814a4405531f910c11149b00830b9196e51a337ca237d06b
6b2700f9dc59589b2297d80dc4c7bdb7e3106239636abd9068225f26e0e62386e857c8cde530b7e3ee62608712cb99794e938473a8f0
03209e37f4b0bf38f84f557affcb39e25d08e710739631d022fdf73b4ed401fa70d825954cd53fd7bcd45a0ff922da8d582e21e17c55
2981ec51998b4e99951dff704ce6d3fcd62dfaf315bb92c2ee1d186b2a25909e313578ee8e1403bb4bb15c430fe50be10fa30466417c
2b210724fd6e0f1aa6778ba159590d0954b22e97884cc7e6e21062e94229651ec047071650238caf464b6b46a386674c25707319e601
2fa5e10211f1d31d87154b1f0b6be601e56fdccb1ce83b4a38b1abace1b444f64527f11c76fe2d858671cea2e35e21969560ae6d9dc3
19ad8bfe22a95d6ad329efe7ee87e6a175dccc8815cc2143872433e7a0ec2494ef16f54fee1e4d28a79f6165ef4d290c0a2b7bdca8f6
d377a5d4aeb3b41be756a35a5a77688a4c1cb02c99a854a33b06248888801ec0bf1a9c0a56260348a4b15328c4a2a5507e5b542e8450
46a0b57e450fe00d6abf5a7cdf7b4ecd85957962ed64829e2951402c0ba044be356a74fb784e02790be55df3a09226b72a67b5faa341
c84d7766d1acd496d20ae2412c733199ce1b85542ceb5f230b38238b042d2f6b1cff48093ae6b6e4144e69ebf431679948c3b9af53ad
3532e134a4624ce1216b995358d797022bb1b86e3d2710cb04c89aac1d9202874d672889c69fcc5721eac16f4d6e289d24470f86e111
98a114cd41d514ca9770d8c8c0049a57c868758d86e5e2057359488ffe8ad7c4d32869b8f2b74b60596ea094fa5279150f2c49e15418
51dc931c03eb714eb12d47f41cdf092ccb01301600fda8b584f21496e418d0f0f91ee3c57e1e305f2667a0f890b5b560593aa89c3ce0
d07357a3432b75b852ff8974fcf82c3f4e2bcf778e9ae72c0562b9a2f2b0fb36e223c787b9bebdb8b5d03e213bba586218bea0e2b07b
76de4028770005b67315d0bb4373a2e2e910f7d8625d7f34b02cf92d852587c97b785418fb8e98e99990d5c765e985ceb1e0b9a48518
87b72c29a5570da13c8150ee9f8595795678b3d746b030398a8fc5ae9d97b662edd09665a5a4a628dc7f50981c83dd2dcc08abaf288c
28d2aa14e7083d7259891a8e1cd1f173a03e4cc3f2ae4211296659f1c1093196e29ca1871c86d370875328cfc98103a13c388b617909
d7e8652e99b92e7894602e3274c57026d1197a38cbb2c26a0c5193d0400e0ca3972e719805e9e5ce2971c692a1d7832f39945856c831
78a8b2a8201d8620f6d75aeba665167ae6e2942a94d39186e1cc1d601e7643284110b20c4b96f7fdd121272667029168240be57414b1
6416cad758ae08522047c8d38246fbd8325b1109fc8d8641e1afa50be574846138b350c2db0db2290d2dea252815eb860f95e07ad796
25b350de80508212283631dbc2ec3597b82a19f1942100ff291d6fa84a00bbb52c393bd6084304300e8516a6c8450f34afea123f33af
ffc51437b3f60cc62ec592313c08e9d440150a4b57b025e205f1ec4e2ca9de3047194d2c5b0455299c0b44d85a6376356749b5423884
f21d8412d486bce4b9f37f3592bf8000bb114b2a519bbc5ac0016ae28066d04b392b4467802591bb6217c370a594e9709fad1de9602e
0874a1c02179b207e7c9080c2f968cc9572194a01b25fd18d11a6dd8c3301c420986879c352739f7314a598ad1195a2c993ac93b0825
90000da773d6655fa38cffa022c30ec32935fd1fd68e34e0cc01e2c88d13c670bc2e438a259343077194402c990e1f14c8abc8a86259
7ad15899034453601020d94b2586134b8ea58c18ae8011c8cdb48efe5d87a11c3cb49e1642090e41ee341147b95a60338c65c9144f89
f5b460289452660ef26bea35c328e06724cbb254284f20946034a8cae169ea6523f6929f21c492c2844a7887256160543293e45e23ab
1430217e18ce30fc86e71b0c4fae771cc3713e46b02c8b86df104ab007281c28393b5187ca90bb45b45852dab5126ef7bd0300f8c87c
f17fb4b6802ca45b967f595be2794793e300ec89d7a9f742357f402162e72c73037267305703f64ace0a363c0fe588b42cc9a903a104
c04d722a375897e588b42c0bd77e3fd75a97861a01201a5897ed116759167aefce2194e020243b2fe1192f439c65596255e2cd098e04
accbb688b22c0be7555e5a5b00d83739d62556f56422cab2845509401a19cf0c56b46522c6b22cb42a117c0e8e4aaa67bc28c5e19111
635916589548a50f0e4dc6b383da531988b02c4bac4a082500d3f3c426f8ddda023611615916589578430290f70ca1564f22ddc53237
13f404a70e0097642c0f86a3271109c3f02ca144a8100057480d3887a32711099665d605c0aa04e02a19cfd253ad358a9b45d2d5b254
4ae5967a805509804d6afab637d616e0a5ab6509ab12005e529f293c4bf174b32c95528fac8d71c0aa048009a46e8ba79b6509ab1200
7ea8144b4a850178c523194d2c4f33cb820270183014af439761786e1d70082500fc2013511cbde62cffb0b66c935c06148083f221f1
b691303b822ec3f0cc2138966701108152eaee344d9f53da0a43f16d9a5b9605437008250011e059a9438f6178ce10fc9db50500c006
59a32080c852b86b905908806452e72d5f585bc0159a8a6541203a00208dd4e9ae87d6167085d69665ce7c656a6253000e8fd63a37ef
02f0d0d41b8ed29d00b403c1e9bc0c31670900a80fd68987692696993f446aca2900c02f529d3c10cb00cd86e14aa97fa669ba69ed08
80610100f990819294b312cf9c9f966289f94a001a83794b3e24cf599e595b0000a0134dc4327375c0336b0b000074a29565993c718c
383100daa394421a440fadc432673d3800a03d104b0f52e72c315f09000fa9cf12c4d28354b14432520078f8917894fbd616708148b1
44e17700d8c0dc3f13d5c51213c600803dd0c2b2845802d00f58964c4814cb6fd6160000e84c0bb14c9d30c67c2500401c121d3c1836
00c0041677f0513d910616f203d0173c833c20f92f00004400b10400fcffeddd3b6f5d4dfbdff7993f1eb84880900202d8400a91f00b
10d5a513f50a4475e948764e25aa73e59b2c5389ec5c1820f90a4416869b0024cb20854820450a03e40692224841b248603bc504235d
4bf7d29e7d5887995973f87e0042cfb3b76e711fd7faad6b66ae41071c2c818231cfd99fa0074b366e07508ad0c972d3b965351a6800
48526a97e17d17fd03588dcb704f421f2cfbbe51cc09039024067880b2912c3d49ed6079efdc02600cfa537a42cd12003ae0321c003a
e06009001d70b004800e385802851ab8d491e6db4b70b004ca35e460c95ce725421f2c79e181e91c0df8cd7c6797205902e5da18f0cc
ae9c5bf053e88365df7993ac360026648c61aef312410f96c61856e40013d05a6ff1bafb95da6538c912f06348bd122b50b304caf465
c0b3ba736ec16fa91d2c59f40f4ce794d77e39b6c2050aa3b5b63b143cf77d567cf756e3321c28cf31efa97f3192e59352eaad73c712
9cdd8071fa5ecd35f8eead162359b2220088eb6cc06f3b716ec11f6224cb3da5d477e78ee5de303116f0436b7dda65649c54b95ef064
698ce9bb7ceac0b905c0d0efdf913d10cac17048e28448718067cfb905c0687307cecbd6bff7caabbb5ef0cb70c5f4212069b6efa531
86b185353858024007b12ec3af9d5b002023b10e9617ce2d2b0c6c870f00c144b90c57fd2fc5af8d310cf4004846aa074bea960092c2
da7000e820e6c1b2d784583a3d034849cccbf0be6da32e8d31ace601908468074b45d221a9ac0000200049444154dd1240c6a8590240
07b10f965f9d5b56d05a73190e2009512fc31597e20032c56538007430c5c1b2d73a712ec501a4608acbf0de3bcf71290e606ad10f96
8aba25800c4d55b37c706e5941f6110180c94c952cb9140790954992e590dd1b592b0e604a534e1dba746e598d3d42004c6692cbf006
033d007291d5a474ad75afed2900c097a993a5dd3ae2bb73c70aa44b00539834591a63ae9c1bd760450f8029a47019de77a0e7dcb905
00029bf432bcd177a04729f5d118c3e83880685239583e29a5de3a77ac40ed12404ca98c86ef38b7aca1b5de5dfd3700c09f2492a51a
76294eba04104d4af32cb79d5bd660641c402cc9244b45ba0490b0d456f07c746e5983f66d0062482a592ad2258044a5b8367c48edf2
c9b911003c4a2e59aa81e95229f566489f4c00e822d5ae436f9c5bd6ebd5791d00fa48f260290971e6dcb186d6fa78f5df00806192bc
0c6f30d8032015a937ffbd766e59436b4ddd128077491f2c8d317bce8deb6d68ad8f623f5600654bfa325cfdbda4b1770f4b2ec701f8
94fcc1520d9f4ac401138037b96c5836642a111b9c01f0268b83a54c25ea3dd8a394dad75a6f39b702404f595c8637b81c073095acf6
0d1fb26e5c319d088007591d2c8d31b661c69d73c77a763a111b9c01182cabcbf0c6d0cb71a5d4a13186411f00bde57ab0dc1cd13863
5b122a0074965bcdf227191d3f71eee8e631f6e30590bf2c936543066e369c3b3a60841c401f591f2cd5b8fa25074c009d6579193e67
d07422c59422003d647fb094c19a33e78e6e36386002e822fbcbf0866c5af6d6b9a39b5763cce6040f1b40268a3958aa91f54b0e9800
5629a166f9dbc8011b2ec9012c55d4c1520c6ae72638600258a8b883a54c58ffecdcd1ddc6c8cb7900052a3159da03e6d58811f29f38
6002682b6a80679ed6da1e343f3977f4f346d22a808a15992c1bb23be48373473fcf5aebdd091e3e8084147db054bf0e983b765a9073
473f37ece703d4ade8cbf0b6314d375a988b0954aaf864d99083dcd884f973a45cfa6902a84835074be5ef80a9a48e79ecdc0aa058d5
5c86b779ba24575c9603f5a82a59363c26cce6b27ccfb9074051aa4c960dadf5bd52ea9d73c730a44ca0605526cb864c2bba74ee18a6
4999a7133c150081559d2c1b3258f39773c7381f8d31ec550e148283a590553a37ce1de3b15c12280007cb9691fb91af620793b63868
02f9e260b940c08e43f6a0b923fb0601c848d5033ccb48c7f5bb25778f61e7763e32dd08c80f07cb258c31bb239b08aff35d0e9a34e8
0032c0657807111b011f1a63387802092259762097e5a33aaf77742e69d3fe1c45f87d003a2259f61070b47ca591bb5602f08064d983
9dfa2307ae936c1e34002f48962368aded14a0b7817fcd832ccb04302192e508c6982da5d476e05fc35a732001244b4f422d97a45e09
a48164e9896d9a2107b6c3229e10803f70b0f4ccce939483a68f09ed3e1a1403f080836520c6982b39686e8f38e851af041241cd3222
59dab8dff53752af04d2c1c172025a6b3b8a7ebb6eda11074b201d5c864fc0b668b3d38e5ab5cd59752f0290190e961393dae656abbe
792d8fc8d7de40003ce0323c03d2fbd2d63b9fe4cf2b1a080300629273d18152ea9304fa237600c13caec2332057ea9bb244dc8e0b5d
b53ab9d99f7bdbd14d9a22010030980d9076f296d6faa57dae915d91f66cd5d8187340a8c422542c332641f274c5cc4d3b547425154e
0e000080dfa402d9fc6c2c7865665295bc72ee019620581644f6a1385e7280686b02e72d43ea0050b60e01b2cd6e6b7cc0b90143112c
0ba5b5de916ae6871ecff081399c00901f19c1b2c171b763806cd8264fa7c69863e71e6000826525b4d6f6a071d4e360d3366b0da9df
3af70200a2901e4ccd229a77037f2755490443b0ac901c988e57cccdec83617500f0ac357cbdeb61ff605b1c38b67b7838f700be3fbb
044bc89694c73d87cdd799352bd899f80d002e99b2d404c8a1d5c7455e652ad4290b37111bc1120eadf58104cdb157c9cbcc5a554e42
27806249786ce63dfabc789f77295549468e30298225d68a1034db7e573a257872b50d2069012b8fab1024912482257a8b1c34db5e25
74de4ae8bc77fe06007836b7e2dac79cc7210892c802c112a349d03c083cccd3c56c2e78720006b09604c7ddd64facaae332af326a43
9044760896f04e86858e3cad3af7c906cf7b2a9e405da413c64e42c171de4c16db5c30fd07b923582238a9061c48d89c6208a98f0709
9ef7543d81f4c9f1a51d1a7706f6eb8de95a566cd31718c521586212095735bb7890e0f9f387930310861c27da3f534fb719622673d2
afa846a206044b24431a02dbcae6a742de95592b80daf0f9440514b56b55189b9fad4c03e322afad216dbeeba812c112494b68615068
afad10fad4aa8652e140f2649385cdb9b098da3c46dfec77f64286b40991802058223badadce721c46f7c99eccfe4fa5d4ffac947a91
30fac2a2240c258b5cda3f3b12184bbfb05b672621924a24b006c112459013e29104ced41708a56026c1b4f96982a9a2529a2fa91caa
56206c42e2660515445faea5d50f732281010896285aabbab997c14ad19c354155b5c26aa3bdb8a9ea8aaa2c46d96cddb4dbfadf4d08
5404c128a8420201102c519dd62e1a7b052d1402b0980d90ffc17ef59552ff835c60dadb8e8c31570bff0b0083fd132f1d6a23c35b4f
adea90921642ef95529f65ebb4191f0c202bf63bfc5529b56d8cd1ad9f2d63ccff688cf957c6984d7bbb4c9bb9d05a1bf9b9900b4e00
2351b1443564fed9456b0ea65dd5b9d7a50fa5cce16cf60aa6ca094c6326f31f6f7d551bb5d6b6c7e45fad9bec1ccb03e65702c3102c
5134a942dccecd57f37ae22074025e357bfe3701324ac093f9afb77373b14f8c31c7ce5f06b014c1124592fe97e773cfedd2187310f3
f94ab0dd4d788f62600ad7a9eed9bfe462f4412e4669e505ac41b04451ec5ca905fd2daf8d317b293ecf56f06cf63aaebd5f20f2376b
ed36955c70ec4a4622eee72a9876facc2e0113588e6089222c0994f624b095fb5ca905e17387d6499848334c5dcd3ef94b463f0898c0
12044b646dc9415fd53837aab56bca6e45dbeac18fbbf676a23504c6beb4d6f70bbe4f3660eed00713f81bc112595a324cd5784f2561
3d59acb0d50aa15b54438bd034abbf6deda8f444f8194f6b6ddb147d5bf00f253bdd06888d6089ec68ad4f95525f163cee99540f6813
12900ccdffef4aa97fbee0b7bc124c477b9040b870bb4df959b46de99931e6c8b9155e49dbb29b25ff2617b5a81ec112d99040b3eca4
3ab38d909d5b118cd6fa76c962a34101bf3594df98ffff4ab63adc71fee33fef8f35fc7fd7daae727ee8d8fb7eeb6b020d6d712292cf
eae392df48c047d50896c8820cdbfe58f2580995135932efacb13dd5f0eb82a6d721440b736b9e0fa172026bc2e5833166d50510502c
b67444f2b4d67b2b42e5eb9a0a16029293e7eb92dff02855368c201d0f9685ca4b42e534e4a2e9e3925ffe4e6bfdc23691a811c11249
9355dfdf573cc65de6544e6e5578bc91f7100368adaf16b4d16acc6237fcc79f64f5fc9973c72f76aef113e112b52158225952ed5ad4
4aa871c644f9e9c97bb0ece46a9d53b9ec4f2a95abb60825542640e653ce963c928dd6c22ba00a044b2449e62f2d5ba8d0600830111d
162bdcc87b8a0ea4adcdb24aa592f636f49a4cc7aacfff5b59e8065481608954ad3b105f32049e9c55554bebcab9050e59a8b6a85762
dba9730b26638cb95a51b5b43ec8c502503c822592237d2a17b5146abb706ec1d4d605c7779c5c3b591bc0a9562669dd31e91b557bd4
806089a4c8817751f3f33f70624d4fc7f7e41b8b199693e0bdeea2eaceb90529e8f2f9a7d28ce2112c911a2a5a79eb127a788f97eb32
6f988baa7c7d92a90e40b1089648cdda6a25159bec112c17907ead6c87593e56f3a368044b24a3474b9a45db08221f1bb41f5a686fd1
8d280ec1124523582225848dfc750dfdbcd7aeaeaf09af5dde36180e47c908964849e7132615aff4f47c4f78ff5ceb16ed34a8d8a7a9
cf679a056c2816c1b212765b3dadb5dd5eccacf9b9c864d52ec349e9212c0ed4f742890bab24f53926f1fea15804cb82d980a8b5beb7
4152b646ec5211b1bb7d3c4bf84c39bcedd3b626397fd5fe0244c402a884c8c2abae15e7a0ec63b1bd80ed6e3fad9f2bdbca8a3e9a88
81605928e987f76c9b528f7886768fe794f7b9654bc74448537b0cd7f77bf689aa6552fa7efebded1a66c36253409022c277e9aef1a1
f5f34976737a94bff79278e10019235816480e18ebb684ebcaee9692ea567c5f98043fbdae4dede7a47cc112ddc0ed49d97d2a015aeb
e301d54a2f9f7fb9f07f1c5040d890c281e10205be112ccbe47b982c56756448e3e75b86c42737e47d7b726e41dffeac6fed9c68e756
4423c7c5de5340c6ee1c26d39c5e468e48356ea85ec227826561e440e7e36033852107db0daa5fd3b173b806ce2d4bb50a3ea5219fff
7da9982132a9d4df0cf8ad3e367838f5dc4cffdcb905188860599e104d965f63eccd2dbf63e6dcb19eaddc50018b4c42e590d63777c6
18de2fd7d0b0fd17e1322e09958f037fa98f63a9f74538322f1f188d60599ebe73ddba8879c019ba08e4ad4c4867583c0209f243fb29
327cbb8031e67e4435eb2f86c5e39051a1a1a1d2becfa95e04301c0e2f08960509344fc65697a29db08c313658be3a7774b321ad9298
8c1e885d2c252b4f87b65699c5fc3c65684ce8d8975eb55c5c052295e121c3df8d33e79661425cecbfe3d8091f08966509112ca7b88a
1dfb3b6fa8def8272d857e8cfc87196e5b41a6835c2fff1b6bbd958b2bf61df7a8e9093cb257ab9d52e4e5f32fd5ed4be78ef1a85a62
34826521a4ed8eefaddecea6980b678cb9f270d0dc97561ab4231a49aa942f1ea6595cca7b8b158c317b23aaf68defd2db90eae5489e
7a022bdfa1cd18132204ee3bb7003d112ccbe1fb20e3edea7a0839680e59c833ef07732f876955697e785881fa10e844582a1f4392ef
a47a49f57e003b2c2cd33e7cf4043e09745175e8dc3212ad873016c1b21cbe17eda47070d9f150b951adb997f75430d76b054a1f551a
251729bcee3dc850a7afd0d054ef09981dc896882f23e752b65d865ab023f3951f9c3bc62158621482650102cca7ba4e61c8527623d9
f2142e9584a4a682c91cb43952a179f1182895bc77ec4f3c8084069f15a92660deb367b4cb2ecc696d89e8ab47e465844abdef7fff03
17e01883605906df43d6c92cb008102e959c34becb49b6ea9d7ba43a792527d41bcf4d97ed5486ad81db15224cb85472d1d0ec195d75
ff4b993ffc249fff310b731689112a9beab6afd5e60daa96188c609939a93cf85cb473925af36a1b4c8c319b01867c94bc76cf7292bd
aaa1922361f2424ea6b63af9c9f94be3d93655844a0f245cbe0ff4cfff259ffd975a42a654e6efe5f3ff6344ebac55bec69c532cf3e1
7d5e7c87e8878c4a102cf3e7f3e0354bb879af92797abeafccdb3eb52a394f254d629779634fad301972f5a7bd38a11f9e4752957ae3
6941db221bad90d954f28b180e950ba96309cf4d653ee4b6b7efa51f6f6c5e479a58c483a10896f9f3793049bec7a05c997f74eef0cf
5631ce5b27da273939253f6c2e27d2a32648b6e68d85a8ccb4bdca4995ed050390cafd56e08babc607998fdc0e9a59040d3bea2015f9
97d685d45f9ea7792c623b1f68b908884e2adb3ef6216f102c3188fd12f0ca654a16a07cf7f4e8afa57f5e36ecd075a061dcae5e657f
67dbd4fa2af6b0af84dc5dd91f7e2fc289739528f3c9ba92615ddf73e6e69d4c15a265cac6fdc4eff94c3eff57d2d83d2a790df6e43b
30e571c0fa9cc282c7917b982ff27eaaa08c7c112c3366ab081ee757bec9713e9c0cd7dd4e7c825dc556105e2404347fb635f359e7e7
76dae7b5d9fad377f37b5f5e535ca0537ab06c4815f1dcb9230daff2796f3efb4fadcf7be35e3efbed9180e673af24346e061eba1e23
a90b2af5f70e59bee6484efe19477e089699f27c65fa75a23941de78aede623d1b1a7653ad66d4122c1bd2a3925d53e27990cf7f9217
e3d236ccc7c5f6ab2c9c043a638e65be7c5d25cf720f954ab681b4f39b42ec44813f34f3283719224b87ad9ac9e73fc4fed1f8db838c
eeec243ec2e3ebfcb041cf5ff445b0cc97af03475113b4ed047639c1bef7dc7ea376330265fa5a01930b2cbf2e65614eea81f22799ef
79eddc314cf28b3a91168265866c1f364f2b7c2fa798741f830d3f3284f3862ace28677242ddca2c50563d7c37778115a2ff6b0d5e65
518ece749f7b5f81f0033b35a10f82659e7c1d308abf1295162d079c647bb10b8eb6e5849add67442ebc68f0fcf705d68e7cfe3f53c5
efe4443efb9b29acf41e4a36ba38f1f4cfd17a089d112c33232d667cb4d6f85adbae287327d96dcf3ddf7277dd0a93bba9edbed49584
ca9b3c1e6d5c320f79b3153243355bcfcd6b2b4cea925641cb73f1f13e331c8ece0896f9f1f1057f2861c1ce1836384980d272a2fd5a
5935c79e6c0e5b27d3bd5cc3648350d99d84ccadd645566dd345ec45e5c75665b2e4963a3ece192ce2416704cbfcf818926058638e0d
da4d35a7b5f8a1a461f3bbb920b9253b751441fa39122a07908bac83d6677f5b86504bbad0ba94c567ba55952f727ef93c8f0b79386f
a013fa5866c4534526b986beb990f0b297c02e1fabbcb67643c9767e581f133609afa679746b979b83849b95abd66e4017742ff89b4c
a17a76eee82fcb8d341017c132233e9a204b45021ec94977477609d909bc4b4eb39bc96df367cd07fa488dd097e122edefddaf9acffe
4ee0e0396bede6735b4bd5d107adb51d12ff36f29f62271eac45b0cc88d67aec9b7558d2f027ea3671a86c7c24dc20175aeba791adea
ec861ab41ec24accb1cc845c6d8ef140a84429a47a3f75a8b46e64281ec8c1d8cfea5b9992052c45b0ccc7d80302273f1421c17db1cf
a57a0a244daaeb633b00d07a082b31149e010f8b76ce726c740dccd35a5f25bc788aef1992e769210f8b78b01415cb3c8ca936be72b2
4309b4d6b789afc8ff22d55420591208bf8e7c7c9c53b01415cb0c8c5cb4c3821d644f4265c8d5f63eddd93e897cea9032adf5fd8815
fc76cefe8e732baaa7a858a66fe4a29d3b422572272b59730995d607396903291b3312f68e453c58866099be31db68b1600759f3d01e
652aefe4b103499206f26316f2707ec1420c85274c1a0fff18f8085948806cc902031bcc36327f176d43fb2d163a205563a65ab1e106
16a16299b6a157842cd841b60a0a954a9ec393ecce04a4e870e863f2d05f1905a26299b01157929f6bd9271a659100f658e8dbba6d8c
61781cc919b19087453c7050b14cd488dd3cae0995c851e1a1d27a64c1031235f47cc3221e380896e91afa45676802d991f9c42587ca
065b402239b290e76ce0e31ab3c0140562283c412316ed9c1863d85a0e59f1b0b3548ee82f8be468ad5f86cc6d66110fdaa858a66948
456346a8446e2a0d958afdc591a841235e54e1d146c5324103af1a59b083acc8c9e8bcf277edd218c34919c918b8cb15bb4de1372a96
89d15aef0d08952cd841560895bfed6badf9ee2225432e743ec8142e806099a0214311543c900d19022654feed93548980c9494bac21
0b79380fe12786c21332b0ddca5763cca9732b90e667dc86cabf9c3ba0649e348dd491842153b258c40345c532397daff866844ae442
6b7d4aa85ce9addd5f5c761e02a6d6bb02c9221e282a9669197085f8d118c3101a92a7b5b6ad75f679a73a617f712441e6ff7eeaf158
58c4032a96a918b068e79250898cb040a53b7b1ce0bb8d14f49df3ff817df141b04c47df2fb05d4d6ae885871cd8ae0532ffea92376c
255bad7ccffecb48812ce439e9f950180eaf1cc1327f7f49c07cd15ab39d2392263d1bb72540e14f76379e4dd95e0f989cccf7ed3bb4
cd79a87204cb74ec8d3cd9dae1b36fad90c9fead4892ad82d800653b1af00efd64a7b568b678440a6c98b4732bedb94429f53ca059fa
86eca8854ab1782731728578e471f5eccc0e4d301f132992cfbbfd6cbeabf00db2dfcd1d16e9606af23db4d3aabe787828b4cdaa1c15
cbc4d8938cddf35b2a18764eda47bbd26ec4a37c6bf762964ae63dbb232025f279b79fc9f795bd31b6a3032bbf31293b47bf5599f411
2a0f0995205826ce561a6dfb8656d0fc3a62c8dc56857eb44226070024c1ce2bac6471cf997c971941c024ec5c7c992e653c8c8cddc9
4592663a071a0c85674c82e1b187fe80f6e0b047f504299061b9a7017be6a7ecc12e82e03b8629c89cfb0b0fdf295bd43866630eac42
c53263b208e2a055cdfc2cf3b6fab293b39fa59279c1ce1f98920c8fdbcfe061016fc4eff641844ac46417d0c84e4eb67af47d44a8b4
a308db729ed92454621d2a9685f2b408c80edbd13a0293d25adf0e58999a824386061193cca1bf18b918ce16278e6cef59e71ea00382
6525e480733ae2047d62171539b70211c8e7f73693e1f16b630cedbe10854c89ba1879f1659ba09f5255870f04cb4a69ad0f2468f63d
5133c70693d15a9f7a5abd1a027b7c230a1991bae8b98f77db9d1cc7594406ef089618b308e8557a643264826812ed7df999ef01421a
d96b928200a22158c221bb265c480fccae68c48ea864a5ebf7895ff54bd9a61208c2f69a1c3857fe52c2e493730f1010c1122bb51601
1df518367f9090c99ec708ce6e3f376248708cffc918f3af7987e19bed3529d5c93e5395b8b84712683784955a3b016d4a4ba3f71d76
026a3762bfa5113b42928532ef47eeb53fc47fe28d852fb602df6a5cfead43a87c9545376fa415d016a112292058a217d921a5bd13d0
e19ade9976a5e2a384cc2b7a642204f95c6eca8916c882f49a7ce9d16bf24efaa2363d258f592c86d4102c318aedd32757ca3664be59
b325dfa75623f65342267c939658f673f87ff1e22245b67556ab71f9cd9a30f92afd509b2d1377996284d4112ce18d0c9bb77702fa28
f32d17f9d20a99f4c7843752c1f9b7bca248859d0ea4b5be9730f963cdc2c8f99d6e68b28fac102c118c9def235bd935d5cc9325f3e0
fe9280f92293d601206b764446e698db30f9b8a23d96bdf8fed8aa4a1eb0921b392358228a058b80b6172c02b24342df5a2193dd4b00
6443c2e49584c9e725bbe1cc2fbad961d10d4a42b0c424ec15f99a454036647e6f85cc5dde2900a9913079da0a938b5a5fddcd0d6fb3
e806c5225822090b16019db51e970d99371232ef65df6800988c9d1bde0a93f3bbe1cc162cba61781b5520582239326c7e34b708a819
366ff7c8bca747268058ec1cf0567ba0f9dd702ee77a4ab2e8065522582279b208a83d6cde2c027ad7ea91794bfb2200bead685cbe68
d10dc3dba81ec112d959b008e85a26c93fd3881dc058d2b8fc69ae71f9cf4537ad20c9a21b60018225b2268b80f6e61601edb41bb1f3
0e0358471a97dfb71a97bf958bd63f16ddacf96780ea112c5194058b8094044c1ab103f883342ebf6d352edf9c5b74b3c7a21ba01f82
258ab56011d0ad9c4468c40e546aaed7a46d5c7ecfa21bc01f8225aad15a046487b44e6585e7138dd881b2cdf59abc524a9db6aa9247
2cba01fc2158a25a365c4a85e24a86c48e69c40e9443bed336341eb5462f7659740384f30f5e5be0d72220a514733081ccc93417bb80
ef5816dbf0bd062222580200b26787bb655e359d20800931140e00c81ef3248134102c010000e005c1120000005e102c010000e005c1
120000005e102c010000e005c1120000005e102c010000e005c1120000005e102c010000e005c1120000005e102c010000e005c11200
00005e102c010000e005c1120000005e102c010000e005c1120000005e102c010000e005c1120000005e102c010000e005c112000000
5e102c010000e005c1120000005e102c010000e005c1120000005e102c010000e005c11245d05aff1f5aeb57adf526ef280000d32058
226b5aeb2dadf57f564afd774aa9ff4629f57f6bad7778570100888f60896c69ad8f95528f4aa97fd67a0eff504afdd05a1ff0ce0200
10d73f78bd911b19eebe574abd5df1d0cf6de5d21873e4dc03000082a06289ac68adf79452cf6b4265e38bd6fad6b91500000441b044
36b4d6ff5e29f5bde7e3fda0b57e726e050000de31148ee4d9053a4aa9ff4d29f55f0f7cac6fb5d62f4aa92d63cc8b732f0000f0828a
2592a6b53e92053a4343a5913f37ec10ba845400001000c112c9d25adb053adf463e3e3df7ff1fb5d6bbcedf020000a3112c911cbb9a
5b6b6d2b8def023db61b59040400003c22582229d29bf24784c7f49d5e970000f8c5e21d244386be43552917b1bd2e378d31a70bee03
00003d51b1c4e4220c7daff24daaa4000060248225261571e87b95bf089700008cc750382633c1d0f72a365c2a630c0113008081a858
22ba8987be57a1720900c008044b44250dcfa71efa5e85700900c0400c85231aadf5adddbb3b8357dc86cb2763cc85730f0000588a60
89e06c4b1fa5d4936cab988b73997349b80400a02386c21194ec70f39c59a86c9cb3fd230000dd112c118cd6da361eff9ef92b7c43b8
0400a01b86c2114462ad84c6b2e172db18f394f7d30000202c2a96f04a6bbd95682ba1b1ee65ae280000588260096fb4d6074aa9c742
5f513b47f4deb9150000fc46b084175a6bbb7afabcf057f3adb44c0200000b102c319aedf9a894daafe495fc20211a0000cc215862b0
d67ccab795bd8afbb2831000006821586290c2e75376f18d36440000fc896089de2a994fd9c50d2bc50100f81b7d2cd14b61fd297db0
afc756fe4f030080f1a858a213994ff942a874bc65310f0000bf102cb196cc257ccc74bfef18f665ce2900005523586225adf5b19d4b
b8eaefe0a773adf50e2f0500a066044b2c25cdc0ff5a763f1c344f070a2217d6007a2058c261573a4bd3f30fce9d586583f996405998
e602f443b0c41f6438f7b9c2a6e7be30df1228c7954c73a1f303d011c112bf4920fac12b321a2722a000c618db4eec4eda8a01e88060
899fb4d6a7343df7eaaaa0e702d4ec56a6b9f09d063a2058a259a4f38557c2ab77ec270e14a19937fd89692ec07a04cb8ac9229d1716
e904f38d2171206fc618bb90f1419e04d35c80350896956a2dd2a1e979580c9f01f96b7f8ff94e032bb0577885b4d67b4aa9efb5bf0e
91fc1c1237c69c56f16c8132b57bfaf29dee40766cdb911f5be5dd941f5b01be8af1fa4901a5f9fdbbadbbb6a4f3c94c1e4fc3becf76
14efde18435fe281b43126cb078e61a4e12f4dcfe3db9621354410e9737e628ca181762564da507b84a7baefb404b5cd56586b82dbd0
91af3363cce8b9e832f7d5164c3e3977fa63bb031c1338d7a362591159a47864af610000200049444154c37cca695ccc5d3103c88b1d
02df6f3de22b095559496cc4ca2e1aed1d2ce5399c46eeb76ccf9d375aebe6ff7b09c525628e6525d84967721f58518a3eec2e4e5a6b
937b7701dbca4c9e47ee43c7f3732b73edfc90e5e7c986495b35b69f2509c6536fe2f1453ed7f77621ac736fc5180a2f9c7ce09f58a4
938457630c07a0084a180a971368dba53126bb8b93529e875afc5cac37c69817e7d6442d790e9331c6e865bf5b86deed68cf3be7cef4
3c1863b2ab608740c5b26032799a95dfe9d828a06a830864a86fdebe5448b20965720c9ab79f7105f3ceb9e5d7906c4e4e127aac8b5e
4ffbb9399200fc239350a9a4826d3fd78b3ef3552158164a4e3e37b5bf0e09fa421f3c74b0eae4742e4382395447563d8f2f199e8817
2ddcd8cfe9394895fd8dad1c3b77c6f7fbf594beca171228bf25f0d886baa97dda13c1b240320cc8f68ce9a26a8975d685463b0af143
16e4a5ac4be0b227e25cf6e25ef67a67f59db643f7763a821d865ef4a3947a6f17a7d8e93bce7fecd7ad04ca7b195ddb0ffcfb62395f
32ea5005826561643f5bda09a5ed53cd071d74d275a1dd874286df9a61c4a4bf172b5acdbc2ba94a658cb17d1c8f644ef847e72ff873
2a813297e1ee3ebed7baa887605910b9ea0bd9c70bfe50b5844f3789ce59ecdb89e2bb5c1ca76ce1bc40dbe3d0b9a50012a6cf023d93
b181d2be175f6df85d54799dff914a6ccc290017ce2d15205816a0b5e77789577da57a5bfb3c1c7867e72c3e155025f9247348537d1e
cb9aa2bfcdbd35d40a2984fdbb050172d7eee0b3a292fc07a9c4fe9e021021647eaa714e3dc13273f2a165e5779ea85ac237dbdbef39
a193d9d0397a1bf23c525ca0b46a3e283b31f975273b1c3521d2eb9c62697bb5eddce1577505048265c6645ed563edaf43c6360aae70
609c8791fffd6322f32e5785b02e7e2438ef72d573e23bedc7652b4c2eab107b21ff7ec870595dfb218265a6e4e0453ba1fc51e1c022
3e1a6e97d2f6e47b4acfa343d58ceff470771228a3bedf122e0f9d3b3008c13243d24e28e73e5ff81b150e2ce26bc8ef7ce25056caf3
98b76c018f92ef34f3a7fbfb6c2b9453fd7263cc4584f64a5520586686764245a2c281793ee7924dd953cff7f34825b0adab28f39dee
c72eca4961815095abb87d235866847642c5a26a893ff85ea420c3c9d1ab41019ec779267347dfb2b55f6787013e2743e5d2a83f6904
cb0c483ba127da09158d608979ab865b87b899a885cfb573cb383709ac7aef1284a85aae772d43d0a908ba50a81604cbc4c901f449da
88a05c6fd98d0773429c70a7a8c894f23cfafa50630fc39e68b956208265c2a487db233d2aab41d512bf05aae4d80b98a81522993be7
7b51c4c6943bf4f418ba6511cf0a090d81c3238265a2647ece8fda5f87ca7c605e16e69c38b78cb73f41753cc4b0f0a78917f37409cb
5c2c2ec70aec42112c13448fcaaa51e1405ba8a1c2d855cbd34041e27cc2ad1fbb0cc76f70b1b8140b650a45b04c0c3d2aabb75fc05e
cff0c418f312a86ab9117b483ce06296d4e7e951b5445508960991033d3d2ac18908bf19638e0355fbf663eec52d55cb9973c778519f
474bd7f981b4884355089689d05adb83d47eedaf037e627538e6859a2211bb254ea8cf76d2554b3a3ea02604cb04488fca0fb5bf0ef8
ed1df3b2d0262bab2f03bc289f22572dedbcba33e78ef1a658f8d667453373a7510d82e584a4f1f90b3d2ab100150efcc11873106828
396ae831c61c057a1e294f2161381cd520584e4416683cd1a3124b7c597c332a17a22a37c5672dc4f3483abc310a815a102c2720bb31
3c132ab10a2722cc33c6d88bd143e78e91627fd6e4797c74ee18ff3c521e72e6fb8c2a102c239303f863554f1a43b13a1c0ed991c777
0ba2e8814c765df11d92a33d8f01bbc6102c5105826544b23290c6e7e88a795958485a10f90c9793849e002139e545902cd044150896
91c810cdf72a9e2cbc61381ccb780e97932d20f41d9227ea69d909df67d480601981eca6735efc134508ac0ec7523e43d994a1c773b8
4c79e72a82258a47b00c4c6b7dca6e3a1881fe77584942d9d7557f2707f23c7cccb94c39bc112c513c826540b245236d6330c606c367
5847b64b1cbbca7af2cf99ccb97cefdc510ee659a27804cb40d8a2111e112cb196ac527e13685ff16864779e37819aa84f8e0b45948e
601980844aae4ce10b272274628c7931c6d83986d739bf62f23cb6026dff08202082a567ecfb8d00f83ca11763ccde80a1f1be7d1983
93ed1ffb0e8d3f39b7a4850b45148d60e991844af6fd86770c9fa12f3b346e8cd14aa9bb9c5f3c3b342ecfa36b159660094c8860e981
ddf75b6bfd42a844409c8c3088316657aa7e2be75e0ed849262aa9c26ee7fe3c3215a237e807d9de3825211e4f8acf332882e5483654
ca1532fb7e232482250693aadfe68a763e59ccc9b47b8ccbf358363c1ead3a3b621421f9a92d3608d9fecb76144e6b6d9452df9cbfe4
c7a3fdf7b5d657b2335dece7698b4207765d843ccf50fda69be7792bbf2fe55eaba3112c4790ab906742252248763711e4c3b6f39161
e5f980799ad9f3b86f3d8f7605f3c2f9cb58c8ee50645be269adef25f4fcfeb14148fa2fc71a85b3dbd77e9f7f1cf2d82ec64e059230
77d504e5d6f37c9630192bec7f90dff7bce0b9dad0795a42e8d4c618e7467478e17e85ca47e70e209c3776b52cafef7ab2db55e88d09
4ea4a977b664fb431b36b3be7091e771252bc963fd4e1b766e9c3bbaf938d5907de6a36cdbb66aeddcba84d6fa2860b535a43b99c292
252a9603c801855089d8a85ac22ba9fc65ffb992e751d53cb6118e321e65eb5b91ce31542a999749b0acc5c8ab54600ce65902e9c8f5
fb98f3a847df2a6fce1d11b27d9f08963d102a31312a324019a60ca539cf43ed1b2c73ed10f0203b5065e91fb93ef0d8ece4df802bc6
802e0896403ab2ac58da79da5a6bbbaadeceb5bccf61de76332c3c605eeaa984cb971c829aacddd8caa017eb4a04cb0e089548043bf0
00182db76ad8d0854e129ab3a95acac2a4ac43a562287c3d42250060012ef4800508962b48cb12422592c1d68e0080943114be44a43e
780080cc70811786bcae5b733f0d3b27f49d34c49f1fcabf9755d43f8792d9d6735a04cb05089548d86ec62b1d815214bd255f28d29c
fd407ede0dfc351b0ba621fcf1ffb5d6ce7fa4949ad926fab22140b62bae73c050f81c422500600d362be840b6527c9adb42f1db8850
3986dd9ef28b52eac7dc568a17b21a1b9e102c5b089500800e180a5f40f61fbf6d05c9f388fb8d0fb56f77d26b05cd22f6eb9e12c152
d8ab16422532c0090d981e154b61e7453655495b0d5c304c9d1b5bd57c96907945c8ec8f60f977a8dc77ee0000e0cff3c566c67b6d7b
615f03095d4676a34bbd2a39d4a756c83cc8f329c4577db0245402007aa8b65a2943dd2f3257f293f317c6b1abbd2f9552874aa98fc6
18bdecc7deaf94faaa94ba8ef8f4cf9b3999ce3df843d5c192500900e8c9c77494ac3a3bb402e50f8fd5dabbb900b9698c3930c65cac
6b1764ef37c69c1a63f69aff5e29f5462975e6fc65fff60998ab551b2c0995008001aa99e72c43def71e03a5ad4abe9730b8ebb3dfa4
ddbed118732421735b29f5e0fc25bff619225facca6049a804000c54c5568ed225e5d9536ba026506ec6e82169f7dc36c6ec48153374
c0b443e4f72cf2f95b75c1925089ccb11a159888af1d7752de1946aa944f1ebba49cc50a94f3a48ab9237332437a278b7caaefdaa16a
0b96844a14801d2380e9ec95fcdadbb99452a5f4b5cafbd00e4f3bb7462673327584eae50d43e315054b42250060241f15a93be79604
48b5ed87c747622b95492d7091eae5a573875fe7b587cb2a8225a1120030866cfb37c55684c149a8bcf1f87b6629542a17b12bcf0997
61151f2c099500000f7c0d832735bf5286bf7d864a2bc950d988182eabdc83bce86049a804007852dcc20c59c9ec3be8be1a63ae9c5b
1323e17216f85165d5afd497628325a11200e091af9d66520a1b1701b6a7cca97178e88b85b7d2b6a92a45064b42250af6c49b0bc4e5
b38d4c2aad86b4d67b01b665543975aeb0fd2e23ecd6e3ab6d53368a0b96727540a844a90896407cbe82e5ab73cb7442cd83ccea1815
639191d63ae939a7be15152c2554567775000008cad7c29d24aa79b260a78a1d843a3a09fcef17ddff745e31c19250894a5439191c98
8a84305f6d8652f9fe060b3a29ef2ab4c2e9f2bbbca82ac417112c09950080407c86b0544257b0452b396e6b68b77e0cddb8bea6ed1e
b30f9632778150895ab0a5231097cf40c01ce974311ae449d6c1523adb7f73ee000a2557d600229006d7be86315f6515720a420ecde6
5a9923587a926db0945079eedc01942bc93d868182f9dc962fa5d18690c7922c779bc9746e6892b20c96844a548a6a6577556ea506ef
4a9c5f195ace2ba0b978f720bb602913600995a811f32bbb23586214cfabc15562c132e490fc464d0b55baaaa9229a55b0940fabefcd
f2815c305403c4e3b5f29658b0087d915a55dfc60eaaaa8466132c651235a11235a36209c4e373b794d48245e8fdbcbf38b7d4edaaa6
679f45b09450f9e8dc01d463c68a70200e191ddbf0f8cb920a1691fa36866e3a9e0d634c55af45f2c1925009fc44b51288c7f7dece29
56ac8e9d5bfcfa22e7ef9c8468c3f4d5b9a57049074badf5262754e027be07e961ce6b81e4bcf3c9e333bb4ba87fe56f32e7f3dab9c3
af6c8680032d387aa8ad5aa932a8583e791e8e00724588e96727a7078ba4f8ae56a6fcddb5adfb5e9d5bfd7997d190b8eff75d65dc2c
7e946483a5d69a5009fcf24af3dede387660a81a86c17f92b996a157707f91ded3c992d6523eabd4d676adf3e2930c965a6b7b127deb
dc01d4895009442001c8e745891d0a4d7a1a8b5cb41e3a77f8759e78b8f47d8cfd98e2f48758920b96122a43ee630ae4a6aa56156345
6cceccbcd7f25453ad6c33c65c440a97a1170cf516607474bbf611a6a482a5d6fa82500938089609a2fd5359b4d67b9e77da5111fa45
7a23e1f263e05ff397d6fa5e16484dca3e06adf58bc7d151db124ed75ca96c24132c6582efbe730750b73b024c6f2cdcc110beab9549
ae065f452a6d6f6c485af1d7c6b2e1fd596b1d62b14c27322cffecb15279628c611b599144b094f2389dfa0117f32bfb8b510da96a8b
b6d2c9f409dfa365d9542bdbec85ac84a413e74ebfbe69ad4dccb99776918e5429cf9d3b87b12beadf1863921be29fd2e4c1523e547f
39770050b99e9c2656658b0f8ce23d18c8d072b6242cd9eae543e0e7702e01f32ad410b91d11b5bf4329f5c36395f2b33166931125d7
a4c152ae127d5d3900a5c96e282d11312a9654920b21732b7d572b4357fba290eae54ea480f94986c86dc87cb245a72141d3eef66387
d9e5df3012287d8e881eca5c4ae6be2f615f9cc5f784fec5bffa46fd70ee00d0f85ae3ae0d63c98924b41386bfca6017930458b4b35d
e245a104bdab0a17d9da21efddd45b47a562928aa5ec1f4aa804562054f617b1d51015cb02c8542cdfa1f2b2d49106a960da80a5a53d
51c85d7b5270267328370995dd450f96ecff0d7472c9cb34489415e1b5f7a92b48888bb72a2e08ed1c52095c36647e2d28645e4a98b4
23ba47cca1ec6f8a8ae53ddbad016bb168679818c1b2f42a4d15a41b89ef73d1758d952d3bbad20a99db995d18cf5af326edcf016172
9ca8732cd95507e8e44126cca3a740f3e5e6d94555ac3ccf984cc77a0cf00cde3364fa2719a5dc933ea1a1bf9bebbcca1cd10b461dc2
f947ac5fc4ae3a4067542b079013588c131727a4fc8518aeaeb25ab98e54ff2e161dd7644ef4ae8c34ec78da05672623a3f67b7acb7b
125f948aa5ecaa430374a003194e424fd236e67b84d7ed23d58e7c05fc9c14b9121ce82bf81c4bd9b6895009745344ffbb89ecc5f8b5
84caec3995330f8a5d090ef415b46219b1820094e20d13c787b10d913d0da5adc2fcca8cc994acfd00cf80ef2d2082cdb19406e8844a
a0bb334e4ec3c85cadd0a15231bf325f52e808112a4ff8defe22dfc366bee496fc2899f37895fb36970dc9373b323f744b76fbb23ff6
e2f6aaf61ec4412a963289fed9b903c02accd11a485ac7fc15e15731bf32535aeb9700ed8566c6982de7d6b25eb7262436616a73e442
dccf296e8728cf73b315169be73cf433630b0547cead1508152c630c490125a9f620e443a436432cacca94d6fa4af6a2f62dc990b44c
c0d7a18fe0d34912799ed51e2fbc2fde915e95844aa01ff69d1e482a0d31da0cdd39b72079b26d638890719d59a8dc49216c85263d4a
8b7f9e29f31a2ce955090cc2dcca71a2ac0697c6cac888848cf3408ff8c0b9256d9b957c768b9e9a90036fc152e6388598180d148d21
f0d1629de0695c9f9f50f3610f33bc18bcaf613b529903cdb6ab13f2122c65a821c6c479a0345f7947878bb81afc8eaa725e649e5d88
cfc6758eab9bede7d7eee76de785caee34c592e7f9b1f4e799aad1c132f0500350b259ed6d293c8855eda55a991119410b31cfeed518
136bea4510765ea85dc96e17962cfa91e099fd7c625bb9ecf03caf9dff10a38d5a152e6d859e02b470006a90d58ad2d4c46c6bc66af0
7c04de98a39a765301cfefc96d3210aaab04abc287b92754028364b5a23451b1aa9597ce2d48928ca0850a952735f53095a91fb58ca8
30cfdda3c1c1525680d356081826b715a52962181cbf4985ed3ed02b62ab6c35b604ab2248b3e9815f8382252bc08151beb210641c59
301863b4e481934e36428da0bdb23f3cd05def6029f35758010e0cf3c0821d2f62558f78af322073e4428da0ed38b70058aa57b00c3c
7f05a841d62b4a5320d5ca18d3705e736c2b531bd9ed2dd4ce4b87ecdf0ff4d3b762c9901030dc574e525ec4aa56b2cd66e2245486da
eded840b0ba0bfcec13260b359a006d70c818fa7b53e8a751ce2fd4a9b9c934285cacb4a17eb00a3fda3cb3f10b0d92c500b5681fbf1
2dd2ef39716e413202572aed0a70beafc0406b2b96b2651a8b7580e13eb30a7c3cad75ac0ae22bd5aa74050e9533568003e3acac584a
5fb01be70e005d9dd1087d3cadb55d99fb25d2af2354262ad40e29c286ca2de75600bdac0c964aa9ffc5b9054057b6b5103b3af8116b
1105fbb7274a6bfd14707eed2b6d85003f96064badf57f544afd4be70e005d31a4e6812cd80955a59ac785406202ee59ddb0a1728be9
2a801f0be7586aadff0da11218e52327aaf1a4776eac053becdf9e1879ff9f0995403e9c60295f64564402c37d651b406f62063d5602
274476797b0cf888ecb4874d4225e097132c9552ffab730b80ae2e99a3e787ac028f3504cefeed099116772177797b60a10e10c6cf60
69e7b0d82fb2d6faff534afdb7bcd6c0200ff4bff343aa55b15681b37f7b42a49d50c81677b64f250b758040fea1b5b6ab2df7798181
515e3959f9218b354256abe6b17f7b02222cd251d2fe8b055a4040ff241596192f32300ac36afedc47fc5d87ecdf3e3dadf541e0453a
4aa63b102a81c09a7643bb8127490325db667e9e1fb2ff7394bdc0651578acfe985822f04e3a8d8f2ca803e2f839c752aed80f79cd81
de3e52f1f24316eb7c8af4ebecd40586c027647753d25a9bc0a1d2b6137a43a84c4a72538664eb6a78f27b55b85cb99ff1c2029d1d72
c2f2435601c75aaca3d865655a7211f123f083b8a39d5092424e77188a8b4c8ffe683724f34fae737f524004270ca3fa21f3eb42ae02
9ec7bcca8948079297081711763e255528ac25bdbb635ed416cfe96329c3430fb5bf30c00a36541e2fbf1b5d49a83c8ff8829d714130
0da94a875ea0f32a739e691f95b0c4869e1975f2cc0996ea57b8dc61a538b0d025a1d28f0942e51dab82e3b31521a95286ae4a3743df
54a3d33779b06c55cf632d16ac8636c62c7cae917a8a0139b9a401ba1f1384ca193badc417b14ff267f6790f431658f96617cf6d4ef8
9c6cb0bd71eef0cc18a3633fb7142cac58aa5f2fc88bf4e67b75ee04ea43a8f4648250f9ca629db8ec7b2c812474a89cc9aa6f42655e
366401575452a5bc8f112a6bb634582ac225d020547a22f3ec6287ca2d5606c7d11af68ef11edb053abcb7e1dd05fa0d5fe422333869
6d752f737cdf25f2ba166b65b054122ea564cd9c4bd48850e9890c8bc65cfd6ded103cc26b55821e234c9f6aaa942cd08923e4e29673
392e78271739b75239ff3145a0acb53fe6da60d990f949ac16474d08959e48e88831d7ae6d9b851ce149308855093aa44a195de82e0a
fb36fcd963c4d02026173676fac595fc5b462e7242efe8b44e95f3ba972ede5946b65c8bb53b0630155a0a79203de2ee27580448a80c
4c6b6d57d87f8bf4eb1ea45b09261071115669aa5c34d8b962d9903e9727ce1d4039be122ac7d35aef451a1a6d7b255486d55a981323
54daf7f33da1725a327253ea74b8cb80abb7df4a05f54a2eb2abd0bb62d99093c677e70e206f8734d01e4f567cc6decd82853a014db0
9a9fef6242a405e173414fe98fa94e81da2a7565a7191e97d2dd6070b054d30e7301217c64efef7126ec7f4b9fca40260894cc6d4e54
21fdad175eb04c1c2cdbb21f71e93d14de669fbcac180fd58e0088a1193e25548e200124f4967d8bdc112afdb3f3eae4641b2b543ec8
6a6f4265a25a5d62ae337be877f2d9d28b426562b2fffc8f0a960dd9ecffab730790be990c9f32276f0459f51db3aad53893e30ffcbc
8f9bcdcada888b357ecfa3641a431e64adc576e23dae1fe47365c3e46e469fadeccf45a386c29d7f8ca171e4e58e5032ce04c3a46d0b
87b4d09f6d20ad94ba8abc6fb20d25f6847fefdc836cc8f0f85502ad7d9454528f86140ab4d64f09ec1b3ee95697be780d960d5a1320
a2ffa294fa67037e1ded84469093c9fd8407e2f70492f122b70c6a10280b25c5a5e388e7ff3b59f4327a1a93e7efc283541eefe5cf9f
ffbb968a7c9060a9febe02fee1dc014cef337b0b0f37f18523fd0c479293ffd5043b9110282b241d6476e567e867ee4142da55c863b7
5c301fcb3cc7f6c8eb4cc2e1ad3c8e7ba64f2d172c5836a85e2221afb2c51f078401261ef656329ff2c8b9159dc83eedb1b7d454120a
f6f8de0175081e2cd5dfd5cb5be65e6242cca71c2891ef2f43df0348b5e862a2f7ee4e02250b72808a785915be8e3d21c88454568e63
0a2784cafe6ca0d45abfc89496a942e583acea24547664f75bb60b116465f7f709debbb30c57e202f0244ac5729eb426893dbf0775a2
d2d55342230ccc85edc88649a5d4e984c755e64f02f8294ac5729e4cbe7f9f780f2ce48d4a574f89542855ab5136a172053bccddaa4c
de4c142a2fe5bddae4bb06404d152cd59fc3e387ce9dc03827ac1ceece2eca4924502aa952d2287b09bb00c7be57ad61ee29da3dbdca
f6a7f6c2ed80f70a40db2443e18b24b0e214f963d5770f89756c6071d502322de13491e6d3f47e05b05632c1b2417b220c742ddb8c61
05092a1709cd71666e5e4bab8fde17e7ce69d89d4ca84a02e82cb960d92060a2878f3e765e2899d6fa34a1b0d2a87e4b46099247f293
4a3bb63b099354fe910cf9aefc9ce2c4f13e6dc906cb0601132b307cbac2c43d0c57b9b473f356dc5fac09b6bceb6af01ecbc0503282
628f057b01e70bcf64b71c1b466f191d092ff960d92060620ead681690b6331713eee1bdca830c7bd7b15feeaf0acb9e04c914df8f33
d96799616e04d7baa8da4be862f755b6370dba55646db209960d16f9548f2ae59c842b938de203a584c803f949b547ef4caa929c4011
45e217baabdcc9451743ee0364172c1b6c135925e652fe392f6f8a7d9ffb282e50b6aa90f6e793f317d243551251c977e42293ef4757
2c62eb21db60d9900ff12d3bf914adda39790db9f23f4ea4edcc3ad9074ab970dd950099c36bdeb89420c95c494455c1b9f8ab31e6d4
b9158e7f38b764464e5e3f578a25baf215c355db9732e1451eab6475012061bdf9c9293cb6112431b9886b2066f2e71443eb6cbad151
f615cb4518262f42155787f6b36a57296656919c975cdb2009e64dd57127e3e0d83693cfc8154372488154299f3c9f6b475d2cc9b1f4
2040d0657e7f474506cb36aa98d979287d3b46093d4772f0cbf5e2e741e61c0d6edd21afc39673c72ff633b0e9dcaad4bf90dbff85fc
b79b855e40ce649eda05d548a44802dc8da787166cb4431ee7958fe384ddc6d4b9118ee2836523c11d47f0a72277609115db7b89b5d8
18e35a7ac26db57e725bf1999aeb56cb132a91489e9c4f7f78789cd1aa803e46320996dd54132cdb3268cf529bac776069ad14de2d28
40c23f0224b2e771f87b926d78b5d6572356acbfa7c1fa7a5506cb36adf571066d5b72673f648baef4b258ecd1da4a2cf7851e48c35d
6b1790eadb67212f1e17ea4cd63e6e44b8a4e55d07d507cb364266108b426512ed6864ee4d131a9b1f867531a5d726744a6593f99548
865c643f7b7a3cdb537ebeb5d643c20fc1b20382e5125aeb235981c9b0e630ffaf52eabf9afb2fff1fd9f9e3df75f917572ceed89c6b
fdd0feff9bcca345a19aede76e194ec714e4bcf8cdd3af9e740ad4c042125b097740b0ec402a5ba70416008962fe26820b30aa3775d5
f2bee779fdc41873ecdc8a3ffc132fc77ab6f46d5be0c88ab037f6c325d5030048819d2f766e8729ed109ffcbcd8766bb21a16f0c1f7
0aee47a9824e425adb9df0c9f08b8ae54899ee9002a04ecd70fa29ab5bd157e0750867c698c94226fca16239922de3db95cdb69a2915
cd6dd939000052b32117c13f16543617cd670662f9d2fa4c5ef179cc1715cbc06415dd81ecb4c28a630039b893aa260b15f09b84bdc7
095e910759e7c01ce20c102c2750c8967e00ea3193fd9bb3ddc8007e8c6c30eed3acb5608d16400921582642c2e69e844d569f0348d9
ab543459215b9909ab965d35f3882f089cd32058268ec089c2ddcd3dbd27f9196bbed7a992ffcf08411854342ba2b53e902e0439b996
b0c9f48ec008961993fe9a6c3388d81e24fcd955c52fcd9f25ae3296563d9bf21d6bc2ea16f3a5d7ba93cd1058795ea84cc3651b55f7
4008960593e0d9deae908a2716994938fcfdc35682c3cced2bcf36a1bf70022f945c78dd163212602f860e38f68d47b0447beb44aa32
e599b5f69ebea7829486ca471baea59ac909bc10b65d956d1754d053bab46d049d5bd109c112dec8c9b26dd95edf63d970f4dfcbbcd3
7f5ef93bf8daec1d6dffe4649dbfd6bceadd4456df86f42021934516052860787c9efd7ceed2e2a81f8225b2c00e47bff782263c564c
6bbd27a173afc085483309992caec89c4c09b92aa81a4f05b3078225922527d1d38a86e45f5b7dd938b9a233a914ed1554e17c95f96e
7c0f32274581d3423e9bef994eb41ec1124909bc176d4aae5b211e9f3caa000013d049444154926116782753530e0aa8f2db90b9c770
7919e473799c7135f390b65aab112c31291932392d7c88fb9266bd48818c021c657c527f904a2655a34264ba13dd67aae9cb112c119d
1c482e0a5d0d4b884436e4bb782027f6dce66c5e4bc8a4e25f182938340bd8529d4fbccd7cf7c5089688a2d030f920fdf9181641115a
27f4e3cce6367f35c69c3ab7a238092d60bb36c6ec39b7826089700a0c9397b26d1d57a9a886cc7bcea5a23993f9980c955764c2799b
6fa898bbfec9b90518c1563cb4d6575a6b7bc5f29871a8b40b064e64b843cb0fbb32a03a76c71c63cca6fd0ed813a952ea2ce1d7c056
597fd8e38f34ed8607f27a2efa79d25a5f2ce8611c959d7a648cd96d8ed5b6822dc7f0d0a8582e40c5125e14b0f3c2ab2c223ae50a14
e84602c569e2dbc53e4815938bc281a450d0c583316627b1c77e117071e8095b95baa8586230db3b4f6bfd22079d1c43e565ab22b929
95194225d091548a7612af66dad0fb28d5369a5c87f54e6b9dd4c245696cbeeddc81600896e8456bbda3b5be9730799ed94a525bb9f8
c8d036e09fbd2833c61c4d301cd9c7b9044c16dc8593dcf42739ce7f74ee4010044b7422f3686c98fc91f8b0d7bc139960ada5b2421b
202002bb4abb3537f3a32cac49c5be04cc7b59090f8f52ac0ccbb1ffceb903de112cb1949d3fd51aeacea581f96cae2ac9f036303119
32df4a3064da8be467099849cd0d4c4cdf4056cbc2292adf0b102ce16855276f3219eabe6ecd95dca22a09a46b2e647e4e64b8fc9dac
267f21607ab191da5ccb00ee984ab518c1123fc9dcc9a78caa9367ad216e567c0219b2dbe2b586cb539893b941c0f4e6839c5392986a
20efa7cff99f2c045b8260593959d9ddcc9d4c7da78d93d610f71143dc40399a3999b2bafc72e22746c0fcd3d0eae35b996a30690893
70ebb382fa9562c67204cb0acd35313f4ffc15688749fa85018593d5e50752c57c3ff17c4c02a61fe753b57b92dff9ec715ad725db87
ae46b0ac48d32a48be649f127ee6ed616ec2245029bb35636b3ee69455cc26602633b41b99af6a5f13309f42864c299e9c06289e9c48
5f4cacc0ce3b15902beddbc417e2dcc9ee180c6f03584a76fbb99af878667be2eed672bc92d7fcc6b9c3af3b394fd99ffbaeafadd67a
cbbe17f2b317f073f19e3de8bb2158164cae08531eea9e4998e4cb0aa097d6bcb929fbea5ed652c1eab1ad6369ceec9cfe4a9ffb200c
8517a8b52027d550f9b5d51a885009a037998bb933f130f97e455b45a6b68b526897cd42d1b29fa67f04cb82241e28ef5af32699f80c
c09bd6629fc3895ed5f3d217f8c88afd14b7e9f4692643de9ab994c3112c0ba0b53e4e3450da03d0a17c49ab998f04601ac6988b0903
66b3c0a7d85198f6369dadd7f9c1f98b79b96bedd6c6289a07ccb1cc58c27328ab9ad80e204d131f23ed4575755bfe69adf75a0b6952
ec8d6ca74d5cb0435b3804cb0c251c284f680f042035131e33eda8cd0ecdb4ff26d305eccf56eb4f9f0bb0ec6b7edfac2eb77f52e488
8b609991482d1ffa7a95ea24c30700923661c0ac66f53840b0cc40a27d2819ee069025adb51da2de8ffcd8b908471558bc9330d93de0
5ef6f14e25545ecaeaee1d4225801cb55691c75c78d22ceea96ede25ea42c53251135d51afc2500e80e248a3f5a7c817efccbd44b1a8
5826a6d58b32955079424f2f00a59246eb365c7e8ef8146d887db4fb593bf70099a362998889ae9a57a9b2550680ba4d305a3493ea25
538b50042a96099003d97322a1b269684ea804501d199d7913718719dbebf159fa3f02d923584ec8b60f4a68d8fb924009007f0c8f7f
8df8727cd75a5f39b7029961287c22b2dadb6753d8a1589403004bc834a5db88c76b5b29dd62681cb9a26219596b71ced4a1f28e4539
00b09a542f7722ee3fbe2143e3bbce3d4006a8584632c155ef323436078001e4387e1f710fec3363cc91732b90302a9611c8a4ece789
43a51d5e794f6373001846aa97766feb93482fe11799360564838a65605a6b5ba5fc30f1c3a075100078a4b5de92ea658c6e1eccbb44
36089681c8fede3f267e18761e25f374002090c8c5838fc6985be75620210c850720bb294c192a9b616f4225000424c7d9586d896eec
0250e7562021542c3d9a6062f7225f8d316c13060011451e1ab75bed1e3bb7020920587a22ad216e267c08acf606808945ec514c0f62
2489a1700fb4d6c71387cacfacf60680e949cfcbb3080f645fe6770249a16239d2c43be83cc8410c009010990b791ee111cda4051290
0482e540329ff229d27c9a456c95927d6501205111bb83102e910c86c20790f994cf13854a3b97f20da11200d2668cb1235a6fa45347
486fb5d64f7c1c900282654f5aeba309e7539e30971200f221bbf5d811ae59e0074db84412180aef41fa537e99e057dbab5d1b283968
0040a622355367581c93225876a4b5b643cf9f26f8d5ec9e03008588742e215c62320c85772057995384caaf844a0028873166cff6a0
0cfc84ded28a0853a162b986cc599962279df732f11b005018adf585ed4519f859d1441dd151b15c61a250399355df844a00289404be
d095cb7dd9c00388868ae5121385ca6b192601005420d29ccb4363cc85732b1000c172818942a56d25c49525005426d26a71a657210a
82e59c894225bbe80040c5229c7b6cdbba2dfa202334e658b6c8bedfb143e57b422500546f27f00e3d76a738568a233882a590a18877
ce1d616d33340100904a62e8f672ef64a30f201882e5df6d1f42cf6f697b9595dfeca40300f8490a0d87815f8d2f5a6b16892298eae7
584a2b86bf9c3bc2619e0b0060a9483d2edf701e420855074bb96afbeedc110ea11200b05684c53cb4b74310d5064badb5dd47f5d1b9
231c422500a09348e728fa5bc2bb9a83e58bac928b81500900e8456b7da494fa16f85563481c5e55b97847763a88152a15a11200d097
31c6aee07e08fcc2b14a1c5e55172ce50a30f4f6596ddb844a00c04007815f38bb9f78e83647a8485543e113ccabfc688ca1212d0060
30e93df925e02b3833c66c39b70203d4162c636ed7c8a4680080175aebd0276bce59f0a29aa170b9e28b152a2ff98202003cfa1af8c5
64ae25bca8a262a9b5b67bb0fe70ee08e3c118b313fb390200ca1661d4edab2c180206ab2558de47da079cb642008020b4d67621cf79
c857d718a39d1b811e8a1f0a9755e03142a5b547a80400842053ac5e43beb872ce04867f864aae586aad379552cfce1d619c1963f842
02008289d034fdd518b3e9dc0a74547ac5f2d8b9258c0742250020b40873203764c81d18a4d860293d2b43f6fd6a63237f00402c9781
7f0fc11283955cb18c55adb4abe89e9c5b01000823743bbb0fd24d05e8adc86029d5ca7de70eff66b4660000c4243bbacd02ff4a46e2
3048a915cb58d54abe7800802984ae5ab26e008314172c23562bedee3af7cead000084771bf837d8453cbbceadc01a25562c635d6571
35070098840c8707ed69c9a81c86283158c658097e46237400c0c4ae02ff7a82257a2b2a58c6eabd45cf4a004002420f87bf657538fa
2aad621923589e38b70000105f8c79feccb3442fc56ce9288b761e9d3bfc7bc330380020055aebd027f16b630c43e2e8aca48a658c6a
e525a112009090bbc00f85a170f45252b08c714545337400404a420f87db79969bceadc01245044b19067fe7dce1d71d7d2b01008989
318a46d5129d9552b18c51ad0cbdcb0100007d855e19ae58c0833e08961d196308960000002b94122c3f38b7f8157a72340000a9a262
89ceb20f9691f632a55a0900481173ff9194122a9631261587de360b0080de22b5c00b3d2a888294102cb79c5bfcbaa377250000c07a
542cd78bb1e20e0000207b2504cbd0257ae6af00009214699d01d059493bef8442b0040000e820eb6019e34acd18f3e4dc0800000007
15cbd5e85f09004819db2d2229b907cbd0154b5683030052b6c9bb839450b15c8df995008094c558bc33736e01962058020090af1815
4bd61aa03382e56af4b00400a4ec1def0e5242b004002043117b5852644167044b0000f2142b5832148ece08960000e4896089e4102c
0100c88cd67a33c296c63f1963180a4767044b0000f2b317e911b351087a21580200909f83488f986a257ac93d58f28107005425e630
38e759f445c57235f6600500a4e628d6e3617e25faca3a5846f8c0b3072b002035b186c12f9d5b8035a85802009009adb55db4f336d2
a3bd706e01d6282158865cb116ab471800005dc41a067f65181c4394102cef9d5bfc618e25002009b28563ac453ba7ce2d40072504cb
17e7167f36623f19000096385e7c73100c836390128265d0527dc44dfe01005848e656c6aa565e1b63d8c61183641f2c23cc01215802
00a61673689a61700c56caaa7016f000008aa4b53e8ab812fc8e453b18a3946019f24b106be80100008731e6d418a3ed8f52eabd52ea
4c293573fea21f31e771a240f6839afdb392799037ce1dfe7ce40a0e00902aadf58e344e1fd3e7d2562b19a5c32845044bf5eb4bf512
7015f7893186ab380040567a06cef7c698902dfc50819276de09d91a21d6f6590000786383a231e6c818b3d50ca7cb90fa67d9b2f155
7ed725a1123e9454b164381c00006042c5542c25f4859accaca85a020000ac56d250b80a3c1cbeefdc02000080df08963d482f310000
002c5054b0942da8ae9d3bfc21580200002c515ac55205de8aeaadd69ab9960000000b14b32abc4d6b6d5b26bc73eef06366db364cf4
d4000000925562c55205de928aaa250000c00245562c55f8aae5ab3166d3b9150000a062a5562c55e0aae586d69a2d1e0100005a8aad
58aa5f554bdb34fd8373873fdbb2121d0000a07a25572c5584f6402157a003000064a5e860291bea5f3a77f8f389853c000000bf143d
14ae7e0d87db4536cfce1d7ebd31c6bcc47e6e00000029297d285c49e03b74eef02be856920000003928be62d908dc7ec83a34c61030
010040b56a0a9676b79c47e70ebf58250e0000aa55fc507843025fe821f15be7160000804a54132cd5af706987aaaf9d3bfcb1db3d5e
c57e5e00000029a86628bc4d6b6d17f46c3877f873628c61671e00005095aa2a962d7bce2d7efda5b50efd3b000000925265b034c6d8
b9905f9d3bfcfa2e0b86000000aa50e5507843e6437e72eef08be6e90000a00a55074bf52b5cdad5e26f9d3bfc7935c66cc67e5e0000
00b1d53ac7b26dc7b9c5af0d09af00000045ab3e58ca30f57be70ebfde122e010040e9aa0f96ea57b8bc8fd03cdd864b1aa803008062
112c85344f0fbd52fc03e1120000948a60d9628c39b5cdcd9d3bfc225c0200802255bf2a7c11adb5ad5eee2fb8cba79931863e970000
a018542c1730c61c28a52edd7bbc62410f0000280a15cb1522552e5f95525b3451070000b92358ae21f3213facfe5b5e6c1b63a86002
00806c3114be86316637c2b0b8f5a8b5de756e050000c804c1b28348732ead1badf591732b0000400608961d49b83c8bf0abbec9dc4e
000080ac30c7b227adf5b152eaaf08bfeac118137a1f730000006f08960368ad6df5f23cc2afb22bc67758d403000072c050f800b2fd
e3c708bf6a4316f530ef120000248f8ae5085a6bbb73ce63a45f77272bd40100009244b01c496bbda994bab73be944f8750c8d030080
6431143e92dd3147f6fcbe8bf0eb181a070000c922587a22c3d427917e9d6d49742fd552000080243014ee99ec9e7313f1577e36c65c
39b70200004446b00c402a894f32741d033d2f0100c0e4180a0f40e65d6e469a7769bdd35a1be9af09000030092a9681c9429b6f117f
e54c568ebf38f70000000444b08c6082a171ebcc18c3ea710000100dc13222adb55d64f329f2af7d6f8cb9776e050000f08c6019d904
abc6ad07a5d42ec3e30000202416ef44668cb935c668097bb1bc534a3d6bad2f6a799d0100407c04cb89487ba0cf917ffb3eabc70100
40280c8527c0eea22355c5d8987f090000bca162998089aa97d60fadf58bd69ae6ea000060342a9689d15adf2aa53e4cf0a85e65810f
154c00003008c132415241bc8ddcf7b2c10a720000300843e109b25543d912f2648247d7ac2067881c0000f442c53271b26bcfed448b
7b1443e40000a02b2a9689b343d2b2b8e7bd84bcd8365a8b7cf66a7f3f0000c07204cb4cb486c70f277ac436607ea70f2600005886a1
f04cc92e3afb133ffa4b630c21130000fc44b0cc5802f32f1b7625f901f3300100a81bc1b20009054cebd018c39ee400005488605990
89fb5fcea31f2600009561f14e415a0b7ca65a41ded6f4c3b48b7d8e9c7b01004071a858162cb10aa635534aed31171300803251b12c
5862154cebadf4c4b455cc5badf596f337000040b6a8585624c10a66e35a5695331f1300808c51b1ac48ab82b92d8b6b52f1a9351ff3
4a56b9030080cc50b1ac58626d8a16b9534a1d31271300803c50b1ac58b30fb93146db5d74127c253eb4e664beb09524000069a36289
3f48783bcfe055b9966ae693730f00009804c1120bc9429f2b59c99d3abbe2fdd4feb000080080e9102cb196d6da6ed1b89fd12b45d0
04006002044b749670bba22e2e2568b21008008040089618446b6d2b825f327ef5ec2e401754350100f087608951a465d195ace0ce5d
13362f58140400407f044b782343e51709f7c51cea5ac2f315d54d00009623582208adf5ae84cc1c56950ff520734e2f98bb090000c1
12115412324f8c31c7cead000054846089a8b4d65b12324b9893d9f68661720040edd8d21151d94531c6985dd946f28dadf449dfc99c
9d112a0100a0628984c890f97186d5cc6d5691030040c5120931c6dc36d54ca9681e4a0ba0945d122a0100f8858a25b2a1b53e908a66
4a8b80deb3221c00805f0896c8960c9d1f29a53e4df41cee6c85d5b91500804a112c510cd905682f6255f3a31dbe776e0500a052044b
14ad15368f3cef08f4608cd9716e0500a062044b5449b69f3c92d0b931e035f86c8cb9726e0500a062044ba04502a75d24b4bba2c239
33c66c39b7020050398225d091844e5be1b44dde2f78dd000068514afdff0a065c21cc50e5d10000000049454e44ae426082
END FIGURE
